% !TEX program = lualatex
% ---------------------------------------------------------------------------------------------
%  Readback note 44: Poincaré series of cut Bessel functions: unfolding, one-sided Poisson summation, the adelic sieve, standard levels
%
%  Build:  python3 tools/rb.py note <slug>
%
%  One unit of a Lean artifact, rewritten as ordinary mathematics. The prose is hand-written. Lean names,
%  line ranges and counts come from data.tex; quoted Lean statements are copied from the artifact by the build.
%
%  Rules of the text: standard terminology; every symbol is explained where it first appears and listed again
%  in the notation at the end; no letter has two meanings; a proof is complete or says what it takes for granted;
%  Lean names appear only in the LEAN lines and in the concordance.
%
%  The signature rule (the build enforces it): a function, relation or predicate that the artifact defines is
%  written only through its notation macro, which the build makes from the spec (`notation`) with one slot for
%  each explicit argument of the Lean definition. No argument is fixed and dropped, also not "for brevity" and
%  not for one section. Every argument and every hypothesis of a stated declaration stands inside its
%  statement, or in a Setting that the statement cites by number; only implicit and instance arguments (the
%  ambient objects) may stand in the conventions.
% ---------------------------------------------------------------------------------------------
\documentclass[11pt,a4paper]{article}
\ifdefined\pdfvariable\pdfvariable suppressoptionalinfo \numexpr 1+2+4+8\relax\fi
\usepackage[a4paper,hmargin=35mm,top=22mm,bottom=27mm,footskip=12mm,nohead]{geometry}

\usepackage{fontspec}
\directlua{
  luaotfload.add_fallback("gfallback", {"STIXTwoMath-Regular.otf:mode=node;", "DejaVuSans.ttf:mode=node;"})
  luaotfload.add_fallback("gmonofallback", {"DejaVuSansMono.ttf:mode=node;", "STIXTwoMath-Regular.otf:mode=node;"})
}
\setmainfont{LibertinusSerif}[Extension=.otf,UprightFont=*-Regular,ItalicFont=*-Italic,
  BoldFont=*-Semibold,BoldItalicFont=*-SemiboldItalic,RawFeature={fallback=gfallback}]
\setsansfont{Inter}[Extension=.otf,UprightFont=*-Regular,ItalicFont=*-Italic,
  BoldFont=*-SemiBold,BoldItalicFont=*-SemiBoldItalic,RawFeature={fallback=gfallback}]
\setmonofont{JetBrainsMono}[Extension=.otf,UprightFont=*-Regular,ItalicFont=*-Italic,
  BoldFont=*-SemiBold,BoldItalicFont=*-SemiBoldItalic,Scale=MatchLowercase,
  RawFeature={-calt;-liga;fallback=gmonofallback}]
\usepackage{amsmath}
\usepackage{amsthm}
\usepackage[math-style=TeX]{unicode-math}
\setmathfont{LibertinusMath-Regular.otf}

\usepackage{xcolor}
\usepackage{tikz}
\usetikzlibrary{calc,arrows.meta,positioning,backgrounds}
\usepackage[protrusion=true,expansion=true]{microtype}
\usepackage{enumitem}
\usepackage{needspace}
\usepackage{booktabs}
\usepackage{array}
\usepackage{longtable}
\usepackage{fancyhdr}
\usepackage{lastpage}
\usepackage{multicol}
\usepackage[hidelinks,pdfusetitle]{hyperref}

\definecolor{ink}{HTML}{0F172A}
\definecolor{inksoft}{HTML}{475569}
\definecolor{inkfaint}{HTML}{94A3B8}
\definecolor{rulegrey}{HTML}{CBD5E1}
\definecolor{arith}{HTML}{C2410C}
\definecolor{est}{HTML}{7C3AED}
\definecolor{ana}{HTML}{0891B2}
\definecolor{ok}{HTML}{047857}
\definecolor{leancol}{HTML}{3730A3}

\title{Poincaré series of cut Bessel functions: unfolding, one-sided Poisson summation, the adelic sieve, standard levels}
\author{grh-compact}
\date{}

% ------------------------------------------------------------------------------------- macros
\newcommand{\leanfont}{\ttfamily\fontsize{8.6}{10.5}\selectfont\color{leancol}}
\newcommand{\lean}[1]{\csname LN@#1\endcsname}
\newcommand{\fillin}[1]{{\color{arith}\itshape[#1]}}  % a place of the template that is not written yet; the build reports every one that is left
\newcommand{\leanquote}[1]{\csname LQ@#1\endcsname}
% generated by rb.py note; do not edit
\newcommand{\ConcDecls}{53}
\newcommand{\Pin}{72fead10}
\newcommand{\PinText}{that of commit {\leanfont 72fead10}}
\newcommand{\ArtifactName}{grh-compact}
\expandafter\newcommand\csname LN@memunipotentiff\endcsname{\href{https://github.com/lsernau/grh-lean/blob/72fead102f663d5c7a407964a13dc0cf9d977207/grh-compact/src/New/Filling/ColumnCoords.lean\#L19-L21}{{\leanfont Filling.\allowbreak{}Fits.\allowbreak{}mem\_unipotent\_iff}}}
\expandafter\newcommand\csname LN@adelicToFinadelicEquivinclX\endcsname{\href{https://github.com/lsernau/grh-lean/blob/72fead102f663d5c7a407964a13dc0cf9d977207/grh-compact/src/New/Filling/ColumnCoords.lean\#L23-L25}{{\leanfont Filling.\allowbreak{}Fits.\allowbreak{}adelicToFin\_adelicEquiv\_inclX}}}
\expandafter\newcommand\csname LN@finiteEquivinclX\endcsname{\href{https://github.com/lsernau/grh-lean/blob/72fead102f663d5c7a407964a13dc0cf9d977207/grh-compact/src/New/Filling/ColumnCoords.lean\#L27-L28}{{\leanfont Filling.\allowbreak{}Fits.\allowbreak{}finiteEquiv\_inclX}}}
\expandafter\newcommand\csname LN@adelicToInfinclX\endcsname{\href{https://github.com/lsernau/grh-lean/blob/72fead102f663d5c7a407964a13dc0cf9d977207/grh-compact/src/New/Filling/ColumnCoords.lean\#L30-L34}{{\leanfont Filling.\allowbreak{}Fits.\allowbreak{}adelicToInf\_inclX}}}
\expandafter\newcommand\csname LN@weylBesselEquiv\endcsname{\href{https://github.com/lsernau/grh-lean/blob/72fead102f663d5c7a407964a13dc0cf9d977207/grh-compact/src/New/Filling/ColumnCoords.lean\#L38-L46}{{\leanfont Filling.\allowbreak{}Fits.\allowbreak{}weylBesselEquiv}}}
\expandafter\newcommand\csname LN@columnCoords\endcsname{\href{https://github.com/lsernau/grh-lean/blob/72fead102f663d5c7a407964a13dc0cf9d977207/grh-compact/src/New/Filling/ColumnCoords.lean\#L51-L53}{{\leanfont Filling.\allowbreak{}Fits.\allowbreak{}columnCoords}}}
\expandafter\newcommand\csname LN@existslaststagemap\endcsname{\href{https://github.com/lsernau/grh-lean/blob/72fead102f663d5c7a407964a13dc0cf9d977207/grh-compact/src/New/Filling/DegreeBounds.lean\#L18-L34}{{\leanfont Data.\allowbreak{}exists\_last\_stage\_map}}}
\expandafter\newcommand\csname LN@cardbesselCoordmulfactorialle\endcsname{\href{https://github.com/lsernau/grh-lean/blob/72fead102f663d5c7a407964a13dc0cf9d977207/grh-compact/src/New/Filling/DegreeBounds.lean\#L42-L50}{{\leanfont Data.\allowbreak{}card\_besselCoord\_mul\_factorial\_le}}}
\expandafter\newcommand\csname LN@rowParametersubmem\endcsname{\href{https://github.com/lsernau/grh-lean/blob/72fead102f663d5c7a407964a13dc0cf9d977207/grh-compact/src/New/Filling/DegreeBounds.lean\#L52-L54}{{\leanfont Data.\allowbreak{}rowParameter\_sub\_mem}}}
\expandafter\newcommand\csname LN@logfillingbesselDegreememIcc\endcsname{\href{https://github.com/lsernau/grh-lean/blob/72fead102f663d5c7a407964a13dc0cf9d977207/grh-compact/src/New/Filling/DegreeBounds.lean\#L61-L78}{{\leanfont Data.\allowbreak{}log\_filling\_besselDegree\_mem\_Icc}}}
\expandafter\newcommand\csname LN@charOnPlaces\endcsname{\href{https://github.com/lsernau/grh-lean/blob/72fead102f663d5c7a407964a13dc0cf9d977207/grh-compact/src/New/Filling/GlobalSieve.lean\#L26-L27}{{\leanfont Filling.\allowbreak{}charOnPlaces}}}
\expandafter\newcommand\csname LN@charOnPlacesevalPlaces\endcsname{\href{https://github.com/lsernau/grh-lean/blob/72fead102f663d5c7a407964a13dc0cf9d977207/grh-compact/src/New/Filling/GlobalSieve.lean\#L29-L31}{{\leanfont Filling.\allowbreak{}charOnPlaces\_evalPlaces}}}
\expandafter\newcommand\csname LN@frequencyGain\endcsname{\href{https://github.com/lsernau/grh-lean/blob/72fead102f663d5c7a407964a13dc0cf9d977207/grh-compact/src/New/Filling/GlobalSieve.lean\#L35-L36}{{\leanfont Filling.\allowbreak{}frequencyGain}}}
\expandafter\newcommand\csname LN@ratPoint\endcsname{\href{https://github.com/lsernau/grh-lean/blob/72fead102f663d5c7a407964a13dc0cf9d977207/grh-compact/src/New/Filling/GlobalSieve.lean\#L40-L40}{{\leanfont Filling.\allowbreak{}ratPoint}}}
\expandafter\newcommand\csname LN@sumfrequenciesle\endcsname{\href{https://github.com/lsernau/grh-lean/blob/72fead102f663d5c7a407964a13dc0cf9d977207/grh-compact/src/New/Filling/GlobalSieve.lean\#L49-L80}{{\leanfont Filling.\allowbreak{}sum\_frequencies\_le}}}
\expandafter\newcommand\csname LQ@sumfrequenciesle\endcsname{\leanquoteline{0}{lemma sum\_frequencies\_le [WeilData E] [T.StarStable]}\leanquoteline{4}{B k X (c : \_ → ℂ) (x : \_ → ℝ) (hB : 1 ≤ B) (hx : ∀ v, ‖c v‖⁻¹ ≤ |x v| ∧}\leanquoteline{2}{|x v| ≤ B) (h : ∀ y ∈ X, y ≠ 0 ∧ y ∈ ratIntegralOff T) : ∑ y ∈ X, (∏ v, dampedPower k ‖c v * Φ v}\leanquoteline{2}{y‖⁻¹) * (√(genIndex (ratToPlaces T y)) \textasciicircum{} k / denomIndex O (ratToPlaces T y) \textasciicircum{} 6) ≤}\leanquoteline{2}{(∏ v, |x v| \textasciicircum{} k * (S.b * B \textasciicircum{} 2) \textasciicircum{} 3) * (S.b * rootDiscr E) \textasciicircum{} (numInfPlaces E * 3)}}
\expandafter\newcommand\csname LN@onesidedpoisson\endcsname{\href{https://github.com/lsernau/grh-lean/blob/72fead102f663d5c7a407964a13dc0cf9d977207/grh-compact/src/New/Filling/GlobalSieve.lean\#L92-L164}{{\leanfont Filling.\allowbreak{}one\_sided\_poisson}}}
\expandafter\newcommand\csname LN@sieve\endcsname{\href{https://github.com/lsernau/grh-lean/blob/72fead102f663d5c7a407964a13dc0cf9d977207/grh-compact/src/New/Filling/GlobalSieve.lean\#L172-L213}{{\leanfont Filling.\allowbreak{}sieve}}}
\expandafter\newcommand\csname LN@Fits\endcsname{\href{https://github.com/lsernau/grh-lean/blob/72fead102f663d5c7a407964a13dc0cf9d977207/grh-compact/src/New/Filling/PoincareSeries.lean\#L26-L32}{{\leanfont Filling.\allowbreak{}Fits}}}
\expandafter\newcommand\csname LN@stdPrincipalLevel\endcsname{\href{https://github.com/lsernau/grh-lean/blob/72fead102f663d5c7a407964a13dc0cf9d977207/grh-compact/src/New/Filling/PoincareSeries.lean\#L36-L37}{{\leanfont Filling.\allowbreak{}stdPrincipalLevel}}}
\expandafter\newcommand\csname LN@JeqratGram\endcsname{\href{https://github.com/lsernau/grh-lean/blob/72fead102f663d5c7a407964a13dc0cf9d977207/grh-compact/src/New/Filling/PoincareSeries.lean\#L45-L46}{{\leanfont Filling.\allowbreak{}Fits.\allowbreak{}J\_eq\_ratGram}}}
\expandafter\newcommand\csname LN@ratEquiv\endcsname{\href{https://github.com/lsernau/grh-lean/blob/72fead102f663d5c7a407964a13dc0cf9d977207/grh-compact/src/New/Filling/PoincareSeries.lean\#L50-L51}{{\leanfont Filling.\allowbreak{}Fits.\allowbreak{}ratEquiv}}}
\expandafter\newcommand\csname LN@adelicEquiv\endcsname{\href{https://github.com/lsernau/grh-lean/blob/72fead102f663d5c7a407964a13dc0cf9d977207/grh-compact/src/New/Filling/PoincareSeries.lean\#L55-L56}{{\leanfont Filling.\allowbreak{}Fits.\allowbreak{}adelicEquiv}}}
\expandafter\newcommand\csname LN@adelicEquivratEquiv\endcsname{\href{https://github.com/lsernau/grh-lean/blob/72fead102f663d5c7a407964a13dc0cf9d977207/grh-compact/src/New/Filling/PoincareSeries.lean\#L58-L59}{{\leanfont Filling.\allowbreak{}Fits.\allowbreak{}adelicEquiv\_ratEquiv}}}
\expandafter\newcommand\csname LN@realHomeo\endcsname{\href{https://github.com/lsernau/grh-lean/blob/72fead102f663d5c7a407964a13dc0cf9d977207/grh-compact/src/New/Filling/PoincareSeries.lean\#L63-L66}{{\leanfont Filling.\allowbreak{}Fits.\allowbreak{}realHomeo}}}
\expandafter\newcommand\csname LN@realHomeomul\endcsname{\href{https://github.com/lsernau/grh-lean/blob/72fead102f663d5c7a407964a13dc0cf9d977207/grh-compact/src/New/Filling/PoincareSeries.lean\#L68-L71}{{\leanfont Filling.\allowbreak{}Fits.\allowbreak{}realHomeo\_mul}}}
\expandafter\newcommand\csname LN@finiteEquiv\endcsname{\href{https://github.com/lsernau/grh-lean/blob/72fead102f663d5c7a407964a13dc0cf9d977207/grh-compact/src/New/Filling/PoincareSeries.lean\#L75-L76}{{\leanfont Filling.\allowbreak{}Fits.\allowbreak{}finiteEquiv}}}
\expandafter\newcommand\csname LN@finiteEquiveqadelicToFin\endcsname{\href{https://github.com/lsernau/grh-lean/blob/72fead102f663d5c7a407964a13dc0cf9d977207/grh-compact/src/New/Filling/PoincareSeries.lean\#L78-L80}{{\leanfont Filling.\allowbreak{}Fits.\allowbreak{}finiteEquiv\_eq\_adelicToFin}}}
\expandafter\newcommand\csname LN@reductionfiniteEquiviff\endcsname{\href{https://github.com/lsernau/grh-lean/blob/72fead102f663d5c7a407964a13dc0cf9d977207/grh-compact/src/New/Filling/PoincareSeries.lean\#L82-L89}{{\leanfont Filling.\allowbreak{}Fits.\allowbreak{}reduction\_finiteEquiv\_iff}}}
\expandafter\newcommand\csname LN@memstdPrincipalLeveliff\endcsname{\href{https://github.com/lsernau/grh-lean/blob/72fead102f663d5c7a407964a13dc0cf9d977207/grh-compact/src/New/Filling/PoincareSeries.lean\#L91-L92}{{\leanfont Filling.\allowbreak{}Fits.\allowbreak{}mem\_stdPrincipalLevel\_iff}}}
\expandafter\newcommand\csname LN@cutBessel\endcsname{\href{https://github.com/lsernau/grh-lean/blob/72fead102f663d5c7a407964a13dc0cf9d977207/grh-compact/src/New/Filling/PoincareSeries.lean\#L96-L97}{{\leanfont Filling.\allowbreak{}Fits.\allowbreak{}cutBessel}}}
\expandafter\newcommand\csname LN@cutBesselratEquivmul\endcsname{\href{https://github.com/lsernau/grh-lean/blob/72fead102f663d5c7a407964a13dc0cf9d977207/grh-compact/src/New/Filling/PoincareSeries.lean\#L99-L103}{{\leanfont Filling.\allowbreak{}Fits.\allowbreak{}cutBessel\_ratEquiv\_mul}}}
\expandafter\newcommand\csname LN@poincarea\endcsname{\href{https://github.com/lsernau/grh-lean/blob/72fead102f663d5c7a407964a13dc0cf9d977207/grh-compact/src/New/Filling/PoincareSeries.lean\#L107-L107}{{\leanfont Filling.\allowbreak{}Fits.\allowbreak{}poincare}}}
\expandafter\newcommand\csname LN@measurablecutBessel\endcsname{\href{https://github.com/lsernau/grh-lean/blob/72fead102f663d5c7a407964a13dc0cf9d977207/grh-compact/src/New/Filling/PoincareSeries.lean\#L111-L112}{{\leanfont Filling.\allowbreak{}Fits.\allowbreak{}measurable\_cutBessel}}}
\expandafter\newcommand\csname LN@tsumcutBesselclassUnip\endcsname{\href{https://github.com/lsernau/grh-lean/blob/72fead102f663d5c7a407964a13dc0cf9d977207/grh-compact/src/New/Filling/PoincareSeries.lean\#L116-L121}{{\leanfont Filling.\allowbreak{}Fits.\allowbreak{}tsum\_cutBessel\_classUnip}}}
\expandafter\newcommand\csname LN@tsumcutBesselle\endcsname{\href{https://github.com/lsernau/grh-lean/blob/72fead102f663d5c7a407964a13dc0cf9d977207/grh-compact/src/New/Filling/PoincareSeries.lean\#L127-L136}{{\leanfont Filling.\allowbreak{}Fits.\allowbreak{}tsum\_cutBessel\_le}}}
\expandafter\newcommand\csname LN@integralpoincare\endcsname{\href{https://github.com/lsernau/grh-lean/blob/72fead102f663d5c7a407964a13dc0cf9d977207/grh-compact/src/New/Filling/PoincareSeries.lean\#L143-L152}{{\leanfont Filling.\allowbreak{}Fits.\allowbreak{}integral\_poincare}}}
\expandafter\newcommand\csname LN@poincareMap\endcsname{\href{https://github.com/lsernau/grh-lean/blob/72fead102f663d5c7a407964a13dc0cf9d977207/grh-compact/src/New/Filling/PoincareSeries.lean\#L156-L171}{{\leanfont Filling.\allowbreak{}Fits.\allowbreak{}poincareMap}}}
\expandafter\newcommand\csname LN@peterssonpoincare\endcsname{\href{https://github.com/lsernau/grh-lean/blob/72fead102f663d5c7a407964a13dc0cf9d977207/grh-compact/src/New/Filling/PoincareSeries.lean\#L177-L196}{{\leanfont Filling.\allowbreak{}Fits.\allowbreak{}petersson\_poincare}}}
\expandafter\newcommand\csname LQ@peterssonpoincare\endcsname{\leanquoteline{0}{lemma petersson\_poincare [StablyUnram E] \{k s\} (hs : Kf.reduceFamily \_ k = Sigma.uncurry s) :}\leanquoteline{2}{inner ℂ (V.formRep \_ k (X.poincareMap.1 (.single a' 1)).1) (X.poincareMap.1 (.single a 1)) =}\leanquoteline{2}{volume.real S.stdPrincipalLevel.carrier * S.amplitudeSum .univ a.1 s a'.1 (a.2, a'.2)}}
\expandafter\newcommand\csname LN@tailVector\endcsname{\href{https://github.com/lsernau/grh-lean/blob/72fead102f663d5c7a407964a13dc0cf9d977207/grh-compact/src/New/Filling/StandardLevel.lean\#L21-L22}{{\leanfont tailVector}}}
\expandafter\newcommand\csname LN@formtailVector\endcsname{\href{https://github.com/lsernau/grh-lean/blob/72fead102f663d5c7a407964a13dc0cf9d977207/grh-compact/src/New/Filling/StandardLevel.lean\#L28-L32}{{\leanfont form\_tailVector}}}
\expandafter\newcommand\csname LN@continuoustailVector\endcsname{\href{https://github.com/lsernau/grh-lean/blob/72fead102f663d5c7a407964a13dc0cf9d977207/grh-compact/src/New/Filling/StandardLevel.lean\#L34-L35}{{\leanfont continuous\_tailVector}}}
\expandafter\newcommand\csname LN@eventuallymassnearquadricle\endcsname{\href{https://github.com/lsernau/grh-lean/blob/72fead102f663d5c7a407964a13dc0cf9d977207/grh-compact/src/New/Filling/StandardLevel.lean\#L40-L61}{{\leanfont HyperbolicBasis.\allowbreak{}eventually\_mass\_near\_quadric\_le}}}
\expandafter\newcommand\csname LN@congGroup\endcsname{\href{https://github.com/lsernau/grh-lean/blob/72fead102f663d5c7a407964a13dc0cf9d977207/grh-compact/src/New/Filling/StandardLevel.lean\#L70-L70}{{\leanfont Level.\allowbreak{}congGroup}}}
\expandafter\newcommand\csname LN@IsLatticeLevel\endcsname{\href{https://github.com/lsernau/grh-lean/blob/72fead102f663d5c7a407964a13dc0cf9d977207/grh-compact/src/New/Filling/StandardLevel.lean\#L75-L82}{{\leanfont IntegralPlace.\allowbreak{}IsLatticeLevel}}}
\expandafter\newcommand\csname LN@IsStdLevel\endcsname{\href{https://github.com/lsernau/grh-lean/blob/72fead102f663d5c7a407964a13dc0cf9d977207/grh-compact/src/New/Filling/StandardLevel.lean\#L86-L88}{{\leanfont IntegralPlace.\allowbreak{}IsStdLevel}}}
\expandafter\newcommand\csname LN@existsisLatticeLevel\endcsname{\href{https://github.com/lsernau/grh-lean/blob/72fead102f663d5c7a407964a13dc0cf9d977207/grh-compact/src/New/Filling/StandardLevel.lean\#L95-L142}{{\leanfont IntegralPlace.\allowbreak{}exists\_isLatticeLevel}}}
\expandafter\newcommand\csname LQ@existsisLatticeLevel\endcsname{\leanquoteline{0}{lemma IntegralPlace.exists\_isLatticeLevel \{N\} [IsHermUnit (latMiddleGram 𝔰 N).splitGram] (hN :}\leanquoteline{2}{N ∈ R ∧ star N = N) (hg : g.1 ∈ R ∧ star g.1 = g) (hm : 𝔰.k = card ι + 2 * m + 1) : ∃ F :}\leanquoteline{2}{Level ι (latMiddleGram 𝔰 N) m, P.IsLatticeLevel F ∧ F.g = g}}
\expandafter\newcommand\csname LN@gauge\endcsname{\href{https://github.com/lsernau/grh-lean/blob/72fead102f663d5c7a407964a13dc0cf9d977207/grh-compact/src/New/Filling/StandardLevel.lean\#L148-L196}{{\leanfont IntegralPlace.\allowbreak{}IsStdLevel.\allowbreak{}gauge}}}
\expandafter\newcommand\csname LN@heisenbergMeasureeqfloorModulus\endcsname{\href{https://github.com/lsernau/grh-lean/blob/72fead102f663d5c7a407964a13dc0cf9d977207/grh-compact/src/New/Filling/StandardLevel.lean\#L200-L216}{{\leanfont IntegralPlace.\allowbreak{}IsStdLevel.\allowbreak{}heisenbergMeasure\_eq\_floorModulus}}}
\expandafter\newcommand\csname LN@ascConvunitaryLevel\endcsname{\href{https://github.com/lsernau/grh-lean/blob/72fead102f663d5c7a407964a13dc0cf9d977207/grh-compact/src/New/Filling/StandardLevel.lean\#L218-L221}{{\leanfont IntegralPlace.\allowbreak{}IsStdLevel.\allowbreak{}ascConv\_unitaryLevel}}}
\expandafter\newcommand\csname LN@unfoldbesselU\endcsname{\href{https://github.com/lsernau/grh-lean/blob/72fead102f663d5c7a407964a13dc0cf9d977207/grh-compact/src/New/Filling/StandardLevel.lean\#L223-L226}{{\leanfont IntegralPlace.\allowbreak{}IsStdLevel.\allowbreak{}unfold\_besselU}}}
\expandafter\newcommand\csname LN@groupMeasureboxSetle\endcsname{\href{https://github.com/lsernau/grh-lean/blob/72fead102f663d5c7a407964a13dc0cf9d977207/grh-compact/src/New/Filling/StandardLevel.lean\#L230-L235}{{\leanfont IntegralPlace.\allowbreak{}IsStdLevel.\allowbreak{}groupMeasure\_boxSet\_le}}}
\expandafter\newcommand\csname LN@lestackMeasureboxSet\endcsname{\href{https://github.com/lsernau/grh-lean/blob/72fead102f663d5c7a407964a13dc0cf9d977207/grh-compact/src/New/Filling/StandardLevel.lean\#L240-L253}{{\leanfont IntegralPlace.\allowbreak{}IsStdLevel.\allowbreak{}le\_stackMeasure\_boxSet}}}
\expandafter\newcommand\csname LN@isHaarMeasuregroupMeasure\endcsname{\href{https://github.com/lsernau/grh-lean/blob/72fead102f663d5c7a407964a13dc0cf9d977207/grh-compact/src/New/Filling/StandardLevel.lean\#L255-L258}{{\leanfont IntegralPlace.\allowbreak{}IsStdLevel.\allowbreak{}isHaarMeasure\_groupMeasure}}}
\newcommand{\ConcordanceRows}{\crow{memunipotentiff}{\href{https://github.com/lsernau/grh-lean/blob/72fead102f663d5c7a407964a13dc0cf9d977207/grh-compact/src/New/Filling/ColumnCoords.lean\#L19-L21}{{\leanfont Filling.\allowbreak{}Fits.\allowbreak{}mem\_unipotent\_iff}}}{ColumnCoords.lean}{19--21}
\crow{adelicToFinadelicEquivinclX}{\href{https://github.com/lsernau/grh-lean/blob/72fead102f663d5c7a407964a13dc0cf9d977207/grh-compact/src/New/Filling/ColumnCoords.lean\#L23-L25}{{\leanfont Filling.\allowbreak{}Fits.\allowbreak{}adelicToFin\_adelicEquiv\_inclX}}}{ColumnCoords.lean}{23--25}
\crow{finiteEquivinclX}{\href{https://github.com/lsernau/grh-lean/blob/72fead102f663d5c7a407964a13dc0cf9d977207/grh-compact/src/New/Filling/ColumnCoords.lean\#L27-L28}{{\leanfont Filling.\allowbreak{}Fits.\allowbreak{}finiteEquiv\_inclX}}}{ColumnCoords.lean}{27--28}
\crow{adelicToInfinclX}{\href{https://github.com/lsernau/grh-lean/blob/72fead102f663d5c7a407964a13dc0cf9d977207/grh-compact/src/New/Filling/ColumnCoords.lean\#L30-L34}{{\leanfont Filling.\allowbreak{}Fits.\allowbreak{}adelicToInf\_inclX}}}{ColumnCoords.lean}{30--34}
\crow{weylBesselEquiv}{\href{https://github.com/lsernau/grh-lean/blob/72fead102f663d5c7a407964a13dc0cf9d977207/grh-compact/src/New/Filling/ColumnCoords.lean\#L38-L46}{{\leanfont Filling.\allowbreak{}Fits.\allowbreak{}weylBesselEquiv}}}{ColumnCoords.lean}{38--46}
\crow{columnCoords}{\href{https://github.com/lsernau/grh-lean/blob/72fead102f663d5c7a407964a13dc0cf9d977207/grh-compact/src/New/Filling/ColumnCoords.lean\#L51-L53}{{\leanfont Filling.\allowbreak{}Fits.\allowbreak{}columnCoords}}}{ColumnCoords.lean}{51--53}
\crow{existslaststagemap}{\href{https://github.com/lsernau/grh-lean/blob/72fead102f663d5c7a407964a13dc0cf9d977207/grh-compact/src/New/Filling/DegreeBounds.lean\#L18-L34}{{\leanfont Data.\allowbreak{}exists\_last\_stage\_map}}}{DegreeBounds.lean}{18--34}
\crow{cardbesselCoordmulfactorialle}{\href{https://github.com/lsernau/grh-lean/blob/72fead102f663d5c7a407964a13dc0cf9d977207/grh-compact/src/New/Filling/DegreeBounds.lean\#L42-L50}{{\leanfont Data.\allowbreak{}card\_besselCoord\_mul\_factorial\_le}}}{DegreeBounds.lean}{42--50}
\crow{rowParametersubmem}{\href{https://github.com/lsernau/grh-lean/blob/72fead102f663d5c7a407964a13dc0cf9d977207/grh-compact/src/New/Filling/DegreeBounds.lean\#L52-L54}{{\leanfont Data.\allowbreak{}rowParameter\_sub\_mem}}}{DegreeBounds.lean}{52--54}
\crow{logfillingbesselDegreememIcc}{\href{https://github.com/lsernau/grh-lean/blob/72fead102f663d5c7a407964a13dc0cf9d977207/grh-compact/src/New/Filling/DegreeBounds.lean\#L61-L78}{{\leanfont Data.\allowbreak{}log\_filling\_besselDegree\_mem\_Icc}}}{DegreeBounds.lean}{61--78}
\crow{charOnPlaces}{\href{https://github.com/lsernau/grh-lean/blob/72fead102f663d5c7a407964a13dc0cf9d977207/grh-compact/src/New/Filling/GlobalSieve.lean\#L26-L27}{{\leanfont Filling.\allowbreak{}charOnPlaces}}}{GlobalSieve.lean}{26--27}
\crow{charOnPlacesevalPlaces}{\href{https://github.com/lsernau/grh-lean/blob/72fead102f663d5c7a407964a13dc0cf9d977207/grh-compact/src/New/Filling/GlobalSieve.lean\#L29-L31}{{\leanfont Filling.\allowbreak{}charOnPlaces\_evalPlaces}}}{GlobalSieve.lean}{29--31}
\crow{frequencyGain}{\href{https://github.com/lsernau/grh-lean/blob/72fead102f663d5c7a407964a13dc0cf9d977207/grh-compact/src/New/Filling/GlobalSieve.lean\#L35-L36}{{\leanfont Filling.\allowbreak{}frequencyGain}}}{GlobalSieve.lean}{35--36}
\crow{ratPoint}{\href{https://github.com/lsernau/grh-lean/blob/72fead102f663d5c7a407964a13dc0cf9d977207/grh-compact/src/New/Filling/GlobalSieve.lean\#L40-L40}{{\leanfont Filling.\allowbreak{}ratPoint}}}{GlobalSieve.lean}{40--40}
\crow{sumfrequenciesle}{\href{https://github.com/lsernau/grh-lean/blob/72fead102f663d5c7a407964a13dc0cf9d977207/grh-compact/src/New/Filling/GlobalSieve.lean\#L49-L80}{{\leanfont Filling.\allowbreak{}sum\_frequencies\_le}}}{GlobalSieve.lean}{49--80}
\crow{onesidedpoisson}{\href{https://github.com/lsernau/grh-lean/blob/72fead102f663d5c7a407964a13dc0cf9d977207/grh-compact/src/New/Filling/GlobalSieve.lean\#L92-L164}{{\leanfont Filling.\allowbreak{}one\_sided\_poisson}}}{GlobalSieve.lean}{92--164}
\crow{sieve}{\href{https://github.com/lsernau/grh-lean/blob/72fead102f663d5c7a407964a13dc0cf9d977207/grh-compact/src/New/Filling/GlobalSieve.lean\#L172-L213}{{\leanfont Filling.\allowbreak{}sieve}}}{GlobalSieve.lean}{172--213}
\crow{Fits}{\href{https://github.com/lsernau/grh-lean/blob/72fead102f663d5c7a407964a13dc0cf9d977207/grh-compact/src/New/Filling/PoincareSeries.lean\#L26-L32}{{\leanfont Filling.\allowbreak{}Fits}}}{PoincareSeries.lean}{26--32}
\crow{stdPrincipalLevel}{\href{https://github.com/lsernau/grh-lean/blob/72fead102f663d5c7a407964a13dc0cf9d977207/grh-compact/src/New/Filling/PoincareSeries.lean\#L36-L37}{{\leanfont Filling.\allowbreak{}stdPrincipalLevel}}}{PoincareSeries.lean}{36--37}
\crow{JeqratGram}{\href{https://github.com/lsernau/grh-lean/blob/72fead102f663d5c7a407964a13dc0cf9d977207/grh-compact/src/New/Filling/PoincareSeries.lean\#L45-L46}{{\leanfont Filling.\allowbreak{}Fits.\allowbreak{}J\_eq\_ratGram}}}{PoincareSeries.lean}{45--46}
\crow{ratEquiv}{\href{https://github.com/lsernau/grh-lean/blob/72fead102f663d5c7a407964a13dc0cf9d977207/grh-compact/src/New/Filling/PoincareSeries.lean\#L50-L51}{{\leanfont Filling.\allowbreak{}Fits.\allowbreak{}ratEquiv}}}{PoincareSeries.lean}{50--51}
\crow{adelicEquiv}{\href{https://github.com/lsernau/grh-lean/blob/72fead102f663d5c7a407964a13dc0cf9d977207/grh-compact/src/New/Filling/PoincareSeries.lean\#L55-L56}{{\leanfont Filling.\allowbreak{}Fits.\allowbreak{}adelicEquiv}}}{PoincareSeries.lean}{55--56}
\crow{adelicEquivratEquiv}{\href{https://github.com/lsernau/grh-lean/blob/72fead102f663d5c7a407964a13dc0cf9d977207/grh-compact/src/New/Filling/PoincareSeries.lean\#L58-L59}{{\leanfont Filling.\allowbreak{}Fits.\allowbreak{}adelicEquiv\_ratEquiv}}}{PoincareSeries.lean}{58--59}
\crow{realHomeo}{\href{https://github.com/lsernau/grh-lean/blob/72fead102f663d5c7a407964a13dc0cf9d977207/grh-compact/src/New/Filling/PoincareSeries.lean\#L63-L66}{{\leanfont Filling.\allowbreak{}Fits.\allowbreak{}realHomeo}}}{PoincareSeries.lean}{63--66}
\crow{realHomeomul}{\href{https://github.com/lsernau/grh-lean/blob/72fead102f663d5c7a407964a13dc0cf9d977207/grh-compact/src/New/Filling/PoincareSeries.lean\#L68-L71}{{\leanfont Filling.\allowbreak{}Fits.\allowbreak{}realHomeo\_mul}}}{PoincareSeries.lean}{68--71}
\crow{finiteEquiv}{\href{https://github.com/lsernau/grh-lean/blob/72fead102f663d5c7a407964a13dc0cf9d977207/grh-compact/src/New/Filling/PoincareSeries.lean\#L75-L76}{{\leanfont Filling.\allowbreak{}Fits.\allowbreak{}finiteEquiv}}}{PoincareSeries.lean}{75--76}
\crow{finiteEquiveqadelicToFin}{\href{https://github.com/lsernau/grh-lean/blob/72fead102f663d5c7a407964a13dc0cf9d977207/grh-compact/src/New/Filling/PoincareSeries.lean\#L78-L80}{{\leanfont Filling.\allowbreak{}Fits.\allowbreak{}finiteEquiv\_eq\_adelicToFin}}}{PoincareSeries.lean}{78--80}
\crow{reductionfiniteEquiviff}{\href{https://github.com/lsernau/grh-lean/blob/72fead102f663d5c7a407964a13dc0cf9d977207/grh-compact/src/New/Filling/PoincareSeries.lean\#L82-L89}{{\leanfont Filling.\allowbreak{}Fits.\allowbreak{}reduction\_finiteEquiv\_iff}}}{PoincareSeries.lean}{82--89}
\crow{memstdPrincipalLeveliff}{\href{https://github.com/lsernau/grh-lean/blob/72fead102f663d5c7a407964a13dc0cf9d977207/grh-compact/src/New/Filling/PoincareSeries.lean\#L91-L92}{{\leanfont Filling.\allowbreak{}Fits.\allowbreak{}mem\_stdPrincipalLevel\_iff}}}{PoincareSeries.lean}{91--92}
\crow{cutBessel}{\href{https://github.com/lsernau/grh-lean/blob/72fead102f663d5c7a407964a13dc0cf9d977207/grh-compact/src/New/Filling/PoincareSeries.lean\#L96-L97}{{\leanfont Filling.\allowbreak{}Fits.\allowbreak{}cutBessel}}}{PoincareSeries.lean}{96--97}
\crow{cutBesselratEquivmul}{\href{https://github.com/lsernau/grh-lean/blob/72fead102f663d5c7a407964a13dc0cf9d977207/grh-compact/src/New/Filling/PoincareSeries.lean\#L99-L103}{{\leanfont Filling.\allowbreak{}Fits.\allowbreak{}cutBessel\_ratEquiv\_mul}}}{PoincareSeries.lean}{99--103}
\crow{poincarea}{\href{https://github.com/lsernau/grh-lean/blob/72fead102f663d5c7a407964a13dc0cf9d977207/grh-compact/src/New/Filling/PoincareSeries.lean\#L107-L107}{{\leanfont Filling.\allowbreak{}Fits.\allowbreak{}poincare}}}{PoincareSeries.lean}{107--107}
\crow{measurablecutBessel}{\href{https://github.com/lsernau/grh-lean/blob/72fead102f663d5c7a407964a13dc0cf9d977207/grh-compact/src/New/Filling/PoincareSeries.lean\#L111-L112}{{\leanfont Filling.\allowbreak{}Fits.\allowbreak{}measurable\_cutBessel}}}{PoincareSeries.lean}{111--112}
\crow{tsumcutBesselclassUnip}{\href{https://github.com/lsernau/grh-lean/blob/72fead102f663d5c7a407964a13dc0cf9d977207/grh-compact/src/New/Filling/PoincareSeries.lean\#L116-L121}{{\leanfont Filling.\allowbreak{}Fits.\allowbreak{}tsum\_cutBessel\_classUnip}}}{PoincareSeries.lean}{116--121}
\crow{tsumcutBesselle}{\href{https://github.com/lsernau/grh-lean/blob/72fead102f663d5c7a407964a13dc0cf9d977207/grh-compact/src/New/Filling/PoincareSeries.lean\#L127-L136}{{\leanfont Filling.\allowbreak{}Fits.\allowbreak{}tsum\_cutBessel\_le}}}{PoincareSeries.lean}{127--136}
\crow{integralpoincare}{\href{https://github.com/lsernau/grh-lean/blob/72fead102f663d5c7a407964a13dc0cf9d977207/grh-compact/src/New/Filling/PoincareSeries.lean\#L143-L152}{{\leanfont Filling.\allowbreak{}Fits.\allowbreak{}integral\_poincare}}}{PoincareSeries.lean}{143--152}
\crow{poincareMap}{\href{https://github.com/lsernau/grh-lean/blob/72fead102f663d5c7a407964a13dc0cf9d977207/grh-compact/src/New/Filling/PoincareSeries.lean\#L156-L171}{{\leanfont Filling.\allowbreak{}Fits.\allowbreak{}poincareMap}}}{PoincareSeries.lean}{156--171}
\crow{peterssonpoincare}{\href{https://github.com/lsernau/grh-lean/blob/72fead102f663d5c7a407964a13dc0cf9d977207/grh-compact/src/New/Filling/PoincareSeries.lean\#L177-L196}{{\leanfont Filling.\allowbreak{}Fits.\allowbreak{}petersson\_poincare}}}{PoincareSeries.lean}{177--196}
\crow{tailVector}{\href{https://github.com/lsernau/grh-lean/blob/72fead102f663d5c7a407964a13dc0cf9d977207/grh-compact/src/New/Filling/StandardLevel.lean\#L21-L22}{{\leanfont tailVector}}}{StandardLevel.lean}{21--22}
\crow{formtailVector}{\href{https://github.com/lsernau/grh-lean/blob/72fead102f663d5c7a407964a13dc0cf9d977207/grh-compact/src/New/Filling/StandardLevel.lean\#L28-L32}{{\leanfont form\_tailVector}}}{StandardLevel.lean}{28--32}
\crow{continuoustailVector}{\href{https://github.com/lsernau/grh-lean/blob/72fead102f663d5c7a407964a13dc0cf9d977207/grh-compact/src/New/Filling/StandardLevel.lean\#L34-L35}{{\leanfont continuous\_tailVector}}}{StandardLevel.lean}{34--35}
\crow{eventuallymassnearquadricle}{\href{https://github.com/lsernau/grh-lean/blob/72fead102f663d5c7a407964a13dc0cf9d977207/grh-compact/src/New/Filling/StandardLevel.lean\#L40-L61}{{\leanfont HyperbolicBasis.\allowbreak{}eventually\_mass\_near\_quadric\_le}}}{StandardLevel.lean}{40--61}
\crow{congGroup}{\href{https://github.com/lsernau/grh-lean/blob/72fead102f663d5c7a407964a13dc0cf9d977207/grh-compact/src/New/Filling/StandardLevel.lean\#L70-L70}{{\leanfont Level.\allowbreak{}congGroup}}}{StandardLevel.lean}{70--70}
\crow{IsLatticeLevel}{\href{https://github.com/lsernau/grh-lean/blob/72fead102f663d5c7a407964a13dc0cf9d977207/grh-compact/src/New/Filling/StandardLevel.lean\#L75-L82}{{\leanfont IntegralPlace.\allowbreak{}IsLatticeLevel}}}{StandardLevel.lean}{75--82}
\crow{IsStdLevel}{\href{https://github.com/lsernau/grh-lean/blob/72fead102f663d5c7a407964a13dc0cf9d977207/grh-compact/src/New/Filling/StandardLevel.lean\#L86-L88}{{\leanfont IntegralPlace.\allowbreak{}IsStdLevel}}}{StandardLevel.lean}{86--88}
\crow{existsisLatticeLevel}{\href{https://github.com/lsernau/grh-lean/blob/72fead102f663d5c7a407964a13dc0cf9d977207/grh-compact/src/New/Filling/StandardLevel.lean\#L95-L142}{{\leanfont IntegralPlace.\allowbreak{}exists\_isLatticeLevel}}}{StandardLevel.lean}{95--142}
\crow{gauge}{\href{https://github.com/lsernau/grh-lean/blob/72fead102f663d5c7a407964a13dc0cf9d977207/grh-compact/src/New/Filling/StandardLevel.lean\#L148-L196}{{\leanfont IntegralPlace.\allowbreak{}IsStdLevel.\allowbreak{}gauge}}}{StandardLevel.lean}{148--196}
\crow{heisenbergMeasureeqfloorModulus}{\href{https://github.com/lsernau/grh-lean/blob/72fead102f663d5c7a407964a13dc0cf9d977207/grh-compact/src/New/Filling/StandardLevel.lean\#L200-L216}{{\leanfont IntegralPlace.\allowbreak{}IsStdLevel.\allowbreak{}heisenbergMeasure\_eq\_floorModulus}}}{StandardLevel.lean}{200--216}
\crow{ascConvunitaryLevel}{\href{https://github.com/lsernau/grh-lean/blob/72fead102f663d5c7a407964a13dc0cf9d977207/grh-compact/src/New/Filling/StandardLevel.lean\#L218-L221}{{\leanfont IntegralPlace.\allowbreak{}IsStdLevel.\allowbreak{}ascConv\_unitaryLevel}}}{StandardLevel.lean}{218--221}
\crow{unfoldbesselU}{\href{https://github.com/lsernau/grh-lean/blob/72fead102f663d5c7a407964a13dc0cf9d977207/grh-compact/src/New/Filling/StandardLevel.lean\#L223-L226}{{\leanfont IntegralPlace.\allowbreak{}IsStdLevel.\allowbreak{}unfold\_besselU}}}{StandardLevel.lean}{223--226}
\crow{groupMeasureboxSetle}{\href{https://github.com/lsernau/grh-lean/blob/72fead102f663d5c7a407964a13dc0cf9d977207/grh-compact/src/New/Filling/StandardLevel.lean\#L230-L235}{{\leanfont IntegralPlace.\allowbreak{}IsStdLevel.\allowbreak{}groupMeasure\_boxSet\_le}}}{StandardLevel.lean}{230--235}
\crow{lestackMeasureboxSet}{\href{https://github.com/lsernau/grh-lean/blob/72fead102f663d5c7a407964a13dc0cf9d977207/grh-compact/src/New/Filling/StandardLevel.lean\#L240-L253}{{\leanfont IntegralPlace.\allowbreak{}IsStdLevel.\allowbreak{}le\_stackMeasure\_boxSet}}}{StandardLevel.lean}{240--253}
\crow{isHaarMeasuregroupMeasure}{\href{https://github.com/lsernau/grh-lean/blob/72fead102f663d5c7a407964a13dc0cf9d977207/grh-compact/src/New/Filling/StandardLevel.lean\#L255-L258}{{\leanfont IntegralPlace.\allowbreak{}IsStdLevel.\allowbreak{}isHaarMeasure\_groupMeasure}}}{StandardLevel.lean}{255--258}}
% the notations of the defined functions: one slot for each explicit argument of the Lean declaration
\newcommand{\wbe}[1]{\mathsf{A}_{#1}}
\newcommand{\colc}[1]{\mathrm{cc}_{#1}}
\newcommand{\charT}[2]{\psi_{#1,#2}}
\newcommand{\fgain}[3]{\beta_{#1,#2}(#3)}
\newcommand{\ratpt}[3]{\mathrm{pt}_{#1}(#2;#3)}
\newcommand{\Fits}[6]{\operatorname{Fits}(#1;#2,#3,#4,#5,#6)}
\newcommand{\stdK}[1]{\mathsf{K}^{\mathrm{std}}_{g}(#1)}
\newcommand{\ratEq}[1]{\mathrm{rat}_{#1}}
\newcommand{\adEq}[1]{\mathrm{ad}_{#1}}
\newcommand{\reH}[1]{\mathrm{re}_{#1}}
\newcommand{\finEq}[1]{\mathrm{fin}_{#1}}
\newcommand{\cutB}[3]{\mathrm{cb}_{#1}(#2;#3)}
\newcommand{\tailv}[3]{\operatorname{tv}(#1,#2,#3)}
\newcommand{\cgp}[1]{\mathsf{K}^{\circ}(#1)}
\newcommand{\IsLat}[1]{#1\in\operatorname{LatLev}}
\newcommand{\IsStd}[1]{#1\in\operatorname{StdLev}}
\newcommand{\UnitMembers}{497}
\newcommand{\UnitInNote}{53}
\newcommand{\UnitProofs}{179}
\newcommand{\UnitSteps}{1371}
\newcommand{\UnitAllSteps}{1479}


\linespread{1.08}                       % 11 pt on 14.7 pt
\setlength{\parindent}{0pt}
\setlength{\parskip}{5.5pt plus 1.5pt minus 0.5pt}
\setlength{\abovedisplayskip}{7.5pt plus 2pt minus 3pt}
\setlength{\belowdisplayskip}{7.5pt plus 2pt minus 3pt}
\setlength{\abovedisplayshortskip}{3pt plus 2pt}
\setlength{\belowdisplayshortskip}{6pt plus 2pt minus 2pt}
\emergencystretch=1.5em
\clubpenalty=10000 \widowpenalty=10000 \displaywidowpenalty=10000
\relpenalty=7000 \binoppenalty=9000
\raggedbottom
\pagestyle{fancy}\fancyhf{}\renewcommand{\headrulewidth}{0pt}
\fancyfoot[C]{\sffamily\fontsize{7.6}{9}\selectfont\color{inkfaint}\thepage\,/\,\pageref*{LastPage}}
\setlist{topsep=2pt,partopsep=0pt,parsep=0pt,itemsep=4pt,leftmargin=2.1em,labelsep=0.55em}

\newcommand{\Q}{\mathbb{Q}}
\newcommand{\Z}{\mathbb{Z}}
\newcommand{\F}{\mathbb{F}}
\newcommand{\len}{\operatorname{len}}
\newcommand{\cok}{\operatorname{cok}}
\newcommand{\End}{\operatorname{End}}
\newcommand{\Hom}{\operatorname{Hom}}
\newcommand{\Gal}{\operatorname{Gal}}
\newcommand{\diag}{\operatorname{diag}}
\newcommand{\kw}[1]{\emph{#1}}
\newcommand{\m}[1]{\mbox{$#1$}}        % a short formula that must not be broken across lines
\newcommand{\head}[1]{\par\addvspace{17pt plus 4pt minus 2pt}\Needspace{6\baselineskip}%
  \noindent{\fontsize{13.2}{16}\selectfont\bfseries\color{ink}#1}\par\nopagebreak\vspace{-1pt}}
\newcommand{\sect}[2]{\refstepcounter{section}\label{#2}\par\addvspace{19pt plus 4pt minus 2pt}\Needspace{7\baselineskip}%
  \noindent\llap{{\fontsize{13.2}{16}\selectfont\bfseries\color{inkfaint}\thesection}\hspace{9pt}}%
  {\fontsize{13.2}{16}\selectfont\bfseries\color{ink}#1}\par\nopagebreak\vspace{-1pt}}
\newcommand{\subhead}[1]{\par\addvspace{9pt plus 2pt}\Needspace{4\baselineskip}\noindent\textbf{#1}\enspace\ignorespaces}
% where a statement lives in the Lean archive
\newcommand{\leantag}{{\sffamily\bfseries\fontsize{6.3}{8}\selectfont\color{leancol}\addfontfeatures{LetterSpace=7}LEAN}}
\newcommand{\sep}{\unskip\nobreak\hspace{5pt}{\color{inkfaint}\textperiodcentered}\hspace{5pt plus 2pt}\ignorespaces}
\newcommand{\leanline}[1]{\par\nopagebreak\vspace{-2.5pt}{\fontsize{8.9}{12.4}\selectfont\color{inksoft}\raggedright
  \noindent\llap{\leantag\hspace{9pt}}#1\par}\vspace{1pt}}
% a displayed statement in a tinted box: #1 colour, #2 text
\newcommand{\statement}[2]{\par\addvspace{7pt}\begingroup\setlength{\fboxsep}{8.5pt}\setlength{\fboxrule}{0.6pt}%
  \noindent\fcolorbox{#1!30!white}{#1!3!white}{\parbox{\dimexpr\linewidth-2\fboxsep-2\fboxrule\relax}{#2}}%
  \endgroup\par\addvspace{7pt}}
\newcommand{\nt}[2]{\par\noindent\hangindent=1.1em\hangafter=1 #1\enspace{\color{inksoft}#2}\par\vspace{1pt}}
\newcommand{\legdot}[1]{\tikz[baseline=-0.62ex]\fill[#1] (0,0) circle (2.5pt);\hspace{3pt}}
\newcommand{\crow}[4]{\csname CS@#1\endcsname & #2 & \textcolor{inksoft}{#3} & \textcolor{inksoft}{#4}\\}
\newcommand{\leanquoteline}[2]{\par\noindent\hangindent=\dimexpr#1\fontdimen2\font+2.2em\relax\hangafter=1
  \hspace*{#1\fontdimen2\font}#2}
\newcommand{\quoted}[1]{\par\vspace{3pt}{\ttfamily\fontsize{7.5}{10.4}\selectfont\color{leancol}\raggedright
  \setlength{\parskip}{0pt}\leanquote{#1}\par}\vspace{4pt}}
\newcommand{\bx}[3]{{\sffamily\bfseries\fontsize{7.3}{9.2}\selectfont\color{#1!82!black}#2}\\[1.5pt]#3}
\newcommand{\cs}[2]{\expandafter\newcommand\csname CS@#1\endcsname{#2}}

% theorems
\newtheoremstyle{plainx}{7pt plus 2pt}{7pt plus 2pt}{\itshape}{}{\bfseries}{.}{0.5em}{}
\newtheoremstyle{uprightx}{7pt plus 2pt}{7pt plus 2pt}{}{}{\bfseries}{.}{0.5em}{}
\theoremstyle{plainx}
\newtheorem{theorem}{Theorem}[section]
\newtheorem{lemma}[theorem]{Lemma}
\newtheorem{proposition}[theorem]{Proposition}
\newtheorem{corollary}[theorem]{Corollary}
\theoremstyle{uprightx}
\newtheorem{definition}[theorem]{Definition}
\newtheorem{setting}[theorem]{Setting}          % hypotheses shared by several statements; each of them cites it by number
\newtheorem{remark}[theorem]{Remark}
\newtheorem{example}[theorem]{Example}
\renewcommand{\qedsymbol}{{\setlength{\fboxsep}{0pt}\setlength{\fboxrule}{0.5pt}\fbox{\rule{0pt}{0.46em}\hspace{0.46em}}}}
% what holds throughout the note: the ambient objects only (implicit and instance arguments), never a hypothesis
\newenvironment{conventions}{\par\addvspace{5pt}\noindent\textbf{Conventions.}\enspace\ignorespaces}{\par\addvspace{3pt}}
\makeatletter
\renewenvironment{proof}[1][Proof]{\par\pushQED{\qed}\normalfont\topsep0pt\partopsep0pt
  \trivlist\item[\hskip\labelsep\itshape #1.]\ignorespaces}{\popQED\endtrivlist\@endpefalse}
\makeatother

% generated by notes/filling-setup/shared/mkspec.py from shared/notation-148ab.json; do not edit.
% \input after data.tex: a macro that the build made from the spec is kept, the others are defined here (same TeX).
\providecommand{\fs}{\mathfrak{s}}                      % the shape (rho, kappa, k_0)
\providecommand{\sT}{\mathsf{T}}                        % a finite set of finite places
\providecommand{\cL}{\mathcal{L}}                       % a level
\providecommand{\AEf}{\mathbb{A}_{E,f}}
\providecommand{\AdE}{\mathbb{A}_{E}}
\providecommand{\hO}{\hat{O}_E}
\providecommand{\GL}{\mathrm{GL}}
\providecommand{\full}{\mathrm{full}}
\providecommand{\cut}{\mathrm{cut}}
\providecommand{\rh}[1]{\varrho(#1)}                    % Shape.ϱ
\providecommand{\kk}[1]{k(#1)}                          % Shape.k = #κ
\providecommand{\nn}[1]{n(#1)}                          % Shape.n = ϱ + 1
\providecommand{\rr}[1]{r(#1)}                          % Shape.r = 2ϱ + k + 1
\providecommand{\Ix}[1]{I(#1)}                          % Shape.BesselIndex
\providecommand{\Mx}[1]{M(#1)}                          % Shape.MiddleIndex
\providecommand{\ninf}[1]{f(#1)}                        % numInfPlaces E
\providecommand{\ET}[1]{E_{#1}}                         % LocProd T
\providecommand{\bS}[1]{b(#1)}                          % S.b
\providecommand{\NS}[1]{N(#1)}                          % S.N
\providecommand{\qS}[2]{q_{#2}(#1)}                     % S.levelPrime i
\providecommand{\piS}[1]{\bar\pi(#1)}                   % S.piBar
\providecommand{\vpS}[1]{\varpi(#1)}                    % S.ϖ
\providecommand{\ChS}[2]{\mathrm{Ch}_{#2}(#1)}          % S.chosenPlaces i
\providecommand{\KcolS}[1]{K^{\mathrm{col}}(#1)}        % S.columnLevel
\providecommand{\BsysS}[2]{\mathfrak{B}_{#2}(#1)}       % S.besselSys v
\providecommand{\mass}[2]{\operatorname{mass}(#1,#2)}   % mass E r
% Lib: forms, groups and frames (letters of the notes of unit 132)
\providecommand{\fr}[1]{\mathfrak{f}_{#1}}              % the frame vector of index a (e_a for a < rho, z for a = rho)
\providecommand{\Trg}[2]{T^{\mathrm{rat}}(#1;#2)}       % Shape.middleGram N: the form diag(1,...,1,-N)
\providecommand{\JT}[1]{J(#1)}                          % Shape.Form.splitGram
\providecommand{\GT}[1]{G(#1)}                          % Shape.Form.G
\providecommand{\Gone}[1]{G_{1}(#1)}                    % Shape.G₁
\providecommand{\UT}[1]{\mathcal{U}(#1)}                % besselU
\providecommand{\UoT}[1]{\mathcal{U}^{0}(#1)}           % besselU0
\providecommand{\HT}[1]{\mathcal{H}(#1)}                % besselH
\providecommand{\ST}[1]{\mathcal{S}(#1)}                % besselS
\providecommand{\Hl}[2]{\mathcal{H}_{#1}(#2)}           % GT l
\providecommand{\Ql}[2]{\mathcal{Q}_{#1}(#2)}           % unitriGroup l
\providecommand{\sd}[1]{\operatorname{sd}(#1)}          % simpleRootSum
\providecommand{\ZT}[2]{Z_{#1}(#2)}                     % Shape.Form.frameGram
\providecommand{\Zr}[1]{Z(#1)}                          % realFrameGram
\providecommand{\lmin}[2]{\mathrm{D}_{#1}(#2)}          % leadingMinor Z l
\providecommand{\weyl}[2]{w_{#1}(#2)}                   % weylElem l p
\providecommand{\gix}[1]{[#1]}                          % genIndex (the generalized index)
\providecommand{\den}[1]{\operatorname{den}(#1)}        % denomIndex
\providecommand{\Kf}[2]{K_{f}(#1,#2)}                   % Kf E n
\providecommand{\Gf}[2]{G_{f}(#1,#2)}                   % Gf E n
\providecommand{\volO}[1]{\operatorname{vol}(\hat O_{#1})} % volÔ E
\providecommand{\rdE}[1]{\operatorname{rd}_{#1}}        % rootDiscr E
\providecommand{\dE}[1]{d_{#1}}                         % discr E
\providecommand{\CC}{\mathbb{C}}
\providecommand{\RR}{\mathbb{R}}
\providecommand{\NN}{\mathbb{N}}
\providecommand{\vol}{\operatorname{vol}}
% ----- the definitions of chunks a and b of unit 148 (and, for this note, its own additions)
\providecommand{\Tlatg}[2]{T^{\mathrm{lat}}(#1;#2)}                  % latMiddleGram (𝔰, N)
\providecommand{\mbc}[2]{\Lambda^{\circ}(#1;#2)}                     % middleBaseChange (𝔰, p)
\providecommand{\blb}[2]{\Lambda(#1;#2)}                             % besselLatticeBasis (𝔰, p)
\providecommand{\fdisc}[2]{\mathsf{g}(#1;#2)}                        % frameDisc (N, Z)
\providecommand{\fpair}[1]{[\mathfrak{f}\mid #1\mathfrak{f}]}        % framePair (γ)
\providecommand{\ifpair}[1]{[\mathfrak{e}\mid #1\mathfrak{e}]}       % isotropicFramePair (γ)
\providecommand{\lmat}[1]{\mathrm{L}(#1)}                            % linkMatrix (γ)
\providecommand{\link}[1]{\operatorname{link}(#1)}                   % link (γ)
\providecommand{\damp}[2]{\operatorname{dp}_{#1}(#2)}                % dampedPower (k, r)
\providecommand{\hw}[3]{\Lambda^{\mathbb{C}}_{#1}(#2;#3)}            % headWeight (K, i, p)
\providecommand{\Sieve}[8]{\operatorname{Sieve}(#1,#2,#3,#4,#5,#6,#7,#8)} % SieveSystem (u, ι, n, J, ψ, En, Λ, β)
\providecommand{\fhw}[3]{\Lambda_{#1}(#2;#3)}                        % finHeadWeight (K, i, p)
\providecommand{\occ}[1]{c_{G}(#1)}                                  % Shape.openCellConst (𝔰)
\providecommand{\bandc}[1]{c_{\mathrm{band}}(#1)}                    % Shape.bandConst (𝔰)
\providecommand{\bsize}[1]{\operatorname{bs}(#1)}                    % besselSize (y)
\providecommand{\InCut}[2]{#1\in\operatorname{Cut}(#2)}              % InCut (y, b)
\providecommand{\frbox}[1]{\operatorname{FBox}(#1)}                  % frameBox (ρ)
\providecommand{\gdr}[1]{\mathsf{g}^{\mathbb{R}}(#1)}                % gramDiscReal (Z)
\providecommand{\cellm}[4]{\operatorname{cm}(#1;#2,#3,#4)}           % cellMass (y, deg, ρ, d)
\providecommand{\bcorr}[3]{\mathrm{J}(#1,#2;#3)}                     % besselCorrelation (μ, W, γ)
\providecommand{\SchurLe}[2]{\operatorname{Sch}(#1)\le #2}           % SchurLe (W, D)
\providecommand{\HasSB}[4]{\operatorname{SchB}(#1,#2;#3,#4)}         % HasSchurBound (W, rep, B, D)
\providecommand{\ncorr}[3]{\mathrm{J}^{\circ}(#1;#2,#3)}             % normCorrelation (y, W, p)
\providecommand{\BandL}[4]{\operatorname{BL}(#1,#2,#3,#4)}           % HasBandLimits (b, W, deg, J)
\providecommand{\CutPair}[4]{\operatorname{CP}(#1,#2,#3,#4)}         % IsCutPair (b, f, rep, W)
\providecommand{\BSys}[3]{\operatorname{BSys}(#1,#2,#3)}             % BesselSystem (𝔰, b, f)
\providecommand{\Fill}[2]{\operatorname{Fill}(#1,#2)}                % Filling (E, 𝔰)
\providecommand{\gS}[1]{g(#1)}                                       % Filling.g (S)
\providecommand{\Trat}[2]{T^{\mathrm{rat}}_{#1}(#2)}                 % Filling.ratMiddle (S, R)
\providecommand{\Tlat}[2]{T^{\mathrm{lat}}_{#1}(#2)}                 % Filling.latMiddle (S, R)
\providecommand{\GS}[2]{G_{#1}(#2)}                                  % Filling.G (S, R)
\providecommand{\GFin}[1]{G^{\mathrm{lat}}_{#1}}                     % Filling.GFin (S)
\providecommand{\GRes}[1]{\mathsf{G}(#1)}                            % Filling.GRes (S)
\providecommand{\Bc}[1]{\mathrm{Bc}(#1)}                             % Filling.BesselCoord (S)
\providecommand{\Cl}[1]{\mathrm{Cl}(#1)}                             % Filling.Classes (S)
\providecommand{\Zf}[2]{Z_{#1}(#2)}                                  % Filling.frameProducts (S, γ)
\providecommand{\KS}[1]{K(#1)}                                       % Filling.latStabilizer (S)
\providecommand{\redS}[1]{\mathrm{red}_{#1}}                         % Filling.reduction (S)
\providecommand{\KgS}[1]{K_g(#1)}                                    % Filling.principalLevel (S)
\providecommand{\redim}[2]{\overline{#2}^{\,#1}}                     % Filling.reducedImage (S, A)
\providecommand{\kapS}[2]{\kappa_{#2}(#1)}                           % Filling.kappa (S, v)
\providecommand{\bscale}[3]{\mathrm{sc}_{#1,#2}(#3)}                 % Filling.basisScale (S, v, i)
\providecommand{\sdc}[1]{\psi_{#1}}                                  % Filling.selfDualChar (S)
\providecommand{\fscale}[3]{\mathrm{fs}_{#1}(#2,#3)}                 % Filling.frameScaleSq (S, v, a)
\providecommand{\fsprod}[1]{\mathrm{fs}(#1)}                         % Filling.frameScaleProd (S)
\providecommand{\ucov}[1]{\nu(#1)}                                   % Filling.unipCovolume (S)
\providecommand{\siegC}[1]{c_{*}(#1)}                                % Filling.siegelConst (S)
\providecommand{\tdisc}[2]{\mathsf{g}_{#1}(#2)}                      % Filling.termDisc (S, γ)
\providecommand{\bsGL}[2]{\mathrm{Sc}_{#1,#2}}                       % Filling.basisScaleGL (S, v)
\providecommand{\latB}[1]{\Lambda(#1)}                               % Filling.latBasis (S)
\providecommand{\tolb}[1]{\mathrm{lb}_{#1}}                          % Filling.toLatBasis (S)
\providecommand{\radl}[1]{\mathrm{r}^{\mathbb{A}}_{#1}}              % Filling.ratToAdelic (S)
\providecommand{\ainf}[1]{\mathrm{a}^{\infty}_{#1}}                  % Filling.adelicToInf (S)
\providecommand{\afin}[1]{\mathrm{a}^{f}_{#1}}                       % Filling.adelicToFin (S)
\providecommand{\rinf}[1]{\mathrm{r}^{\infty}_{#1}}                  % Filling.ratToInf (S)
\providecommand{\rfin}[1]{\mathrm{r}^{f}_{#1}}                       % Filling.ratToFin (S)
\providecommand{\incH}[1]{\mathrm{inc}_{#1}}                         % Filling.inclHFin (S)
\providecommand{\hc}[2]{h_{#1}(#2)}                                  % Filling.h (S, c)
\providecommand{\pos}[4]{\mathrm{pos}_{#1}(#2,#3;#4)}                % Filling.termPosition (S, c, c', γ)
\providecommand{\terms}[4]{\mathrm{Tm}_{#1}(#2,#3,#4)}               % Filling.termsOver (S, c, s, c')
\providecommand{\clU}[2]{\mathcal{U}_{#1,#2}}                        % Filling.classUnip (S, c)
\providecommand{\clUo}[2]{\mathcal{U}^{0}_{#1,#2}}                   % Filling.classUnip₀ (S, c)
\providecommand{\bxs}[1]{\operatorname{bx}(#1)}                      % Filling.boxSize (S)
\providecommand{\IsBP}[3]{#3\in\operatorname{BP}(#1,#2)}             % Filling.IsBoxPartition (S, c, ω)
\providecommand{\besfn}[4]{W^{#1}_{#2}(#3)_{#4}}                     % Filling.besselFn (σ, S, y, i)
\providecommand{\besker}[4]{\mathcal{W}^{#1}_{#2}(#3;#4)}            % Filling.besselKernel (σ, S, γ, y)
\providecommand{\cdim}[1]{\mathrm{cd}(#1)}                           % Filling.columnDim (S)
\providecommand{\bdeg}[1]{\deg(#1)}                                  % Filling.besselDegree (S)
\providecommand{\mainE}[1]{\mathsf{d}(#1)}                           % Filling.mainEnergy (S)
\providecommand{\uord}[1]{\mathrm{uo}(#1)}                           % Filling.unipOrder (S)
\providecommand{\packC}[1]{a_{*}(#1)}                                % Filling.packetConst (S)
\providecommand{\besE}[1]{\mathsf{d}_{1}(#1)}                        % Filling.besselEnergy (S)
\providecommand{\tailS}[1]{\mathrm{tl}(#1)}                          % Filling.tailScale (S)
\providecommand{\codimF}[1]{\mathrm{cf}(#1)}                         % Filling.codimFactor (S)
\providecommand{\IsAdm}[1]{#1\text{ is admissible}}                  % Filling.IsAdmissible (S)
\providecommand{\cutoff}[2]{\beta_{#1,#2}}                           % Filling.classCutoff (S, c)
\providecommand{\amp}[5]{a^{#1}_{#2}(#3,#4;#5)}                      % Filling.amplitude (σ, S, c, c', γ)
\providecommand{\ampabs}[5]{|a|^{#1}_{#2}(#3,#4;#5)}                 % Filling.amplitudeAbs (σ, S, c, c', γ)
\providecommand{\ampsum}[6]{A^{#1}_{#2}(#3;#4,#5,#6)}                % Filling.amplitudeSum (σ, S, T, c, s, c')
\providecommand{\gid}[1]{\hat{g}(#1)}                                % Filling.gIdele (S)
\providecommand{\uchar}[2]{\psi^{\mathcal{U}}_{#1}(#2)}              % Filling.unipotentChar (S, u)
\providecommand{\Agr}[2]{#2\in\operatorname{Agr}(#1)}                % Filling.Agrees (S, s)
\providecommand{\IsStab}[2]{#2\in\operatorname{Stab}(#1)}            % Filling.IsStable (S, T)
\providecommand{\tlink}[2]{\mathrm{Tm}^{\mathrm{lk}}_{#1}(#2)}       % Filling.termsOfLink (S, Λ)
\providecommand{\clB}[2]{\mathrm{B}_{#1}(#2)}                        % Filling.classBasis (S, c)
\providecommand{\ioia}[2]{\operatorname{adj}_{#1}^{#2}}              % Filling.idOrInvAdjoint (S, b)
\providecommand{\rag}[2]{\mathrm{red}_{#1,#2}}                       % Filling.reduceAboveG (S, x)
\providecommand{\vfix}[3]{\mathrm{Fix}_{#1,#2}(#3)}                  % Filling.vectorFixer (S, w, d)
\providecommand{\FixN}[2]{\operatorname{fix}(#1)\ge #2}              % Filling.FixesNormAtLeast (s, x)
\providecommand{\fixset}[3]{\mathrm{Fix}_{#1}(#2,#3)}                % Filling.fixingSet (S, Sg, d)
\providecommand{\zfp}[1]{\mathrm{L}_{z}(#1)}                         % Filling.zFormPair (T)
\providecommand{\clink}[2]{\operatorname{link}^{\circ}(#1,#2)}       % commonLink (γ, γ')
\providecommand{\HCFV}[4]{\operatorname{CF}_{#4}(#1;#2,#3)}          % HasCommonFlagVector (γ, c, c', a)
\providecommand{\IsNF}[2]{#2\in\operatorname{NF}_{#1}}               % IsNormalForm (j, γ)
\providecommand{\ratG}[1]{J_{#1}}                                    % Filling.ratGram (S)
\providecommand{\lform}[1]{Q^{\pm}_{#1}}                             % Filling.linkForm (S)
\providecommand{\clC}[2]{\mathrm{C}_{#1}(#2)}                        % Filling.classCoord (S, c)
\providecommand{\clCX}[2]{\mathrm{C}^{\circ}_{#1}(#2)}               % Filling.classCoordX₀ (S, c)
\providecommand{\cw}[2]{\operatorname{cw}_{#2}(#1)}                  % Filling.countWeight (S, i)
\providecommand{\HasLev}[3]{#3\in\operatorname{Lv}_{#2}(#1)}         % Filling.HasLevel (S, j, γ)
\providecommand{\Cong}[5]{#4\sim_{#1;#2,#3}#5}                       % Filling.CongruentTerms (S, c, c', γ, γ')
\providecommand{\tdnorm}[2]{\|\mathsf{g}\|_{#1}(#2)}                 % Filling.termDiscNorm (S, γ)
\providecommand{\bxp}[1]{\operatorname{BxP}(#1)}                     % Filling.boxedPoints (S)
\providecommand{\IsNr}[2]{#2\in\operatorname{Near}(#1)}              % Filling.IsNear (S, γ)
\providecommand{\ChA}[2]{#2\in\operatorname{ChA}(#1)}                % Filling.CharsAgree (S, γ)
\providecommand{\fadm}[1]{(#1)_{\mathbb{A}_{f}}}                     % Filling.finAdeleMat (M)
\providecommand{\slb}[2]{\mathrm{m}_{#1}(#2)}                        % Filling.sublatticeBasis (H, V)
\providecommand{\sld}[2]{\operatorname{disc}_{#1}(#2)}               % Filling.sublatticeDisc (u, V)
\providecommand{\lmet}[1]{Q^{\mathrm{link}}_{#1}}                    % Filling.linkMetric (S)
\providecommand{\ldsq}[3]{\mathrm{ld}_{#1}(#2,#3)}                   % Filling.linkDiscSq (S, v, e)
\providecommand{\HIL}[3]{#3\in\operatorname{IL}(#1;#2)}              % Filling.HasIntegralLifts (c, γ, m)
\providecommand{\IdxB}[1]{\mathrm{Ib}(#1)}                           % Filling.IndexBelow (j)
\providecommand{\xrows}[1]{#1^{\mathrm{x}}}                          % Filling.xRows (Z)
\providecommand{\padl}[2]{\operatorname{pad}_{#1}(#2)}               % Filling.paddedLeading (j, Z)
\providecommand{\fmid}[1]{(#1\mathfrak{f})^{\mathrm{mid}}}           % Filling.frameMiddle (γ)
\providecommand{\latd}[2]{\operatorname{disc}_{#1,#2}}               % Filling.latticeDisc (S, c)
\providecommand{\fvi}[5]{\mathrm{o}_{#1,#2}(#3;#4,#5)}               % Filling.frameValueIndex (S, c, k, a, V)
\providecommand{\xri}[3]{\mathrm{i}_{#1,#2}(#3)}                     % Filling.xRowsIndex (S, j, k)
\providecommand{\sati}[2]{\mathfrak{d}_{#1}(#2)}                     % Filling.saturationIndex (S, k)
\providecommand{\pinv}[5]{\mathrm{pi}_{#1}(#2,#3;#4,#5)}             % Filling.pairInvariant (S, c, c', γ, γ')
\providecommand{\telem}[3]{\mathrm{el}_{#1}(#2;#3)}                  % Filling.termElem (S, c, p)
\providecommand{\mnorm}[3]{\mathrm{mn}_{#1}(#2;#3)}                  % Filling.minorNorm (S, j, γ₁)
\providecommand{\cgauge}[5]{\mathrm{cg}_{#1}(#2;#3,#4;#5)}           % Filling.classGauge (S, j, c, c', γ)
\providecommand{\IsNRep}[4]{#4\in\operatorname{NRep}_{#1}(#2,#3)}    % Filling.IsNearRep (S, c, c', γ)
\providecommand{\ipw}[4]{\mathrm{ipw}_{#1}(#2;#3,#4)}                % Filling.invPowWeight (S, c, n, p)
\providecommand{\termw}[7]{\mathrm{w}_{#1}(#2,#3;#4;#5,#6;#7)}       % Filling.termWeight (S, j, i, c, p, p', V)
\providecommand{\pairw}[6]{\mathrm{pw}_{#1}(#2;#3,#4;#5,#6)}         % Filling.pairWeight (S, j, c, c', γ, γ')
\providecommand{\IsLab}[5]{#5\in\operatorname{Lab}_{#1}(#2;#3,#4)}   % Filling.IsLabelledTerm (S, j, c, t, p)
\providecommand{\energy}[3]{\mathrm{En}_{#1}(#2;#3)}                 % Filling.energy (S, A, Ω)
\providecommand{\linkE}[5]{\mathrm{LE}^{#1}_{#2}(#3;#4,#5)}          % Filling.linkEnergy (σ, S, T, Ω', R)
\providecommand{\csize}[3]{\operatorname{cs}_{#1}(#2;#3)}            % Filling.classSize (S, j, x)
\providecommand{\bcut}[2]{\beta^{0}_{#1,#2}}                         % Filling.boxCutoff (S, c)
\providecommand{\schw}[4]{\mathrm{sw}_{#1}(#2;#3,#4)}                % Filling.schurWeight (S, c, γ, y)
\providecommand{\schi}[4]{\mathrm{SI}_{#1}(#2,#3;#4)}                % Filling.schurIntegral (S, c, c', γ')
\providecommand{\charT}[2]{\psi_{#1,#2}}                             % Filling.charOnPlaces (S, T)
\providecommand{\fgain}[3]{\beta_{#1,#2}(#3)}                        % Filling.frequencyGain (S, T, ξ)
\providecommand{\ratpt}[3]{\mathrm{pt}_{#1}(#2;#3)}                  % Filling.ratPoint (T, l, t)
\providecommand{\Fits}[6]{\operatorname{Fits}(#1;#2,#3,#4,#5,#6)}    % Filling.Fits (S, V, l, kw, Y, η)
\providecommand{\stdK}[1]{K^{\mathrm{std}}_{g}(#1)}                  % Filling.stdPrincipalLevel (S)
\providecommand{\tailv}[3]{\operatorname{tv}(#1,#2,#3)}              % tailVector (π, t, x)
\providecommand{\ratEq}[1]{\mathrm{rat}_{#1}}                        % Filling.Fits.ratEquiv (X)
\providecommand{\adEq}[1]{\mathrm{ad}_{#1}}                          % Filling.Fits.adelicEquiv (X)
\providecommand{\reH}[1]{\mathrm{re}_{#1}}                           % Filling.Fits.realHomeo (X)
\providecommand{\finEq}[1]{\mathrm{fin}_{#1}}                        % Filling.Fits.finiteEquiv (X)
\providecommand{\cutB}[3]{\mathrm{cb}_{#1}(#2;#3)}                   % Filling.Fits.cutBessel (X, a, g)
\providecommand{\Poin}[4]{\mathrm{P}_{#1}(#2,#3)(#4)}                % Filling.Fits.poincare (X, φ, Δ, g)
\providecommand{\Pmap}[1]{\mathrm{Q}_{#1}}                           % Filling.Fits.poincareMap (X)
\providecommand{\wbe}[1]{\mathsf{A}_{#1}}                            % Filling.Fits.weylBesselEquiv (h)
\providecommand{\colc}[1]{\mathrm{cc}_{#1}}                          % Filling.Fits.columnCoords (h)
\providecommand{\cgp}[1]{\mathrm{K}^{\circ}(#1)}                     % Level.congGroup (F)
\providecommand{\IsLat}[1]{#1\in\operatorname{LatLev}}               % IntegralPlace.IsLatticeLevel (F)
\providecommand{\IsStd}[1]{#1\in\operatorname{StdLev}}               % IntegralPlace.IsStdLevel (F)
\providecommand{\stdg}[2]{\Gamma_{#1}(#2)}                           % IntegralPlace.IsStdLevel.gauge (F, α)

% ----- hand macros of this note. Notions of other units are written as the notes of those units write them
% ----- (unit 132: notes/real-place-local-theorem; unit 141: notes/seesaw-identity; unit 152: notes/admissible-data).
\newcommand{\dd}{\mathrm{d}}
\newcommand{\cj}[1]{\overline{#1}}
\newcommand{\Nb}[2]{\mathcal{N}_{<#1}(#2)}            % unipotentBelow l
\newcommand{\Nf}[2]{\mathcal{N}_{#1}(#2)}             % floorUnipotent l
\newcommand{\modA}{\operatorname{mod}_A}              % the module of a place
\newcommand{\wo}[2]{w^{\circ}_{#1}(#2)}               % floorWeyl l c
\newcommand{\gL}[1]{g_{#1}}                           % the unit of a level
\newcommand{\BoxL}[1]{\operatorname{Box}(#1)}         % boxSet
\newcommand{\Posn}[1]{\operatorname{Posn}(#1)}        % positions
\newcommand{\Ball}[1]{\mathfrak{B}(#1)}               % the ball of a level
\newcommand{\EL}[1]{E(#1)}                            % the basis matrix of a level
\newcommand{\quadric}[2]{\delta^{(#2)}(#1)}           % floorQuadric
\newcommand{\Levels}[3]{\operatorname{Lev}(#1,#2,#3)} % Level ι T m
\newcommand{\HypB}[3]{\operatorname{Hyp}(#1,#2,#3)}   % HyperbolicBasis ι T m
\newcommand{\floorMeasure}[2]{\theta_{#2}(#1)}        % Level.floorMeasure
\newcommand{\stackMeasure}[1]{\mu^{#1}}               % Level.stackMeasure (indices a, i as subscripts)
\newcommand{\groupMeasure}[1]{\mathsf{M}(#1)}         % Level.groupMeasure
\newcommand{\upperMeasure}[2]{\mathsf{M}_{\ge #2}(#1)} % Level.upperMeasure
\newcommand{\IsNear}[3]{#3\in\operatorname{Near}_{#2}(#1)} % Level.IsNear
\newcommand{\floorModulus}[3]{\mathsf{d}_{#2}(#1;#3)} % Level.floorModulus
\newcommand{\Gauges}[1]{\operatorname{Gau}(#1)}       % Gauge F
\newcommand{\volumeExponent}[1]{\mathsf{s}(#1)}       % Level.volumeExponent
\newcommand{\uLev}[2]{K_{#2}(#1)}                     % unitaryLevel O_A g J(T)
\newcommand{\etab}[1]{\eta_{<#1}}                     % ascConv heisenbergMeasure j
\newcommand{\zetaA}[1]{\zeta_{A}(#1)}                 % idealZeta O_A
\newcommand{\Inp}[1]{\textup{(#1)}}
% the data D (Chapters III and V of the write-up)
\newcommand{\scx}[1]{x(#1)}
\newcommand{\scb}[1]{b(#1)}
\newcommand{\kD}[1]{k(#1)}
\newcommand{\fD}[1]{f(#1)}
\newcommand{\aD}[1]{a(#1)}
\newcommand{\hhD}[1]{h_2(#1)}
\newcommand{\hD}[1]{h(#1)}
\newcommand{\Dm}[1]{m(#1)}
\newcommand{\Dr}[1]{r(#1)}
\newcommand{\ED}[1]{E(#1)}
\newcommand{\wt}[3]{w_{#1,#2}(#3)}                    % Data.w D j σ
\newcommand{\nw}[2]{w_{#1}(#2)}                       % Data.nominalWeight
\newcommand{\rowPar}[3]{\mathrm{par}_{#1,#2}(#3)}     % Data.rowParameter
\newcommand{\parW}[3]{\mathrm{par}^{W}_{#1,#2}(#3)}   % Data.paramW
\newcommand{\Ridx}[1]{\underline{R}(#1)}              % Data.RowIndex
\newcommand{\SD}[1]{S(#1)}                            % Data.filling
\newcommand{\Ys}[2]{Y_{#2}(#1)}                       % Data.spaceY
\newcommand{\Stg}[2]{\mathrm{St}_{#2}(#1)}            % Data.Stage
\newcommand{\FormY}[2]{\mathcal{F}_{#2}(#1)}          % Data.FormY
\newcommand{\Col}[1]{\mathrm{Col}(#1)}                % Data.Column
\newcommand{\chcomb}[4]{\mathrm{C}^{(#4)}_{#1,#2}[#3]} % chainComb with the translated vacua
\newcommand{\chainY}[3]{\mathcal{L}^{(#3)}(#1,#2)}    % Data.liftChainY
\newcommand{\iot}[2]{\iota_{#1,#2}}                   % Data.ι
\newcommand{\Bmu}[3]{B^{#3}_{#1,#2}}                  % Data.Bμ
\newcommand{\gYlast}[1]{\mathcal{Y}_{\mathrm{last}}(#1)} % Data.lastGrading
\newcommand{\incl}[1]{i_{#1}}                         % Grading.inclX
\newcommand{\Uni}[1]{U(#1)}                           % Grading.unipotent
\newcommand{\Bes}[1]{B(#1)}                           % Grading.besselGroup
\newcommand{\coef}[2]{\langle #1,#2\rangle}           % algFormCoeff
% where each declaration of the concordance is treated: \cs{key}{place}, the place naming a \label or a lettered item
\cs{memunipotentiff}{Lemma~\ref{l:inclx}}
\cs{adelicToFinadelicEquivinclX}{Lemma~\ref{l:inclx}}
\cs{finiteEquivinclX}{Lemma~\ref{l:inclx}}
\cs{adelicToInfinclX}{Lemma~\ref{l:inclx}}
\cs{weylBesselEquiv}{Definition~\ref{d:colc}}
\cs{columnCoords}{Definition~\ref{d:colc}}
\cs{existslaststagemap}{Proposition~\ref{p:last}}
\cs{cardbesselCoordmulfactorialle}{Theorem F}
\cs{rowParametersubmem}{Lemma~\ref{l:rowpar}}
\cs{logfillingbesselDegreememIcc}{Theorem F}
\cs{charOnPlaces}{Definition~\ref{d:sievedefs}}
\cs{charOnPlacesevalPlaces}{Lemma~\ref{l:charT}}
\cs{frequencyGain}{Definition~\ref{d:sievedefs}}
\cs{ratPoint}{Definition~\ref{d:sievedefs}}
\cs{sumfrequenciesle}{Proposition~\ref{p:freq}}
\cs{onesidedpoisson}{Theorem C}
\cs{sieve}{Theorem D}
\cs{Fits}{Definition~\ref{d:fits}}
\cs{stdPrincipalLevel}{Definition~\ref{d:fits}}
\cs{JeqratGram}{Lemma~\ref{l:dict}}
\cs{ratEquiv}{Definition~\ref{d:equivs}}
\cs{adelicEquiv}{Definition~\ref{d:equivs}}
\cs{adelicEquivratEquiv}{Lemma~\ref{l:dict}}
\cs{realHomeo}{Definition~\ref{d:equivs}}
\cs{realHomeomul}{Lemma~\ref{l:dict}}
\cs{finiteEquiv}{Definition~\ref{d:equivs}}
\cs{finiteEquiveqadelicToFin}{Lemma~\ref{l:dict}}
\cs{reductionfiniteEquiviff}{Lemma~\ref{l:dict}}
\cs{memstdPrincipalLeveliff}{Lemma~\ref{l:dict}}
\cs{cutBessel}{Definition~\ref{d:poin}}
\cs{cutBesselratEquivmul}{Lemma~\ref{l:cbrat}}
\cs{poincarea}{Definition~\ref{d:poin}}
\cs{measurablecutBessel}{Proposition~\ref{p:unfold}}
\cs{tsumcutBesselclassUnip}{Lemma~\ref{l:classunip}}
\cs{tsumcutBesselle}{Theorem A}
\cs{integralpoincare}{Proposition~\ref{p:unfold}}
\cs{poincareMap}{Theorem A}
\cs{peterssonpoincare}{Theorem B}
\cs{tailVector}{Definition~\ref{d:tail}}
\cs{formtailVector}{Lemma~\ref{l:tail}}
\cs{continuoustailVector}{Proposition~\ref{p:slab}}
\cs{eventuallymassnearquadricle}{Proposition~\ref{p:slab}}
\cs{congGroup}{Definition~\ref{d:stdlevel}}
\cs{IsLatticeLevel}{Definition~\ref{d:stdlevel}}
\cs{IsStdLevel}{Definition~\ref{d:stdlevel}}
\cs{existsisLatticeLevel}{Theorem E}
\cs{gauge}{Definition~\ref{d:gauge}}
\cs{heisenbergMeasureeqfloorModulus}{Proposition~\ref{p:heis}}
\cs{ascConvunitaryLevel}{Corollary~\ref{c:unfoldU}}
\cs{unfoldbesselU}{Corollary~\ref{c:unfoldU}}
\cs{groupMeasureboxSetle}{Proposition~\ref{p:mass}}
\cs{lestackMeasureboxSet}{Proposition~\ref{p:stack}}
\cs{isHaarMeasuregroupMeasure}{Corollary~\ref{c:haar}}

\begin{document}
\color{ink}

{\raggedright
{\sffamily\bfseries\fontsize{7.2}{9}\selectfont\color{leancol}\addfontfeatures{LetterSpace=8}READBACK NOTE 44\par}
\vspace{5pt}
{\fontsize{21.5}{25.5}\selectfont\bfseries Poincar\'e series of cut Bessel functions: unfolding, one-sided Poisson
summation, the adelic sieve, standard levels\par}
\vspace{5pt}
{\fontsize{12.4}{16}\selectfont\itshape\color{inksoft} One unit of the Lean artifact \ArtifactName, read back into
mathematics\par}}
\vspace{8pt}
{\color{rulegrey}\hrule height 0.6pt}
\vspace{7pt}
\thispagestyle{empty}

\textbf{Summary.} Theorem~B of the write-up (\emph{enough forms}) rests on a second moment of Poincar\'e series on an
indefinite unitary group. This note treats four pieces of its proof that do not depend on each other. First, the
Poincar\'e series themselves: they are sums of Bessel functions at the infinite places, cut off along an arithmetic
unipotent group; they converge uniformly, they are automorphic forms, and their Petersson products unfold into the
\emph{amplitude sums} of the companion notes (Theorems A and~B). Second, a sieve: a tower of positive kernels on
$\CC^{f}\times\prod_{w\in\sT}E_w$, each the integral of the next over a common last entry, is summed over rational
points by Poisson summation used as an inequality, one entry at a time, and the loss is a factor close to~$1$
(Theorems C and~D). Third, at a finite place the lattice of the filling data carries the structure that the tower
method of the real-place unit needs, with exact masses (Theorem~E). Fourth, the degree and the number of coordinates
of the Bessel systems attached to admissible data are computed from Weyl's dimension formula (Theorem~F).

\statement{leancol}{\fontsize{9.6}{13.2}\selectfont\textbf{What this note is.} The artifact \ArtifactName{} is a
body of Lean code whose final theorems state the Riemann hypothesis and its analogue for Dirichlet $L$-functions.
That is a claim: its authors report that Lean's kernel accepts the code, a replay of that check by the owner of this
write-up is recorded elsewhere, and nothing was compiled for this note; no mathematician has read the proof through.
This note treats five files of one unit (the folder of the filling theorem): \UnitInNote{} of its \UnitMembers{}
declarations. Every statement was compared, argument by argument, with the Lean text. Proofs are complete at the
level of a research paper, with these exceptions, each said where it occurs: results of the other notes of the unit,
of other units and of the background library are quoted as \emph{Inputs}; the verification that an explicit record
satisfies the axioms of a level or of a gauge is described axiom by axiom and not carried out in every detail
(Theorem~E, Lemma~\ref{l:gauge}); and comparisons of exponents that the formal text leaves to automation are replaced
by explicit ones of ours. The Lean text read is \PinText; Appendix~A is the concordance of \ConcDecls{} declarations.}

\sect{The statements}{s:intro}
The objects are defined in Section~\ref{s:setting}. In Theorems A and~B, $S$ is filling data over a CM field~$E$
(Note~S of this unit, \emph{The filling data: Bessel subgroups, links of frames, Siegel's lemma and strong
approximation}), $V$ is a hermitian space with frame whose rational form is that of~$S$, and $X$ is the record of
the compatibilities between them (Setting~\ref{set:fits}).

\statement{arith}{\phantomsection\label{t:A}\textbf{Theorem A (Poincar\'e series converge and are forms).}\enspace
In Setting~\ref{set:fits}, let $E$ be stably unramified and $S$ admissible.
\begin{enumerate}[label=(\alph*),leftmargin=2.2em,itemsep=1pt,topsep=2pt]
\item For every pair $a=(c,i)\in\Cl{S}\times\Bc{S}$ there is a real number $M\ge0$ such that
\[\sum_{\gamma\in\Gamma(V)}\bigl|\cutB{X}{a}{\gamma g}\bigr|\le M\qquad\text{for all $g\in G(V)$.}\]
\item The formal text fixes a $\CC$-linear map $\Pmap{X}$ from $\CC^{\Cl{S}\times\Bc{S}}$ to the space of the forms on
$V$ of type $\epsilon_l$ that are fixed by the level $\stdK{S}$, such that for all coefficient vectors $\phi$, with
entries $\phi_a$ indexed by the pairs~$a$, and all points $g\in G(V)$
\[\Pmap{X}(\phi)(g)\;=\;\Poin{X}{\phi}{\Gamma(V)}{g}\;=\;\sum_{a}\phi_a\sum_{\gamma\in\Gamma(V)}\cutB{X}{a}{\gamma g}.\]
\end{enumerate}}
\leanline{\lean{tsumcutBesselle}\sep\lean{poincareMap}}

\statement{arith}{\phantomsection\label{t:B}\textbf{Theorem B (the Petersson product of two Poincar\'e series).}\enspace
In Setting~\ref{set:fits}, let $E$ be stably unramified and $S$ admissible. Let $a=(c,i)$ and $a'=(c',i')$ be pairs
in $\Cl{S}\times\Bc{S}$, let $s\in\GRes{S}$, and let $k\in\Kf{E}{\Ix{\fs}}$ be an integral finite point whose reduction
at every chosen place $w\in\ChS{S}{j}$, $j\in\{0,1\}$, is the component $s_{j,w}$ of~$s$. Then, with $e_a$ the unit
vector at~$a$, $R(k)$ the right translation by~$k$, and $\langle\,,\rangle$ the Petersson product, conjugate
linear in the first variable,
\[\bigl\langle R(k)\,\Pmap{X}(e_{a'}),\;\Pmap{X}(e_a)\bigr\rangle\;=\;
\vol\bigl(\stdK{S}\bigr)\cdot\ampsum{\full}{S}{\GS{S}{E}}{c}{s}{c'}_{(i,i')} .\]
Here $\ampsum{\full}{S}{\GS{S}{E}}{c}{s}{c'}\in\CC^{\Bc{S}\times\Bc{S}}$ is the amplitude sum of Note~S for the full
Bessel functions over all rational elements that are terms over $(c,s,c')$, and $\vol$ is the Haar measure of the
finite adelic unitary group that gives its integral points mass~$1$.}
\leanline{\lean{peterssonpoincare}}

In Theorems C and~D, $\sT$ is a finite set of finite places of~$E$, $\ET{\sT}=\prod_{w\in\sT}E_w$, and a
\emph{sieve system} is a tower of kernels in the sense of Note~S, recalled in Input~\Inp{S4}.

\statement{arith}{\phantomsection\label{t:C}\textbf{Theorem C (one-sided Poisson summation along a tower).}\enspace
Let $S\in\Fill{E}{\fs}$ be admissible filling data, $\sT$ a finite set of finite places of~$E$, and $\mathcal{A}$ a
finite set of units of $E$ all of which are $\sT$-units; assume that $\gS{S}^{-1}$ is a $\sT$-integer. Put
$C=\CC^{f}\times\ET{\sT}$, the product over the infinite places and over~$\sT$, with its Haar measure, and
$U=(\CC^{\times})^{f}\times\ET{\sT}^{\times}$. Let $\mathsf{n}\in\NN$, let $J_i(p,q)\in\CC$ and
$\Lambda_i(p)\in\RR$ be given for $i\in\NN$ and sequences $p,q\colon\NN\to U$, let $\mathrm{En}\in\RR$ and
$\beta\colon C\to\RR$, and suppose that \[\Sieve{u}{\iota}{\mathsf{n}}{J}{\psi}{\mathrm{En}}{\Lambda}{\beta}\]
holds on $C$, where $u\colon C\to U$ replaces every coordinate that is not a unit by~$1$,
$\iota(p)=\bigl((p_\tau^{-1})_\tau,\ \gS{S}^{-1}p_{\sT}^{-1}\bigr)$ and
$\psi(x)=\prod_\tau\exp(x_\tau-\bar x_\tau)\cdot\charT{S}{\sT}(x_{\sT})$. Let $\mathsf{C}$ and $\mathsf{e}$ be real
numbers with $\mathsf{e}\ge0$, and suppose that $\Lambda_i(p)\ge0$ for all $i$,~$p$ and that the following
\emph{frequency bound} holds: for every $a\in\NN$, every family $c=(c_\tau)$ of complex numbers with
$|c_\tau|^{-1}\le|\kapS{S}{\tau}|/\gS{S}$ for all~$\tau$, and every finite set $\mathcal{X}$ of non-zero
$\sT$-integers of~$E$,
\[\sum_{y\in\mathcal{X}}\beta\Bigl(\ratpt{\sT}{(c_\tau/\fscale{S}{\tau}{a})_\tau}{y}\Bigr)\le\mathsf{e}.\]
For $i\in\NN$, $p\in\mathcal{A}^{i}$ and $t\in E^{\times}$ let $\mathrm{hd}_i(p,t)\colon\NN\to U$ be the sequence
\begin{gather*}\mathrm{hd}_i(p,t)_a=u\Bigl(\ratpt{\sT}{(\fscale{S}{\tau}{a})_\tau}{(p\,|\,t)_a}\Bigr),\\
(p\,|\,t)_a=p_a\ (a<i),\qquad (p\,|\,t)_i=t,\qquad (p\,|\,t)_a=1\ (a>i),\end{gather*}
and put $c_a=\volO{E}\big/\prod_\tau\fscale{S}{\tau}{a}^{2}$. If
\[\Bigl(\sum_{p\in\mathcal{A}^{i}}\Lambda_i\bigl(\mathrm{hd}_i(p,1)\bigr)\Bigr)^{2}\le\mathsf{C}\prod_{a<i}c_a
\qquad\text{for all }i\le\mathsf{n},\]
then for all $i\le\mathsf{n}$ and all subsets $\mathcal{V}\subseteq\mathcal{A}$
\[\operatorname{Re}\sum_{t\in\mathcal{V}}\ \sum_{p,q\in\mathcal{A}^{i}}J_{i+1}\bigl(\mathrm{hd}_i(p,t),\mathrm{hd}_i(q,t)\bigr)
\;\le\;\Bigl(\prod_{a\le i}c_a\Bigr)\bigl(\mathrm{En}+(2i+1)\,\mathsf{C}\,\mathsf{e}\bigr).\]}
\leanline{\lean{onesidedpoisson}}

\statement{arith}{\phantomsection\label{t:D}\textbf{Theorem D (the sieve).}\enspace
Let $E$ carry Weil data, let $S\in\Fill{E}{\fs}$ be admissible filling data, $\sT$ a finite set of finite places of~$E$
that is stable under complex conjugation, and $\mathcal{A}$ a finite set of units of $E$ all of which are $\sT$-units;
assume that $\gS{S}^{-1}$ is a $\sT$-integer, and let $\mathcal{A}'\subseteq\mathcal{A}$. Write $\varrho=\rh{\fs}$,
$n=\nn{\fs}$, $k=\kk{\fs}$, $f=\ninf{E}$, $b=\bS{S}$. Let $J'_i(p,q)\in\CC$, for $i\in\NN$ and sequences
$p,q\colon\NN\to\ET{\sT}^{\times}$, and $e\in\RR$ be such that
\[\Sieve{u_{\sT}}{\iota_{\sT}}{\varrho}{J'}{\charT{S}{\sT}}{e}{\fhw{n^{6}}{\,\cdot\,}{\,\cdot\,}}{\fgain{S}{\sT}{\,\cdot\,}}\]
holds on $\ET{\sT}$, where $u_{\sT}$ replaces a non-unit by~$1$ and $\iota_{\sT}(p)=\gS{S}^{-1}p^{-1}$. For
$p\in\mathcal{A}^{\varrho}$, $t\in\mathcal{A}$ and an infinite place $\tau$ let $x_\tau(p,t)$ be the sequence of
the non-zero complex numbers $\fscale{S}{\tau}{a}\cdot\bigl((p\,|\,t)_a\bigr)_\tau$, $a\in\NN$, and $y(p,t)$ the
sequence of the images of the $(p\,|\,t)_a$ in $\ET{\sT}^{\times}$. Then, with $J^{\tau}$ the band-limit forms and
$\deg_\tau(S)$ the degree of the Bessel system $\BsysS{S}{\tau}$,
\begin{multline*}
\operatorname{Re}\sum_{t\in\mathcal{A}'}\ \sum_{p,q\in\mathcal{A}^{\varrho}}\ \prod_\tau J^{\tau}_{n}\bigl(x_\tau(p,t),x_\tau(q,t)\bigr)\cdot
J'_{n}\bigl(y(p,t),y(q,t)\bigr)\\
\le\;\prod_{a<n}\frac{\volO{E}}{\prod_\tau\fscale{S}{\tau}{a}^{2}}\cdot
\Bigl(e\prod_\tau\frac{1+b^{-10}}{\deg_\tau(S)}+\frac{\bdeg{S}^{-1}}{b^{6kf+10}}\Bigr).
\end{multline*}}
\leanline{\lean{sieve}}

In Theorem~E, $A$ is a finite product of non-archimedean local fields with an involution, $O_A$ its ring of
integers, and a \emph{level} is the structure of the tower method (unit~132, recalled in Section~\ref{s:setting}).

\statement{arith}{\phantomsection\label{t:E}\textbf{Theorem E (the lattice level at a finite place).}\enspace
Let $A$ be a semilocal field with an involution and an integral place structure, $\fs$ a shape, $\Omega$ a finite set
and $\mathfrak{m}$ a natural number with $\kk{\fs}=|\Omega|+2\mathfrak{m}+1$. Let $N\in O_A$ with $\bar N=N$ be such that
the split Gram matrix $\JT{\Tlatg{\fs}{N}}$ is a hermitian unit, and let $g$ be a unit $g$ of $A$ with $g\in O_A$ and
$\bar g=g$. Then there is a level $\cL\in\Levels{\Omega}{\Tlatg{\fs}{N}}{\mathfrak{m}}$ that is a lattice level,
\m{\IsLat{\cL}} (Definition~\ref{d:stdlevel}), with unit $\gL{\cL}=g$.}
\leanline{\lean{existsisLatticeLevel}}

Section~\ref{s:level} adds the measure identities of a standard level: the Heisenberg measures of the congruence
subgroups (Proposition~\ref{p:heis}) and the two-sided bound
$\vol(O_A^{\times})^{\varrho+1}\le\groupMeasure{\cL}(\BoxL{\cL})\le1+2^{-n}$ for the mass of the integral group
(Propositions \ref{p:mass} and~\ref{p:stack}).

\statement{arith}{\phantomsection\label{t:F}\textbf{Theorem F (the Bessel systems of admissible data).}\enspace
Let $P$ be parameters, $D$ admissible data for~$P$, and $\SD{D}$ the filling data of~$D$ (Note~G of this unit). Write
$k=\kD{D}$, $f=\fD{D}$, $m=\Dm{D}$, $r=\Dr{D}$.
\begin{enumerate}[label=(\alph*),leftmargin=2.2em,itemsep=1pt,topsep=2pt]
\item $\displaystyle\#\Bc{\SD{D}}\cdot\bigl((k-1)!\bigr)^{f}\;\le\;\bigl(\scb{D}^{5/4}\bigr)^{kf}$.
\item With $\nw{D}{\tau}$ the nominal weight of the infinite place~$\tau$,
\begin{gather*}\sum_\tau m\,(k-1)\log\nw{D}{\tau}\;\le\;\log\bdeg{\SD{D}}+f\sum_{l=k-1}^{r-1}\log l!\\
\le\;\sum_\tau m\,r\,\Bigl(\log\nw{D}{\tau}+\frac1{\scx{D}}\Bigr),\end{gather*}
the sums over the $f$ infinite places of~$\ED{D}$.
\end{enumerate}}
\leanline{\lean{cardbesselCoordmulfactorialle}\sep\lean{logfillingbesselDegreememIcc}}

Theorems A and~B are used by the notes on the filling theorem (Note~F) and on the seesaw identity (unit~141);
Theorem~D is the sieve of the main-term estimate of Note~F; Theorem~E and the masses feed Notes T and~G; Theorem~F is
used when the filling data of $D$ are shown to be admissible (Note~G). The file of Theorem~F also holds the last
step of the seesaw argument (Proposition~\ref{p:last}), which belongs to unit~141 and is stated here because it is a
member of this unit.

\sect{The setting}{s:setting}
\begin{conventions}\phantomsection\label{c:conventions}
$E$ is a CM field with a fixed CM type; $\tau$ runs over its $f=\ninf{E}$ infinite places, and $x_\tau\in\CC$ is the
image of $x\in E$ under the embedding of the CM type at~$\tau$; $w$ is a finite place, $E_w$ the completion, $\F_w$
the residue field; $\AdE=\CC^{f}\times\AEf$ is the ring of adeles, $\AEf$ the finite adeles, $\hO$ their integral
elements; $\rdE{E}$ is the root discriminant. $\fs=(\varrho,\kappa,k_0)$ is a shape, with $\rh{\fs}=\varrho$,
$\kk{\fs}=|\kappa|$, $\nn{\fs}=\varrho+1$, $\rr{\fs}=2\varrho+|\kappa|+1$ and the index set $\Ix{\fs}$ of the basis
$e_0,\dots,e_{\varrho-1}$, $x_j$ ($j\in\kappa$), $z$, $e^*_{\varrho-1},\dots,e^*_0$. In Sections \ref{s:poincare}
and~\ref{s:columns} the index set $\Ix{\fs}$ carries signs, a frame for them is fixed, and $S$, $V$, $l$,
$\mathsf{k}$, $\mathcal{Y}$, $\eta$, $X$ are as in Setting~\ref{set:fits}. In Section~\ref{s:level}, $A$ is a
commutative ring with an involution $x\mapsto\bar x$, a topology and a Haar measure; from Definition~\ref{d:stdlevel}
on it is a semilocal field; $\Omega$ is a finite set, $\mathfrak{m}$ a natural number, $T$ an $\fs$-form over $A$
whose split Gram matrix is a hermitian unit, $P$ the place structure of~$A$, $\cL$ a level; and, as in the
\texttt{variable} lines of the formal file, $\pi$ is a unit of~$A$, $N$ and $t$ are elements of $A$, $x$ is a vector
in $A^{\kappa}$, $g$ a unit of~$A$. The letters $k$, $s$, $i$, $c$, $M$ are used for group elements, indices and
constants as each statement says. Integrals of functions with values in $[0,\infty]$ are Lebesgue integrals; sums of
such functions over infinite index sets are unconditional. Subtraction of natural numbers in exponents and indices is
truncated at~$0$.
\end{conventions}

\subhead{Recollections from Note~S.}\phantomsection\label{s:recall}
These objects are defined in Note~S of this unit, where their Lean names are given; we use its notation. A
\emph{$\fs$-form} over a commutative ring $R$ with involution is a square matrix $T$ on the indices of the $x_j$ and
of~$z$; $\JT{T}$ is its split Gram matrix on $\Ix{\fs}$ (the $e_a$ are paired with the $e^*_a$, the middle block
is~$T$), $\GT{T}$ the group of the invertible matrices $g$ with $g^{*}\JT{T}g=\JT{T}$, and $\UT{T}$, $\HT{T}$,
$\ST{T}$ its Bessel unipotent subgroup, the subgroup that fixes all $e_a$, $e^*_a$ and~$z$, and the Bessel subgroup
$\ST{T}=\UT{T}\HT{T}$; $\Hl{l}{T}$ fixes the basis vectors of floor less than~$l$ (so $\Hl{\varrho}{T}$ is the
unitary group of $\langle x_j,z\rangle$ and $\Hl{n}{T}=\HT{T}$), $\Nb{l}{T}$ is the unipotent radical of the
stabiliser of the first $l$ steps of the flag $\langle e_0\rangle\subset\langle e_0,e_1\rangle\subset\cdots$, and
$\Nf{l}{T}=\Hl{l}{T}\cap\Nb{l+1}{T}$. The frame is $\fr{0},\dots,\fr{\varrho}$ with $\fr{a}=e_a$ for $a<\varrho$ and
$\fr{\varrho}=z$; $\sd{u}$ is the sum of the entries of $u$ at the places $(\fr{a},\fr{a+1})$. For $N\in R$,
$\Trg{\fs}{N}$ is the form $\mathrm{diag}(1,\dots,1,-N)$ and $\Tlatg{\fs}{N}$ its Gram matrix in the lattice basis;
for a unit $p$ of $R$, $\blb{\fs}{p}$ is the matrix of that basis ($x_{k_0}$ replaced by
$x_{k_0}/\bar p+z/(p\bar p)$). $\Gone{\fs}$ is the group of $\Trg{\fs}{1}$ over~$\CC$.

\emph{Filling data} $S\in\Fill{E}{\fs}$ consist of a prime $\NS{S}$, primes $\qS{S}{0},\qS{S}{1}$ with product
$\gS{S}$, a finite idele $\piS{S}$ with $\piS{S}\,\overline{\piS{S}}=\NS{S}$, a generator $\vpS{S}$ of the different
with $\overline{\vpS{S}}=-\vpS{S}$, a real number $\bS{S}$, finite sets $\ChS{S}{0},\ChS{S}{1}$ of \emph{chosen
places} above the two primes, a level $\KcolS{S}$ of the definite space on~$\kappa$, and a Bessel system
$\BsysS{S}{\tau}$ at every infinite place, with the size conditions of Note~S; $S$ may be \emph{admissible}
(the standing hypotheses of Note~S). From $S$: the forms $\Trat{S}{R}=\Trg{\fs}{\NS{S}}$ and
$\Tlat{S}{R}=\Tlatg{\fs}{\NS{S}}$ over a ring~$R$; the groups $\GS{S}{R}=\GT{\Trat{S}{R}}$ and
$\GFin{S}=\GT{\Tlat{S}{\AEf}}$; the lattice basis $\latB{S}=\blb{\fs}{\piS{S}}$ and the isomorphism
$\tolb{S}\colon\GS{S}{\AEf}\to\GFin{S}$, $x\mapsto\latB{S}^{-1}x\,\latB{S}$; the integral subgroup $\KS{S}$ of
$\GFin{S}$, the finite group $\GRes{S}=\prod_{j}\prod_{w\in\ChS{S}{j}}\GL_{\Ix{\fs}}(\F_w)$, the reduction
$\redS{S}\colon\KS{S}\to\GRes{S}$ and its kernel $\KgS{S}$; the diagonal map
$\radl{S}\colon\GS{S}{E}\to\GS{S}{\AdE}$; the map $\ainf{S}\colon\GS{S}{\AdE}\to\prod_\tau\Gone{\fs}$, conjugation
of the component at $\tau$ by the diagonal matrix $\bsGL{S}{\tau}$ of the scales; the map
$\afin{S}\colon\GS{S}{\AdE}\to\GFin{S}$, $\tolb{S}$ of the finite part; $\rinf{S}=\ainf{S}\circ\radl{S}$ and
$\rfin{S}=\afin{S}\circ\radl{S}$; the inclusion $\incH{S}$ of the finite adelic unitary group $\Gf{E}{\kappa}$ of the
definite space into $\GFin{S}$; the set $\Cl{S}$ of the double cosets of $\Gf{E}{\kappa}$ for the rational points and
$\KcolS{S}$, with $\hc{S}{c}=\incH{S}(\text{the chosen representative of }c)$; the position
$\pos{S}{c}{c'}{\gamma}=\hc{S}{c}^{-1}\rfin{S}(\gamma)\hc{S}{c'}$ of a rational $\gamma\in\GS{S}{E}$ and the set
$\terms{S}{c}{s}{c'}$ of the \emph{terms over $(c,s,c')$}, the $\gamma$ whose position lies in $\KS{S}$ with reduction
$s\in\GRes{S}$; the arithmetic unipotent group $\clU{S}{c}$ of the $u\in\UT{\Trat{S}{E}}$ with
$\hc{S}{c}^{-1}\rfin{S}(u)\hc{S}{c}\in\KgS{S}$; the cutoff $\cutoff{S}{c}\colon\prod_\tau\Gone{\fs}\to[0,\infty]$; the
set $\Bc{S}=\prod_\tau I_\tau(S)$ of the tuples of coordinates of the Bessel systems, and for a side
$\sigma\in\{\full,\cut\}$ the product Bessel function $\besfn{\sigma}{S}{y}{i}=\prod_\tau W^{\sigma}_{S,\tau}(y_\tau)_{i_\tau}$,
$y\in\prod_\tau\Gone{\fs}$, $i\in\Bc{S}$, where $W^{\sigma}_{S,\tau}\colon\Gone{\fs}\to\CC^{I_\tau(S)}$ are the two
functions of $\BsysS{S}{\tau}$; the kernel
$\besker{\sigma}{S}{\gamma}{y}_{(i,i')}=\besfn{\sigma}{S}{y}{i}\,\overline{\besfn{\sigma}{S}{\rinf{S}(\gamma)^{-1}y}{i'}}$,
the amplitude
$\amp{\sigma}{S}{c}{c'}{\gamma}=\int\cutoff{S}{c}(y)\,\cutoff{S}{c'}(\rinf{S}(\gamma)^{-1}y)\,\besker{\sigma}{S}{\gamma}{y}\,\dd y$
and, for a set $\Xi$ of rational elements, the amplitude sum
$\ampsum{\sigma}{S}{\Xi}{c}{s}{c'}=\sum_{\gamma\in\Xi\cap\terms{S}{c}{s}{c'}}\amp{\sigma}{S}{c}{c'}{\gamma}$; the
additive character $\sdc{S}(x)=\psi_f(-x/\vpS{S})$ of $\AEf$, where $\psi_f$ is the standard character; the real
numbers $\kapS{S}{\tau}=\operatorname{Re}\bigl(\gS{S}\vpS{S}_\tau/2\pi i\bigr)$ and
$\fscale{S}{\tau}{a}=\kapS{S}{\tau}^{2(\varrho-a)}/\NS{S}$ for $a\in\NN$; the degree
$\bdeg{S}=\prod_\tau\deg_\tau(S)$, the product of the degrees of the Bessel systems. Finally, for a real number $\rho$ and natural numbers $j$, $K$, $i$, the damped power is $\damp{j}{\rho}=\rho^{j}\min(1,\rho)^{6}$, and the head weights are
\begin{gather*}\hw{K}{i}{p}=\prod_{l<i}\damp{K}{|p_0\cdots p_l|^{-1}}\quad(p\colon\NN\to\CC^{\times}),\\
\fhw{K}{i}{p}=\prod_{l<i}\bigl(\gix{p_0\cdots p_l}^{1/2}\bigr)^{K}\quad(p\colon\NN\to A^{\times}),\end{gather*}
the factor at $l$ of the second being $0$ if $p_0\cdots p_l\notin O_A$; here, for a semilocal field~$A$,
$\gix{\xi}$ is the generalized index (the inverse of the module of $\xi$; it is $[O_A:\xi O_A]$ for $\xi\in O_A$) and
$\den{\xi}=[O_A:O_A\cap\xi^{-1}O_A]$.

\statement{ana}{\textbf{Inputs from the other notes of this unit.}\enspace
\Inp{S1} (Note~S) For admissible $S$: $1\le|\kapS{S}{\tau}|\le\bS{S}^{23}$ and
$|\kapS{S}{\tau}|=\gS{S}\,|\vpS{S}_\tau|/2\pi$ for every~$\tau$; $\NS{S}\le\bS{S}^{60n}$;
$\prod_\tau\bandc{\fs}\,\bS{S}^{7k}\bigl(|\kapS{S}{\tau}|^{2\varrho+1}/\gS{S}\NS{S}\bigr)^{k-n^{7}}\le\bdeg{S}^{-1}$.
For all filling data: $f\le\bS{S}/1000$, $k\le\bS{S}$, $\rdE{E}\le\bS{S}$, $1000\le\bS{S}$, $1000\,(2\varrho+2)^{9}\le k$
and $1000\,(2\varrho+2)\le k$.

\Inp{S2} (Note~S) For $a\in\AEf$: $\sdc{S}(ra)=1$ for all $r\in\hO$ if and only if $a\in\hO$.

\Inp{S3} (Note~S) For a commutative ring $R$ with involution, a unit $p$ and $N=p\bar p$:
$\blb{\fs}{p}^{*}\,\JT{\Trg{\fs}{N}}\,\blb{\fs}{p}=\JT{\Tlatg{\fs}{N}}$. If $R'\subseteq R$ is a subring, $N\in R'$ and
$\bar N=N$, then $\Tlatg{\fs}{N}$ is unimodular over~$R'$: hermitian, with entries and inverse over~$R'$.

\Inp{S4} (Note~S; sieve systems) Let $C$ be a commutative ring with a measure, $U$ a topological space,
$u\colon C\to U$, $\iota\colon U\to C$, $\mathsf{n}\in\NN$, $J_i(p,q)\in\CC$ and $\Lambda_i(p)\in\RR$ for $i\in\NN$
and sequences $p,q\colon\NN\to U$, $\psi\colon C\to\CC$, $\mathrm{En}\in\RR$ and $\beta\colon C\to\RR$. The
statement $\Sieve{u}{\iota}{\mathsf{n}}{J}{\psi}{\mathrm{En}}{\Lambda}{\beta}$ means: $\psi(0)=1$; for $i\le\mathsf{n}$, all
$p,q$ and $\xi\in C$ the function $c\mapsto\psi(\xi c)J_{i+1}(p[i\mapsto u(c)],q[i\mapsto u(c)])$ is integrable,
with integral $J_i(p,q)$ if $\xi=0$ and of absolute value at most $\beta(\xi)\Lambda_i(p)\Lambda_i(q)$ if $\xi$ is a
unit; for finitely many sequences with the same entry at~$\mathsf{n}$ the matrix of the values of
$J_{\mathsf{n}+1}$ is positive semidefinite; $J_{\mathsf{n}+1}$ is continuous; $|J_0|\le\mathrm{En}$; and for
$i<\mathsf{n}$, if $\xi=\iota(p_i)-\iota(q_i)$ is a unit then $|J_{i+1}(p,q)|\le\beta(\xi)\Lambda_{i+1}(p)\Lambda_{i+1}(q)$.
Here $p[i\mapsto x]$ is $p$ with its entry $i$ replaced by~$x$. Let this hold, let $i\le\mathsf{n}$ and let
$(a_p)_p$ be a finite family of sequences. (a) $F(c)=\sum_{p,q}J_{i+1}(a_p[i\mapsto u(c)],a_q[i\mapsto u(c)])$ is
integrable, $\int F=\sum_{p,q}J_i(a_p,a_q)$, and $|\int\psi(\xi c)F(c)\,\dd c|\le\beta(\xi)(\sum_p\Lambda_i(a_p))^{2}$
for every unit~$\xi$. (b) If $C$ is first countable and $(b_p(x))_p$ is a finite family of sequences depending on
$x\in C$ that share their entry~$i$ for each $x$, then $\sum_{p,q}J_{i+1}(b_p(x),b_q(x))$ is a non-negative real
number, and as a function of $x$ it is lower semicontinuous at every point where all $b_p$ are continuous.
(c) (products) If the measure of $C$ is $s$-finite and, for each infinite place~$\tau$, a system
$\Sieve{u_{\CC}}{\iota'}{\mathsf{n}}{J^{\tau}}{\psi_{\CC}}{\mathrm{En}_\tau}{\Lambda'}{\gamma_\tau}$ holds on~$\CC$,
with $u_{\CC}$ replacing $0$ by~$1$ and $\psi_{\CC}(c)=\exp(c-\bar c)$, then the product system holds on
$\CC^{f}\times C$: with $u$, $\iota$ taken factor by factor, the forms, the characters, the constants
$\mathrm{En}$, the weights $\Lambda$ and the functions $\beta$ multiplied.

\Inp{S5} (Note~S; band limits) For every filling data $S$ and every $\tau$, the Bessel system $\BsysS{S}{\tau}$ carries
forms $J^{\tau}$ with
$\Sieve{u_{\CC}}{\xi\mapsto\xi^{-1}}{\varrho}{J^{\tau}}{\psi_{\CC}}{(1+b^{-10})/\deg_\tau(S)}{\hw{n^{6}}{\,\cdot\,}{\,\cdot\,}}{\beta_{\CC}}$
on $\CC$, where $b=\bS{S}$ and $\beta_{\CC}(\xi)=\bandc{\fs}\,b^{2000n^{9}}\,\damp{k-n^{7}}{|\xi|^{-1}}$, and
$\bandc{\fs}>0$.

\Inp{S6} (Note~S; lattices) Let $E$ be stably unramified. For a subspace $L$ of $E^{2}$ let $\sld{1}{L}>0$ be
the discriminant of the lattice $L\cap\hO^{2}$ of Note~S (formed with the identity matrix). Then
$\sum_{L}\sld{1}{L}^{-3}\le(200\cdot2f\cdot\rdE{E})^{3f}$, the sum over the lines $L\subset E^{2}$. If $y\in E$,
$L$ is the line spanned by $(1,y)$ and $w\in\AEf^{\times}$ is such that $w$ and $yw$ lie in~$\hO$, then
$\sld{1}{L}\le\prod_\tau(1+|y_\tau|^{2})\cdot\gix{w}$.

\Inp{B1} (Note~B, \emph{Box partitions of arithmetic unipotent groups: frame counts, the covolume, amplitude sums}) For
admissible $S$, stably unramified~$E$, a class $c$ and $y\in\prod_\tau\Gone{\fs}$:
$\sum_{u\in\clU{S}{c}}\cutoff{S}{c}(\rinf{S}(u)\,y)=1$.

\Inp{B2} (Note~B) For $u\in\clU{S}{c}$, a side $\sigma$ and all $y$, $i$:
$\besfn{\sigma}{S}{\rinf{S}(u)y}{i}=\besfn{\sigma}{S}{y}{i}$.

\Inp{B3} (Note~B) Let $E$ be stably unramified and $S$ admissible, $\Xi$ a set of rational elements, $c,c'$ classes,
$s\in\GRes{S}$ and $\sigma$ a side. There is $M\ge0$ with
$\sum_{\gamma\in\Xi\cap\terms{S}{c}{s}{c'}}\cutoff{S}{c'}(\rinf{S}(\gamma)^{-1}y)\prod_\tau|W^{\sigma}_{S,\tau}((\rinf{S}(\gamma)^{-1}y)_\tau)|\le M$
for all~$y$, and $\sum_{\gamma\in\Xi\cap\terms{S}{c}{s}{c'}}\ampabs{\sigma}{S}{c}{c'}{\gamma}<\infty$, where
$\ampabs{\sigma}{S}{c}{c'}{\gamma}$ is the integral of the norm of the integrand of the amplitude.

\Inp{B4} (Note~B) If $(g_j)_j$ is a countable family of measurable functions with values in $\CC^{I}$, $I$ finite,
and $\sum_j\int|g_j|<\infty$, then each coordinate of $\sum_j\int g_j$ is the integral of the corresponding coordinate
of $\sum_jg_j$.}

\subhead{Hermitian spaces with frame, forms and gradings.}\phantomsection\label{s:spaces}
This vocabulary belongs to the background library; it is recalled as the notes of Chapters V and~VII of the write-up
state it. A \emph{hermitian space with frame} $V$ over $E$ consists of a finite index set $I$, signs
$\operatorname{sgn}_V\colon I\to\{\pm1\}$, a matrix $J_V$ over $E$ and a frame $\mathrm{fr}_V\in\GL_I(\AdE)$ such that
$\mathrm{fr}_V^{*}J_V\mathrm{fr}_V$ is the diagonal matrix of the signs at every infinite place and the identity at the
finite places; $\mathrm{fr}_{V,f}$ is its finite part. Its groups are written in the coordinates of the frame:
$G_\infty(V)=\prod_\tau\mathrm{U}(\operatorname{sgn}_V)$, $G_f(V)=\Gf{E}{I}=\{g\in\GL_I(\AEf):g^{*}g=1\}$,
$G(V)=G_\infty(V)\times G_f(V)$, with $\Kf{E}{I}=\Gf{E}{I}\cap\GL_I(\hO)$; each carries the Haar measure fixed in the
library (on $G_f(V)$ the one with $\vol\Kf{E}{I}=1$). The \emph{rational coordinates} of $g\in G(V)$ are
$\mathrm{fr}_Vg\,\mathrm{fr}_V^{-1}$, an adelic point of the unitary group of~$J_V$; the \emph{rational points}
$\Gamma(V)$ are the $g$ whose rational coordinates have entries in~$E$, a discrete subgroup. For a subgroup
$\Delta\subseteq\Gamma(V)$, integrals over $\Delta\backslash G(V)$ are taken for the measure induced by the Haar
measure (in the formal text: a measure on $G(V)$ itself, the Haar measure times a cutoff for~$\Delta$). For a
function $\epsilon$ on $G_\infty(V)$, a \emph{form of type~$\epsilon$} on $V$ is a bounded measurable function
$F\colon G(V)\to\CC$ that is left invariant under $\Gamma(V)$, right invariant under an open subgroup of~$G_f(V)$,
and satisfies $\int_{G_\infty(V)}\epsilon(u)F(g\cdot(u,1))\,\dd u=F(g)$ for all $g$ if $\epsilon$ is integrable; the
group $G_f(V)$ acts on the forms of type $\epsilon$ by right translation, $(R(k)F)(g)=F(gk)$, and the forms fixed by
an open subgroup carry the Petersson product $\langle F,F'\rangle=\int_{\Gamma(V)\backslash G(V)}\overline{F}F'$.
A function $\epsilon$ on $G_\infty(V)$ is an \emph{idempotent kernel} if it is continuous and integrable,
$\epsilon(g^{-1})=\overline{\epsilon(g)}$, $\epsilon(1)\ne0$ and
$\int\epsilon(z)\epsilon((zx)^{-1}y)\,\dd z=\epsilon(x^{-1})\epsilon(y)/\epsilon(1)$. For a family
$l=(l_\tau)_\tau$ of integer vectors indexed by~$I$, the \emph{type} $\epsilon_l$ is the product over $\tau$ of the
functions $\epsilon_{l_\tau}$ on $\mathrm{U}(\operatorname{sgn}_V)$, the formal degree times the complex conjugate of
the canonical matrix coefficient of the discrete series with Harish-Chandra parameter~$l_\tau/2$.

Let $N\in E$ with $N\ne0$, $\bar N=N$, and let $I_0$ be a finite set. A \emph{grading} $\mathcal{Y}$ of $V$ as
$I_0$ plus a string of length $\lambda$ with norm $N$ is a bijection of $I_0\sqcup\{0,\dots,\lambda-1\}$ with $I$
(and a constant that plays no role here) such that $J_V$ is the identity on $I_0$, pairs the string vectors $z_a$,
$z_{a'}$ to $1$ if $a+a'=\lambda-1$, $a\ne a'$, to $-N$ if $a=a'$ and $2a=\lambda-1$, and to $0$ otherwise, and pairs $I_0$
with the string to~$0$; the \emph{grade} is $0$ on $I_0$ and $\lambda-1-2a$ at~$z_a$. With matrices in rational
coordinates, the \emph{unipotent subgroup} $\Uni{\mathcal{Y}}\subseteq G(V)$ consists of the $u$ for which $u-1$
strictly raises grades, the \emph{Bessel group} $\Bes{\mathcal{Y}}$ of the points that fix every string vector of
grade $\ge0$ up to vectors of strictly higher grade, and $\incl{\mathcal{Y}}$ is the inclusion into $G(V)$ of the
adelic group $\prod_\tau\mathrm{U}(I_0)\times\Gf{E}{I_0}$ of the standard definite space on~$I_0$, acting on the span of
$I_0$ and fixing the string. For a shape $\fs$, the \emph{standard bijection} of $\kappa\sqcup\{0,\dots,2\varrho\}$
with $\Ix{\fs}$ sends $j\in\kappa$ to $x_j$ and the string to $e_0,\dots,e_{\varrho-1},z,e^*_{\varrho-1},\dots,e^*_0$
in this order.

\statement{ana}{\textbf{Inputs from the background library and from other units.}\enspace
\Inp{L1} (\emph{Poincar\'e series}; Lib, unit~102) Let $V$ be a hermitian space with frame, $\epsilon$ an idempotent
kernel on $G_\infty(V)$, $K$ an open subgroup of $G_f(V)$, $c\colon G(V)\to\CC$ right invariant under~$K$, and
$(W_q)_q$ a countable family of measurable functions on $G(V)$ with $\sum_{\gamma\in\Gamma(V)}c(\gamma g)=\sum_qW_q(g)$
for all $g$ and $\sum_q|W_q(g)|\le M$ for all $g$, with a finite~$M$. Then there is a form of type $\epsilon$ on
$V$, fixed by~$K$, whose value at $g$ is $\sum_q\int_{G_\infty(V)}\epsilon(u)W_q(g\cdot(u,1))\,\dd u$. Reference to be
confirmed (the file names A.~Borel, \emph{Introduction to automorphic forms}, Proc.\ Sympos.\ Pure Math.~9 (1966), as
the source for Poincar\'e series).

\Inp{L2} (\emph{unfolding}; Lib, unit~114) Let $G$ be a second countable locally compact Hausdorff group with left Haar
measure, $\Gamma$ a discrete subgroup and $\phi\colon G\to\CC$ integrable. Then
$g\mapsto\sum_{\gamma\in\Gamma}\phi(\gamma g)$ is integrable on $\Gamma\backslash G$, with integral $\int_G\phi$.
The catalogue of the write-up names Weil's integration formula (N.~Bourbaki, \emph{Int\'egration}, Chap.~VII), place
not confirmed, of which this is a special case; the second reader derives it from the general formula for proper
actions in D.~Gentner and G.~Last, \emph{Palm pairs and the general mass-transport principle}, §2, (2.5)--(2.9).

\Inp{L3} (\emph{one-sided Poisson summation}; Lib, unit~123) Let $F\colon\AdE\to\RR$ be integrable and non-negative,
and lower semicontinuous at every point of a finite set $\mathcal{V}$ of elements of~$E$ (embedded diagonally). If a
real number $\Sigma$ bounds $\sum_{\xi\in\mathcal{Y}'}\bigl|\int_{\AdE}\psi_E(-\xi x)F(x)\,\dd x\bigr|$ for every finite
set $\mathcal{Y}'\subset E$, where $\psi_E(x)=\exp\bigl(-4\pi i\sum_\tau\operatorname{Re}x_\tau\bigr)\psi_f(x_f)$ is
the standard character, then $\sum_{t\in\mathcal{V}}F(t)\le\Sigma$. By the second reader's verdict this is a weaker
form of the Poisson summation formula of J.~Tate, \emph{Fourier analysis in number fields and Hecke's
zeta-functions}, Theorem~4.1.4 and Lemmas 4.2.1--4.2.4: an inequality, for a non-negative function that need not be
continuous.

\Inp{L4} (Lib, units 57 and~123) (a) For $y\in E$ there is $w\in\AEf^{\times}$ with
$\{t\in\AEf:t\in\hO,\ yt\in\hO\}=w\hO$. (b) (\emph{product formula}) For a non-zero $\sT$-integer $y$ of~$E$:
$\gix{y_{\sT}}^{1/2}\le\prod_\tau|y_\tau|$, where $y_{\sT}$ is the image of $y$ in $\ET{\sT}$. (c) The projection
$\AEf\to\ET{\sT}$ carries the Haar measure restricted to the open subring of the finite adeles that are integral at
every place outside~$\sT$ to $\volO{E}$ times the Haar measure of $\ET{\sT}$; here $\volO{E}=(2/\rdE{E})^{f}$ is
the volume of~$\hO$. (d) For $x\in\ET{\sT}$, extended by $0$ to a finite adele $x^{0}$, and $z\in\AEf$ with image
$z_{\sT}$ in $\ET{\sT}$: $z-(z_{\sT})^{0}\in\hO$ if and only if $z$ is integral at every place outside~$\sT$.
Elementary; no reference is needed.

\Inp{R1} (\emph{the tower}; unit~132, recalled from its Note~T, \emph{The tower method at a place: floor measures,
gauges and the bounds on each floor}) For a place $A$ and an $\fs$-form $T$ with $\JT{T}$ a hermitian unit: a
\emph{hyperbolic basis} $F\in\HypB{\Omega}{T}{\mathfrak{m}}$ is a matrix whose columns form a basis of
$V_\varrho=\langle x_j,z\rangle$ made of $\mathfrak{m}+1$ hyperbolic pairs and of orthonormal vectors indexed
by~$\Omega$, with $k+1=|\Omega|+2\mathfrak{m}+2$ and $\modA(\det T)=1$; $\quadric{F}{b}$ is the measure on the quadric
$\{v\in V_b:h_T(v,v)=h_T(\fr{b},\fr{b})\}$ that is the image of the flat chart measure of the library under the basis
of $V_b$ built from~$F$; the Haar measure of $V_\varrho$ disintegrates along $v\mapsto h_T(v,v)$ into such quadric
measures, with fibre integrals of continuous compactly supported functions that are continuous in the fibre. A
\emph{level} $\cL\in\Levels{\Omega}{T}{\mathfrak{m}}$ consists of a hyperbolic basis $\EL{\cL}$, two $s$-finite
measures $\floorMeasure{\cL}{\varrho}$, $\floorMeasure{\cL}{n}$ on $\GT{T}$, the second non-zero, a unit
$\gL{\cL}=\overline{\gL{\cL}}$, an open set $\Ball{\cL}\ni0$ of vectors with compact closure, a compact set
$\BoxL{\cL}\subseteq\GT{T}$ and a measurable set $\Posn{\cL}\subseteq\GT{T}$, such that
$\floorMeasure{\cL}{\varrho}$ is carried by $\Hl{\varrho}{T}$ with image $\quadric{\cL}{\varrho}$ under
$\sigma\mapsto\sigma^{-1}z$, and $\floorMeasure{\cL}{n}$ is carried by $\HT{T}$, invariant under right translations
by $\HT{T}$, and gives mass at most $1$ to $\{h:u\,h\,g\in\BoxL{\cL}\text{ for some }u\in\UT{T}\}$ for every~$g$. For
$l<\varrho$, $\eta_l$ is the Heisenberg measure, the image of $\dd v\,\dd\beta$ on $V_{l+1}\times A^{-}$ ($A^{-}$ the
skew elements with their measure $\nu_A$) under the parametrisation $(v,\beta)\mapsto\mathrm{n}_l(v,\beta)$ of
$\Nf{l}{T}$, and $\floorMeasure{\cL}{l}=\varpi_l*\eta_l$ with $\varpi_l$ the image of
$\modA(c)^{k+1+2(\varrho-l-1)}\dd c$ under $c\mapsto\wo{l}{c}^{-1}$, the inverse of the Weyl element of the pair
$(e_l,e^*_l)$; $\eta_l=\delta_1$ for $l\ge\varrho$, $\etab{j}=\eta_0*\cdots*\eta_{j-1}$, and
\begin{gather*}\stackMeasure{\cL}_{a,i}=(\eta_{a+i-1}*\floorMeasure{\cL}{a+i-1})*\cdots*(\eta_a*\floorMeasure{\cL}{a}),\\
\groupMeasure{\cL}=\stackMeasure{\cL}_{0,\varrho+2},\qquad
\upperMeasure{\cL}{l}=\floorMeasure{\cL}{n}*\cdots*\floorMeasure{\cL}{l}.\end{gather*}
The measure $\groupMeasure{\cL}$ is right invariant and non-zero. A point $g$ is \emph{near at~$j$},
\m{\IsNear{\cL}{j}{g}}, if $\xi g\in\BoxL{\cL}$ for some $\xi\in\Ql{j}{T}$;
$\floorModulus{\cL}{l}{g}=\modA(g)^{k+1+2(\varrho-(l+1))}\modA(g)^{1/2}$ for a unit~$g$;
$\volumeExponent{\cL}=|\Omega|-(3n+1+\mathfrak{m})$. A \emph{gauge} $\Gamma\in\Gauges{\cL}$ consists of sets
$\Gamma_0,\Gamma_1,\dots\subseteq A$ and a constant $\Lambda(\Gamma)$ subject to the axioms (G1)--(G12) listed in
the proof of Lemma~\ref{l:gauge}.

\Inp{R2} (unit~132, Note~T) Let $\cL\in\Levels{\Omega}{T}{\mathfrak{m}}$. (a) $\groupMeasure{\cL}=\etab{\varrho+2}*\upperMeasure{\cL}{0}$.
(b) Let $\Gamma\in\Gauges{\cL}$, $u$ a unit of $A$, $a<\varrho$, $(v,\beta)\in V_{a+1}\times A^{-}$ and $i\in\NN$. If
for some $\nu\in\Nb{a}{T}$ the column at $e^*_a$ of the matrix $u^{-1}(\nu\,\mathrm{n}_a(v,\beta)-1)$ has all its
coordinates in $\Gamma_i$, then the coordinates of $u^{-1}v$ lie in $\Gamma_i$ and
$u^{-1}(\beta-\tau_Ah_T(v,v))\in\Gamma_i$. (c) Let $\Gamma\in\Gauges{\cL}$ and $M'\in[0,\infty]$, and suppose that for
every $b\le\varrho$ the flat quadric measure on the hyperbolic planes of the chart of $V_b$ (there are
$\mathfrak{m}+\varrho-b$ of them) gives the vectors with all coordinates in $\Gamma_{12}$ mass at most~$M'$. Then
$\upperMeasure{\cL}{l}\{g:\IsNear{\cL}{\varrho+2}{gc}\}\le\bigl(M'\vol(\Gamma_{11})^{|\Omega|}\bigr)^{n-l}$ for all
$l$ and all $c\in\GT{T}$; here $\Ql{\varrho+2}{T}=\UT{T}$. (Three lemmas of Note~T combined: the fibre bound of the
floor quadrics from the cone bound, the masses of the floors, and the induction over the floors.)

\Inp{R3} (\emph{finite places}; Lib, unit~123) A \emph{semilocal field} is a finite product of non-archimedean local
fields, with ring of integers $O_A$ (compact, open, of volume~$1$). An \emph{integral place structure} on a semilocal
field $A$ with involution is a place structure in the sense of unit~132 (module $\modA$, character $\psi_A$, an
element $\tau_A$ with $\tau_A+\bar\tau_A=1$, a skew unit $\delta_A$, the measure $\nu_A$ on the skew elements~$A^{-}$)
with $\tau_A,\delta_A,\delta_A^{-1}\in O_A$, Plancherel constant~$1$, $\nu_A(A^{-}\cap O_A)=1$, $O_A$ self-dual for
$\psi_A$, $\modA(x)=\gix{x}^{-1}$, and with the property that $O_A$ contains an element $c$ with $c\bar c=-1$. Then:
(a) $\vol\{x:ux\in O_A\}=\modA(u)^{-1}$ and, for $\bar u=u$, $\nu_A\{\beta:u\beta\in O_A\}=\modA(u)^{-1/2}$, for
every unit~$u$; $\modA(u)\le1$ for $u\in O_A$; $\modA(\det M)=1$ for $M\in\GL(O_A)$. (b) For a hermitian unit $J$ the
groups $\uLev{g}{J}$ of the integral isometries congruent to $1$ modulo a unit $g$ are open, and $\uLev{1}{J}$ is
compact. (c) (\emph{cone bound}) If $\zetaA{\mu}\le2$ and $\mu+1\le|\Omega'|$, then on $|\Omega'|+1$ hyperbolic planes
the flat quadric measure of the vectors $v$ with $\beta v$ integral is at most
$\zetaA{\mu}\,\modA(\beta)^{-(2|\Omega'|+1)}$, for every unit~$\beta$; here $\zetaA{\mu}=\sum_I[O_A:I]^{-\mu}$ over the
ideals of finite index. (d) (\emph{content}) With $\mathrm{ct}(p)$ the volume of $\{c\in A:cp\text{ integral}\}$ for
a vector $p$ one of whose marked entries is a unit: if $\zetaA{\mu}\le2$, the integral of the $s$-th power of the
content of $s+j$ marked coordinates of $M^{*}w$, $j\ge\mu$, $M\in\GL_{\Omega'}(O_A)$, over the $w$ with $\beta w$
integral is at most $\zetaA{\mu}\,\modA(\beta)^{-(|\Omega'|-s)}$. (e) (\emph{Eichler elements}) Let $J=\JT{T}$, and let
$x$ be a vector with $h_T(x,\cdot)$ the $z$-coordinate and $z$-coordinate~$0$. For a vector $e$ with
$z$-coordinate~$\beta$ put $v=e-z$ and $\varsigma'=h_T(v,e)-\tau_A(h_T(e,e)-h_T(z,z))$. Then
$\mathrm{E}(e)=1+\beta^{-1}\bigl(\beta x\,h_T(v,\cdot)-(v+\varsigma'x)\,h_T(x,\cdot)\bigr)$, with $\beta^{-1}$ read as
$0$ for a non-unit, lies in $\GT{T}$; it lies in $\Hl{\varrho}{T}$ if $e\in V_\varrho$; and
$\mathrm{E}(e)^{-1}z=e$ if $h_T(e,e)=h_T(z,z)$ and $\beta$ is a unit. No reference: these are computations at a finite
place; the catalogue marks (c) as a local density bound that the write-up must read.

\Inp{R4} (unit~133, \emph{DiscreteSeries}) A Bessel system $\mathfrak{B}$ of $\fs$ \emph{has type $(l,\mathsf{k})$}, for a
Harish-Chandra parameter $l$ of the signs and a weight $\mathsf{k}$ of $\mathrm{U}(\kappa)$, if (i) its degree is
Weyl's number $\prod_{l_j<l_i}\frac{l_i-l_j}{2}\big/\prod_{i<r}i!$; (ii) there is an invertible matrix
$\mathsf{A}$ from the coordinates of $\mathfrak{B}$ to the basis of the irreducible representation $\mathbb{S}_{\mathsf{k}}$ of
$\GL_\kappa(\CC)$ of highest weight $\mathsf{k}$ (the Weyl coordinates) with
$\mathsf{A}\,W^{\full}(hy)=\mathbb{S}_{\mathsf{k}}(h^{\mathsf{t}})^{\mathsf{t}}\,\mathsf{A}\,W^{\full}(y)$ for all
$h\in\mathrm{U}(\kappa)$, embedded in $\Gone{\fs}$, and all~$y$; (iii)
$\int\epsilon_l(x)W^{\full}(y\,\mathsf{p}x\mathsf{p}^{-1})\,\dd x=W^{\full}(y)$, the integral over the unitary group
of the signs, with $\mathsf{p}$ the fixed frame.

\Inp{D1} (units 141, 145, 147, 152; the data) For admissible data $D$ for parameters~$P$: (a) (unit~141) the filling
data $\SD{D}$ fit the space $\Ys{D}{0}$ with the parameter $\parW{D}{\tau}{\cdot}$, the weights of the columns, its
grading and a central element, in the sense of Definition~\ref{d:fits}; the shape of $\SD{D}$ has $|\kappa|=\kD{D}$,
and the field is $\ED{D}$ with $\fD{D}$ infinite places. (b) (unit~141, Chapter~V) The \emph{row parameter} is
$\rowPar{D}{\tau}{i}=\kD{D}/2-1-i$ for $i+1<\kD{D}$ and
$\rowPar{D}{\tau}{i}=\hhD{D}-\wt{D}{i+2-\kD{D}}{\tau}$ otherwise, and $\parW{D}{\tau}{j}=2\,\rowPar{D}{\tau}{\nu(j)}$
for a fixed numbering $\nu$ of $\Ridx{D}$ by $\{0,\dots,\Dr{D}-1\}$. (c) (Chapter~III) $\Dm{D}\ge1$, $\Dr{D}\ge1$, $\kD{D}=\Dr{D}+1-\Dm{D}$ is
even, $\hhD{D}=\Dr{D}^{8}$, $\hD{D}=2\hhD{D}$, and $\wt{D}{j}{\tau}=b_\ast+j\,\hD{D}+\hhD{D}$ with a natural number
$b_\ast$ that depends on $\tau$ only; $\nw{D}{\tau}\le\wt{D}{j}{\tau}-\hhD{D}$ for all $j$;
$\log\wt{D}{j}{\tau}\le\log\nw{D}{\tau}+1/\scx{D}$ for $j\le2\Dm{D}$; $\nw{D}{\tau}>0$; and
$(\kD{D}+\aD{D})\,\kD{D}\le\scb{D}^{5/4}$ (a comparison of sizes by the calculus of Chapter~III). (d) (unit~145) The
weight of the columns at $\tau$ is that of the symmetric power of order $a_\tau$ of the standard representation,
with $a_\tau\in\{0,\aD{D}\}$.}

\sect{Poincar\'e series of cut Bessel functions}{s:poincare}
\begin{definition}[filling data that fit a space; the principal level]\label{d:fits}
Let $E$ be a CM field, $\fs$ a shape, and let the index set $\Ix{\fs}$ carry signs with a frame $\mathsf{p}$: a unitary
matrix with $\mathsf{p}^{*}\JT{\Trg{\fs}{1}}\mathsf{p}$ the diagonal matrix of the signs. Let $S\in\Fill{E}{\fs}$, let
$V$ be a hermitian space with frame over $E$ with index set $\Ix{\fs}$ and these signs, let $l=(l_\tau)_\tau$ be a
family of Harish-Chandra parameters for the signs and $\mathsf{k}=(\mathsf{k}_\tau)_\tau$ a family of weights of
$\mathrm{U}(\kappa)$, let $\mathcal{Y}$ be a grading of $V$ as $\kappa$ plus a string of length $2\varrho+1$ with
norm~$\NS{S}$, and let $\eta\in\Gf{E}{\kappa}$. The statement $\Fits{S}{V}{l}{\mathsf{k}}{\mathcal{Y}}{\eta}$, in words
\emph{$S$ fits $V$ with $(l,\mathsf{k},\mathcal{Y},\eta)$}, means:
\begin{enumerate}[label=(F\arabic*),leftmargin=2.6em,itemsep=1pt,topsep=2pt]
\item for every $\tau$ the Bessel system $\BsysS{S}{\tau}$ has type $(l_\tau,\mathsf{k}_\tau)$ in the sense of \Inp{R4};
\item the bijection of $\mathcal{Y}$ is the standard bijection of~$\fs$;
\item the type $\epsilon_l$ is an idempotent kernel on $G_\infty(V)$;
\item there are a matrix $k_0\in\GL_{\Ix{\fs}}(\hO)$ whose reduction at every chosen place of $S$ is~$1$ and a finite
  idele $u$ with $\mathrm{fr}_{V,f}=u\cdot\dot\eta\cdot\latB{S}\cdot k_0$, where $\dot\eta$ acts by $\eta$ on the span of
  the $x_j$ and fixes the other basis vectors;
\item $\eta$ commutes with every element of $\Gf{E}{\kappa}$.
\end{enumerate}
For $S\in\Fill{E}{\fs}$, the \emph{principal level at the chosen places} $\stdK{S}$ is the open subgroup of
$\Gf{E}{\Ix{\fs}}$ of the elements of $\Kf{E}{\Ix{\fs}}$ whose reduction at every chosen place $w\in\ChS{S}{0}\cup\ChS{S}{1}$
is~$1$.
\end{definition}
\leanline{\lean{Fits}\sep\lean{stdPrincipalLevel}}

\begin{setting}[filling data that fit a space]\label{set:fits}
$E$ is a CM field, $\fs$ a shape, the index set $\Ix{\fs}$ carries signs with a frame~$\mathsf{p}$, and $S$, $V$,
$l$, $\mathsf{k}$, $\mathcal{Y}$, $\eta$ are as in Definition~\ref{d:fits}, such that $S$ fits $V$:
$\Fits{S}{V}{l}{\mathsf{k}}{\mathcal{Y}}{\eta}$ holds. The letter $X$ stands for this hypothesis; the objects built from
it carry $X$ as an index. We write $g=(g_\infty,g_f)$ for the points of $G(V)$.
\end{setting}

\begin{definition}[the four dictionaries]\label{d:equivs}
In Setting~\ref{set:fits}:
\begin{enumerate}[label=(\alph*),leftmargin=2.2em,itemsep=1pt,topsep=2pt]
\item $\ratEq{X}\colon\GS{S}{E}\to\Gamma(V)$ is the isomorphism that sends $\gamma$ to the point with rational
  coordinates~$\gamma$, that is $\mathrm{fr}_V^{-1}\gamma\,\mathrm{fr}_V$;
\item $\adEq{X}\colon G(V)\to\GS{S}{\AdE}$ is the isomorphism $g\mapsto\mathrm{fr}_Vg\,\mathrm{fr}_V^{-1}$, the rational
  coordinates;
\item $\reH{X}\colon G_\infty(V)\to\prod_\tau\Gone{\fs}$ is the homeomorphism
  $\reH{X}(g)_\tau=\bsGL{S}{\tau}^{-1}\,\mathrm{fr}_{V,\tau}\,g_\tau\,\mathsf{p}^{-1}$;
\item $\finEq{X}\colon\Gf{E}{\Ix{\fs}}\to\GFin{S}$ is the isomorphism
  $g\mapsto\incH{S}(\eta)^{-1}\cdot\tolb{S}\bigl(\mathrm{fr}_{V,f}\,g\,\mathrm{fr}_{V,f}^{-1}\bigr)\cdot\incH{S}(\eta)$.
\end{enumerate}
They are well defined by Lemma~\ref{l:dict}(a).
\end{definition}
\leanline{\lean{ratEquiv}\sep\lean{adelicEquiv}\sep\lean{realHomeo}\sep\lean{finiteEquiv}}

\begin{lemma}[the dictionary]\label{l:dict}
In Setting~\ref{set:fits}:
\begin{enumerate}[label=(\alph*),leftmargin=2.2em,itemsep=1pt,topsep=2pt]
\item $J_V=\JT{\Trat{S}{E}}$;
\item $\adEq{X}(\ratEq{X}(\gamma))=\radl{S}(\gamma)$ for $\gamma\in\GS{S}{E}$;
\item $\reH{X}\bigl((gg')_\infty\bigr)=\ainf{S}\bigl(\adEq{X}(g)\bigr)\cdot\reH{X}(g'_\infty)$ for $g,g'\in G(V)$;
\item $\finEq{X}(g_f)=\incH{S}(\eta)^{-1}\cdot\afin{S}\bigl(\adEq{X}(g)\bigr)\cdot\incH{S}(\eta)$ for $g\in G(V)$;
\item for $k\in\Gf{E}{\Ix{\fs}}$ and $s\in\GRes{S}$: $\finEq{X}(k)$ lies in $\KS{S}$ and has reduction
  $\redS{S}(\finEq{X}(k))=s$ if and only if $k$ lies in $\Kf{E}{\Ix{\fs}}$ and its reduction at every chosen place
  $w\in\ChS{S}{j}$ is $s_{j,w}$;
\item for $k\in\Gf{E}{\Ix{\fs}}$: $k\in\stdK{S}$ if and only if $\finEq{X}(k)\in\KgS{S}$.
\end{enumerate}
\end{lemma}
\begin{proof}
(a) By (F2) the Gram matrix $J_V$ is the Gram matrix of a grading for the standard bijection, with norm $\NS{S}$: the
identity on the $x_j$, the string vector at position $a$ paired with the one at position $2\varrho-a$ to~$1$, and
the middle one, $z$, with itself to $-\NS{S}$. Positions $a$ and $2\varrho-a$ are $e_a$ and $e^*_a$, so this is
$\JT{\Trg{\fs}{\NS{S}}}$. (b) holds by the definition of $\ratEq{X}$. (c) At $\tau$ the left side is
$\bsGL{S}{\tau}^{-1}\mathrm{fr}_{V,\tau}g_\tau g'_\tau\mathsf{p}^{-1}$, which is the product of
$\bsGL{S}{\tau}^{-1}(\mathrm{fr}_{V,\tau}g_\tau\mathrm{fr}_{V,\tau}^{-1})\bsGL{S}{\tau}$, the component of
$\ainf{S}(\adEq{X}(g))$, and of $\reH{X}(g'_\infty)_\tau$. (d) $\afin{S}$ is $\tolb{S}$ of the finite part. (e) With
$k_0$, $u$ as in (F4), and since $\incH{S}(\eta)=\latB{S}^{-1}\dot\eta\,\latB{S}$, the scalar $u$ and the factors
$\dot\eta$, $\latB{S}$ cancel:
\begin{multline*}\finEq{X}(k)=\latB{S}^{-1}\dot\eta^{-1}\latB{S}\cdot\latB{S}^{-1}\bigl(u\,\dot\eta\,\latB{S}k_0\,k\,k_0^{-1}\latB{S}^{-1}\dot\eta^{-1}u^{-1}\bigr)\latB{S}\\
\cdot\latB{S}^{-1}\dot\eta\,\latB{S}=k_0\,k\,k_0^{-1}.\end{multline*}
Since $k_0\in\GL(\hO)$, the matrix $k_0kk_0^{-1}$ is integral with integral inverse exactly when $k$ is, and since
$k_0$ reduces to $1$ at the chosen places, the two have the same reductions there. (f) is (e) for $s=1$, by the
definitions of the two levels.
\end{proof}
\leanline{\lean{JeqratGram}\sep\lean{adelicEquivratEquiv}\sep\lean{realHomeomul}\sep\lean{finiteEquiveqadelicToFin}\sep
\lean{reductionfiniteEquiviff}\sep\lean{memstdPrincipalLeveliff}}

\begin{definition}[cut Bessel functions and Poincar\'e series]\label{d:poin}
In Setting~\ref{set:fits}, for a pair $a=(c,i)\in\Cl{S}\times\Bc{S}$ and $g\in G(V)$, with $y=\reH{X}(g_\infty)$,
\[\cutB{X}{a}{g}=\begin{cases}\cutoff{S}{c}(y)\cdot\besfn{\full}{S}{y}{i}&\text{if }
\finEq{X}^{-1}\bigl(\hc{S}{c}^{-1}\incH{S}(\eta)\bigr)\cdot g_f\in\stdK{S},\\[2pt] 0&\text{otherwise,}\end{cases}\]
where the value of the cutoff is read as $0$ where it is infinite. For a coefficient vector
$\phi\in\CC^{\Cl{S}\times\Bc{S}}$, a subgroup $\Delta$ of $G(V)$ and $g\in G(V)$, the \emph{Poincar\'e series} is
\[\Poin{X}{\phi}{\Delta}{g}=\sum_{a}\phi_a\sum_{\gamma\in\Delta}\cutB{X}{a}{\gamma g},\]
the sum over the pairs $a$; an inner sum that does not converge unconditionally is read as~$0$.
\end{definition}
\leanline{\lean{cutBessel}\sep\lean{poincarea}}

\begin{lemma}\label{l:cbrat}
In Setting~\ref{set:fits}, let $a=(c,i)$ be a pair in $\Cl{S}\times\Bc{S}$, $g\in G(V)$ and $\gamma\in\GS{S}{E}$. Put
$y=\rinf{S}(\gamma)\cdot\reH{X}(g_\infty)$. Then $\cutB{X}{a}{\ratEq{X}(\gamma)\,g}=\cutoff{S}{c}(y)\,\besfn{\full}{S}{y}{i}$
if $\hc{S}{c}^{-1}\,\rfin{S}(\gamma)\,\incH{S}(\eta)\,\finEq{X}(g_f)\in\KgS{S}$, and
$\cutB{X}{a}{\ratEq{X}(\gamma)\,g}=0$ otherwise.
\end{lemma}
\begin{proof}
By Lemma~\ref{l:dict}(f) and since $\finEq{X}$ is a homomorphism, the condition in the definition of
$\cutB{X}{a}{g}$ reads $\hc{S}{c}^{-1}\incH{S}(\eta)\,\finEq{X}(g_f)\in\KgS{S}$. Replace $g$ by $\ratEq{X}(\gamma)g$: by
(d) and (b) of the lemma, $\finEq{X}\bigl((\ratEq{X}(\gamma)g)_f\bigr)=\incH{S}(\eta)^{-1}\rfin{S}(\gamma)\incH{S}(\eta)\,\finEq{X}(g_f)$,
and by (c) and (b), $\reH{X}\bigl((\ratEq{X}(\gamma)g)_\infty\bigr)=\rinf{S}(\gamma)\,\reH{X}(g_\infty)$.
\end{proof}
\leanline{\lean{cutBesselratEquivmul}}

\begin{lemma}\label{l:classunip}
In Setting~\ref{set:fits}, let $E$ be stably unramified and $S$ admissible, let $a=(c,i)$ be a pair in
$\Cl{S}\times\Bc{S}$ and $g\in G(V)$. Then
\[\sum_{u\in\clU{S}{c}}\cutB{X}{a}{\ratEq{X}(u)\,g}=\begin{cases}\besfn{\full}{S}{\reH{X}(g_\infty)}{i}&\text{if }
\hc{S}{c}^{-1}\,\incH{S}(\eta)\,\finEq{X}(g_f)\in\KgS{S},\\[2pt]0&\text{otherwise.}\end{cases}\]
\end{lemma}
\begin{proof}
Let $y=\reH{X}(g_\infty)$ and $u\in\clU{S}{c}$. By the definition of $\clU{S}{c}$,
$\hc{S}{c}^{-1}\rfin{S}(u)\hc{S}{c}\in\KgS{S}$, so the condition of Lemma~\ref{l:cbrat} for $u$,
\[\bigl(\hc{S}{c}^{-1}\rfin{S}(u)\hc{S}{c}\bigr)\cdot\hc{S}{c}^{-1}\incH{S}(\eta)\finEq{X}(g_f)\in\KgS{S},\]
does not depend on~$u$. When it holds, the term of $u$ is
$\cutoff{S}{c}(\rinf{S}(u)y)\,\besfn{\full}{S}{\rinf{S}(u)y}{i}=\cutoff{S}{c}(\rinf{S}(u)y)\,\besfn{\full}{S}{y}{i}$ by
\Inp{B2}, and the cutoffs sum to $1$ over $u$ by \Inp{B1}.
\end{proof}
\leanline{\lean{tsumcutBesselclassUnip}}

\begin{proof}[Proof of Theorem A(a)]
Let $M$ be the constant of \Inp{B3} for $\Xi=\GS{S}{E}$, the class $c$ in both places, $s=1$ and the full side. Fix
$g$ and put $y_0=\reH{X}(g_\infty)$. If no $\delta\in\GS{S}{E}$ satisfies the condition
$\hc{S}{c}^{-1}\rfin{S}(\delta)\incH{S}(\eta)\finEq{X}(g_f)\in\KgS{S}$ of Lemma~\ref{l:cbrat}, all terms
$\cutB{X}{a}{\gamma g}$ vanish. Otherwise fix such a $\delta$, call $\mathsf{m}_0\in\KgS{S}$ that element, and put
$y=\rinf{S}(\delta)y_0$. Every element of $\Gamma(V)$ is $\ratEq{X}(\gamma^{-1}\delta)$ for a unique
$\gamma\in\GS{S}{E}$, and by Lemma~\ref{l:cbrat} its term is
$\cutoff{S}{c}(\rinf{S}(\gamma)^{-1}y)\,\besfn{\full}{S}{\rinf{S}(\gamma)^{-1}y}{i}$ if
$\hc{S}{c}^{-1}\rfin{S}(\gamma)^{-1}\hc{S}{c}\cdot\mathsf{m}_0\in\KgS{S}$, and $0$ otherwise. The condition says
that $\pos{S}{c}{c}{\gamma}\in\KgS{S}$, that is, $\gamma\in\terms{S}{c}{1}{c}$. Since a coordinate of a vector is
bounded by its norm, $|\besfn{\full}{S}{y'}{i}|\le\prod_\tau|W^{\full}_{S,\tau}(y'_\tau)|$. So the sum is at most the sum
bounded by $M$ in \Inp{B3}.
\end{proof}

\begin{proposition}[unfolding]\label{p:unfold}
In Setting~\ref{set:fits}, let $E$ be stably unramified and $S$ admissible. Let $\phi\in\CC^{\Cl{S}\times\Bc{S}}$, let
$\Delta\subseteq\Gamma(V)$ be a subgroup, $q$ a function on $G_\infty(V)$ and $z$ a form of type $q$ on~$V$. Then
the function $g\mapsto z(g)\,\Poin{X}{\phi}{\Delta}{g}$ is integrable on $\Delta\backslash G(V)$, and
\[\int_{\Delta\backslash G(V)}z(g)\,\Poin{X}{\phi}{\Delta}{g}\,\dd g=\sum_a\phi_a\int_{G(V)}z(g)\,\cutB{X}{a}{g}\,\dd g.\]
\end{proposition}
\begin{proof}
Fix a pair $a$. The function $\cutB{X}{a}{\,\cdot\,}$ is measurable: it is the product of measurable functions of
$g_\infty$ with the indicator of the preimage of the open and closed set $\stdK{S}$. It is integrable on $G(V)$:
by \Inp{L2} for non-negative functions and Theorem~A(a),
$\int_{G(V)}|\cutB{X}{a}{g}|\,\dd g=\int_{\Gamma(V)\backslash G(V)}\sum_\gamma|\cutB{X}{a}{\gamma g}|\,\dd g\le M\cdot\vol(\Gamma(V)\backslash G(V))$,
which is finite because the quotient has finite volume (reduction theory, proved in the background library;
reference to be confirmed). As $z$ is bounded and measurable, $z\cdot\cutB{X}{a}{\,\cdot\,}$ is integrable, and \Inp{L2}
for the discrete subgroup $\Delta$ shows that $g\mapsto\sum_{\delta\in\Delta}z(\delta g)\cutB{X}{a}{\delta g}$ is
integrable on $\Delta\backslash G(V)$ with integral $\int_{G(V)}z\cdot\cutB{X}{a}{\,\cdot\,}$. Since $z(\delta g)=z(g)$ for
$\delta\in\Gamma(V)$, this function is $z(g)\sum_\delta\cutB{X}{a}{\delta g}$. Multiply by $\phi_a$ and sum over the
finitely many~$a$.
\end{proof}
\leanline{\lean{integralpoincare}\sep\lean{measurablecutBessel}}

\begin{proof}[Proof of Theorem A(b)]
Fix $a=(c,i)$. Every $\gamma'\in\Gamma(V)$ is $\ratEq{X}(u\,\gamma_q^{-1})$ for a unique coset
$q\in\GS{S}{E}/\clU{S}{c}$, with its chosen representative $\gamma_q$, and a unique $u\in\clU{S}{c}$. By
Theorem~A(a) the sum over $\Gamma(V)$ converges absolutely, so it may be grouped:
\begin{gather*}\sum_{\gamma'\in\Gamma(V)}\cutB{X}{a}{\gamma'g}=\sum_qW_q(g),\qquad\sum_q|W_q(g)|\le M,\\
W_q(g)=\sum_{u\in\clU{S}{c}}\cutB{X}{a}{\ratEq{X}(u)\,\ratEq{X}(\gamma_q)^{-1}g}.\end{gather*}
The function $\cutB{X}{a}{\,\cdot\,}$ is right invariant under $\stdK{S}$, because the condition on $g_f$ is. By (F3)
and \Inp{L1} there is a form $y_a$ of type $\epsilon_l$ on $V$, fixed by $\stdK{S}$, with
$y_a(g)=\sum_q\int\epsilon_l(x)W_q(g\cdot(x,1))\,\dd x$. We claim that each integral is $W_q(g)$. By
Lemma~\ref{l:classunip}, $W_q(g)$ is $\besfn{\full}{S}{\reH{X}(g'_\infty)}{i}$ or~$0$, according to a condition on
$g'_f$, where $g'=\ratEq{X}(\gamma_q)^{-1}g$. Replacing $g$ by $g\cdot(x,1)$ leaves $g'_f$ unchanged and multiplies
$\reH{X}(g'_\infty)_\tau$ on the right by $\mathsf{p}x_\tau\mathsf{p}^{-1}$. As $\epsilon_l$ and the Bessel function
are products over~$\tau$, the integral is the product over $\tau$ of
$\int\epsilon_{l_\tau}(x)\,W^{\full}_{S,\tau}(y_\tau\mathsf{p}x\mathsf{p}^{-1})_{i_\tau}\,\dd x=W^{\full}_{S,\tau}(y_\tau)_{i_\tau}$,
by (F1) and \Inp{R4}(iii). So $y_a(g)=\sum_{\gamma'}\cutB{X}{a}{\gamma'g}$. Let $\Pmap{X}$ be the linear map with
$\Pmap{X}(e_a)=y_a$ for all~$a$; then $\Pmap{X}(\phi)(g)=\sum_a\phi_ay_a(g)=\Poin{X}{\phi}{\Gamma(V)}{g}$. (In the
formal text $\Pmap{X}$ is a definition: the forms $y_a$ are chosen once, by the axiom of choice.)
\end{proof}

\begin{proof}[Proof of Theorem B]
The function $z=\overline{R(k)\Pmap{X}(e_{a'})}$ is a form on $V$ (of the conjugate type): bounded, measurable,
invariant under $\Gamma(V)$. By the definition of the Petersson product, Theorem~A(b) and
Proposition~\ref{p:unfold} for $\phi=e_a$ and $\Delta=\Gamma(V)$,
\[\bigl\langle R(k)\Pmap{X}(e_{a'}),\Pmap{X}(e_a)\bigr\rangle=\int_{G(V)}z(g)\,\cutB{X}{a}{g}\,\dd g,\qquad
z(g)=\sum_{\gamma\in\GS{S}{E}}\overline{\cutB{X}{a'}{\ratEq{X}(\gamma)^{-1}gk}},\]
where the elements of $\Gamma(V)$ were written as $\ratEq{X}(\gamma)^{-1}$. Put $y=\reH{X}(g_\infty)$ and
$\mathsf{m}=\hc{S}{c}^{-1}\incH{S}(\eta)\finEq{X}(g_f)$, so that $\cutB{X}{a}{g}=\cutoff{S}{c}(y)\besfn{\full}{S}{y}{i}$ if
$\mathsf{m}\in\KgS{S}$ and $0$ otherwise. By Lemma~\ref{l:cbrat}, the term of $\gamma$ in $z(g)$ is the conjugate of
$\cutoff{S}{c'}(\rinf{S}(\gamma)^{-1}y)\,\besfn{\full}{S}{\rinf{S}(\gamma)^{-1}y}{i'}$ if
\[\hc{S}{c'}^{-1}\rfin{S}(\gamma)^{-1}\incH{S}(\eta)\finEq{X}(g_f)\finEq{X}(k)=\pos{S}{c}{c'}{\gamma}^{-1}\cdot\mathsf{m}\cdot\finEq{X}(k)\in\KgS{S}\]
and $0$ otherwise. By Lemma~\ref{l:dict}(e) and the hypothesis on $k$, $\finEq{X}(k)\in\KS{S}$ has reduction~$s$. So
for $\mathsf{m}\in\KgS{S}$ the displayed condition says that $\pos{S}{c}{c'}{\gamma}$ lies in $\KS{S}$ and has
reduction $s$, that is $\gamma\in\terms{S}{c}{s}{c'}$: if the product lies in $\KgS{S}$, then
$\pos{S}{c}{c'}{\gamma}\in\mathsf{m}\,\finEq{X}(k)\,\KgS{S}$; conversely, if $\pos{S}{c}{c'}{\gamma}\in\KS{S}$ has
reduction $s$, the product lies in $\KS{S}$ and reduces to $s^{-1}\cdot1\cdot s=1$. Hence the integrand is
\[\mathbf{1}\bigl[\mathsf{m}\in\KgS{S}\bigr]\sum_{\gamma\in\terms{S}{c}{s}{c'}}\cutoff{S}{c}(y)\,\cutoff{S}{c'}(\rinf{S}(\gamma)^{-1}y)\,
\besker{\full}{S}{\gamma}{y}_{(i,i')}.\]
The Haar measure of $G(V)$ is the product of those of $G_\infty(V)$ and $G_f(V)$. The condition
$\mathsf{m}\in\KgS{S}$ says, by Lemma~\ref{l:dict}(f), that $g_f$ lies in a left translate of $\stdK{S}$, of volume
$\vol(\stdK{S})$. The map $\reH{X}$ carries the Haar measure of $G_\infty(V)$ to that of $\prod_\tau\Gone{\fs}$,
being at each place a conjugation by the frame, which preserves the Haar measures fixed in the library, followed by
a left translation. The sum over $\gamma$ and the integral over $y$ may be exchanged, and the coordinate $(i,i')$
taken inside, by \Inp{B4}, since $\sum_\gamma\ampabs{\full}{S}{c}{c'}{\gamma}<\infty$ by \Inp{B3}. This gives
$\vol(\stdK{S})\sum_{\gamma\in\terms{S}{c}{s}{c'}}\amp{\full}{S}{c}{c'}{\gamma}_{(i,i')}$, which is the claim.
\end{proof}

\sect{Columns in Bessel coordinates}{s:columns}
The lemma of this section identifies the groups of a grading with the Bessel subgroups of Note~S, and the two
definitions describe the algebraic modular forms on the definite part (the \emph{columns} of the write-up) in the
coordinates in which the Poincar\'e series are written. They are used by the note on the seesaw identity (unit~141).

\begin{lemma}\label{l:inclx}
In Setting~\ref{set:fits}, where the hypothesis is $\Fits{S}{V}{l}{\mathsf{k}}{\mathcal{Y}}{\eta}$, let $g\in G(V)$,
and let $h=(h_\infty,h_f)$ be a point $h$ of the adelic group $\prod_\tau\mathrm{U}(\kappa)\times\Gf{E}{\kappa}$ of the
standard definite space on~$\kappa$.
\begin{enumerate}[label=(\alph*),leftmargin=2.2em,itemsep=1pt,topsep=2pt]
\item $g\in\Uni{\mathcal{Y}}$ if and only if $\adEq{X}(g)\in\UT{\Trat{S}{\AdE}}$, and $g\in\Bes{\mathcal{Y}}$ if and only
  if $\adEq{X}(g)\in\ST{\Trat{S}{\AdE}}$.
\item $\afin{S}\bigl(\adEq{X}(\incl{\mathcal{Y}}(h))\bigr)=\incH{S}(h_f)$.
\item $\finEq{X}\bigl(\incl{\mathcal{Y}}(h)_f\bigr)=\incH{S}(h_f)$.
\item $\ainf{S}\bigl(\adEq{X}(\incl{\mathcal{Y}}(h))\bigr)$ is the family of the elements of $\Gone{\fs}$ that act by
  $h_\tau$ on the span of the $x_j$ and fix the other basis vectors.
\end{enumerate}
\end{lemma}
\begin{proof}
(a) By (F2) the basis vectors $e_a$, the $x_j$ and $z$, and $e^*_a$ have the grades $2(\varrho-a)$, $0$ and
$-2(\varrho-a)$. So ``$u-1$ strictly raises grades'' says that the matrix of $u$ is block upper triangular with
identity diagonal blocks for the order $e_0,\dots,e_{\varrho-1}$, $\langle x_j,z\rangle$, $e^*_{\varrho-1},\dots,e^*_0$,
which defines $\UT{\cdot}$; and ``fixes $e_a$ and $z$ up to vectors of higher grade'' says that
$u\fr{a}-\fr{a}$ lies in the span of the earlier frame vectors, which defines $\ST{\cdot}$. (b), (d) The rational
coordinates of $\incl{\mathcal{Y}}(h)$ are the block matrix that is $h$ on the $x_j$ and $1$ on the string. Its
finite part in the lattice basis is $\incH{S}(h_f)$ by the definition of $\incH{S}$; at $\tau$, the scales of
$\bsGL{S}{\tau}$ are $1$ on the $x_j$, so the conjugation of $\ainf{S}$ changes nothing. (c) By
Lemma~\ref{l:dict}(d) and (b), $\finEq{X}(\incl{\mathcal{Y}}(h)_f)=\incH{S}(\eta)^{-1}\incH{S}(h_f)\incH{S}(\eta)$, and
$\eta$ commutes with $h_f$ by (F5).
\end{proof}
\leanline{\lean{memunipotentiff}\sep\lean{adelicToFinadelicEquivinclX}\sep\lean{finiteEquivinclX}\sep
\lean{adelicToInfinclX}}

\begin{definition}[column coordinates]\label{d:colc}
Let $E$ be a CM field, $\fs$ a shape, $S\in\Fill{E}{\fs}$, and let $V$, $l$, $\mathsf{k}$, $\mathcal{Y}$, $\eta$ be as
in Definition~\ref{d:fits} with $\Fits{S}{V}{l}{\mathsf{k}}{\mathcal{Y}}{\eta}$; $h$ stands for this hypothesis.
\begin{enumerate}[label=(\alph*),leftmargin=2.2em,itemsep=1pt,topsep=2pt]
\item Let $\mathbb{S}_{\mathsf{k}}=\bigotimes_\tau\mathbb{S}_{\mathsf{k}_\tau}$, with its basis, the \emph{Weyl coordinates}, and
  for each $\tau$ let $\mathsf{A}^{(\tau)}$ be a matrix as in \Inp{R4}(ii), chosen once. Then
  $\wbe{h}\colon\mathbb{S}_{\mathsf{k}}\to\CC^{\Bc{S}}$ is the linear isomorphism whose matrix is the transpose of the
  Kronecker product of the~$\mathsf{A}^{(\tau)}$. It satisfies, for all families $u=(u_\tau)$ in $\mathrm{U}(\kappa)$, all
  $y\in\prod_\tau\Gone{\fs}$ and all $v\in\mathbb{S}_{\mathsf{k}}$,
  \[\wbe{h}(v)\cdot\besfn{\full}{S}{\dot u\,y}{\,\cdot\,}=\wbe{h}\bigl(\mathbb{S}_{\mathsf{k}}(u^{\mathsf{t}})v\bigr)\cdot\besfn{\full}{S}{y}{\,\cdot\,},\]
  where the dot is the bilinear pairing $\sum_ia_ib_i$ on $\CC^{\Bc{S}}$, $\dot u$ acts by $u_\tau$ on the span of
  the $x_j$ at $\tau$, and $\mathbb{S}_{\mathsf{k}}(u^{\mathsf{t}})=\bigotimes_\tau\mathbb{S}_{\mathsf{k}_\tau}(u_\tau^{\mathsf{t}})$.
\item $\colc{h}$ is the linear isomorphism from the space of the algebraic modular forms of weight
  $(\mathsf{k},0)$ on the standard definite space on $\kappa$ that are fixed by $\KcolS{S}$ to
  $\CC^{\Cl{S}\times\Bc{S}}$ that sends a form $F$ to the family $\bigl(\wbe{h}(F(g_c))\bigr)_{c\in\Cl{S}}$, where $g_c$ is
  the chosen representative of the class~$c$. It is an isomorphism because the level is free: the rational points
  meet every conjugate of $\KcolS{S}$ trivially (a condition on filling data).
\end{enumerate}
\end{definition}
The displayed property follows from \Inp{R4}(ii) at each place: with $\mathsf{A}=\bigotimes_\tau\mathsf{A}^{(\tau)}$,
$\wbe{h}(v)\cdot W(\dot uy)=v^{\mathsf{t}}\mathsf{A}W(\dot uy)=v^{\mathsf{t}}\,\mathbb{S}_{\mathsf{k}}(u^{\mathsf{t}})^{\mathsf{t}}\mathsf{A}W(y)$.
\leanline{\lean{weylBesselEquiv}\sep\lean{columnCoords}}

\sect{The adelic sieve}{s:sieve}
Throughout this section $\sT$ is a finite set of finite places of~$E$, $\ET{\sT}=\prod_{w\in\sT}E_w$ with ring of
integers $O_{\sT}$, and $y_{\sT}\in\ET{\sT}$ is the image of $y\in E$ or of $y\in\AEf$. An element of $E$ is a
\emph{$\sT$-integer} if it is integral at every finite place outside~$\sT$, and a \emph{$\sT$-unit} if it and its
inverse are $\sT$-integers. For $x\in\ET{\sT}$, $x^{0}\in\AEf$ is its extension by zero.

\begin{definition}[the character, the gain and the points at $\sT$]\label{d:sievedefs}
Let $E$ be a CM field, $\fs$ a shape, $S\in\Fill{E}{\fs}$ and $\sT$ a finite set of finite places of~$E$.
\begin{enumerate}[label=(\alph*),leftmargin=2.2em,itemsep=1pt,topsep=2pt]
\item $\charT{S}{\sT}$ is the additive character $x\mapsto\sdc{S}(x^{0})$ of $\ET{\sT}$.
\item For $\xi\in\ET{\sT}$, the \emph{frequency gain} is
  $\fgain{S}{\sT}{\xi}=\bigl(\gix{\xi}^{1/2}\bigr)^{\kk{\fs}-\nn{\fs}^{7}}\big/\den{\xi}^{6}\cdot\bS{S}^{\nn{\fs}^{9}\ninf{E}}$.
\item For a family $l=(l_\tau)_\tau$ of complex numbers and $t\in E$,
  $\ratpt{\sT}{l}{t}=\bigl((l_\tau t_\tau)_\tau,\ t_{\sT}\bigr)\in\CC^{f}\times\ET{\sT}$.
\end{enumerate}
\end{definition}
\leanline{\lean{charOnPlaces}\sep\lean{frequencyGain}\sep\lean{ratPoint}}

\begin{lemma}\label{l:charT}
Let $S\in\Fill{E}{\fs}$ and let $\sT$ be a finite set of finite places of~$E$. If $x\in\AEf$ is a finite adele $x$
that is integral at every finite place outside~$\sT$, then $\charT{S}{\sT}(x_{\sT})=\sdc{S}(x)$.
\end{lemma}
\begin{proof}
By \Inp{L4}(d), $x-(x_{\sT})^{0}\in\hO$, and $\sdc{S}$ is trivial on $\hO$ by \Inp{S2} for $a=1$.
\end{proof}
\leanline{\lean{charOnPlacesevalPlaces}}

\begin{proposition}[the sum over frequencies]\label{p:freq}
Let $E$ carry Weil data, let $S\in\Fill{E}{\fs}$, and let $\sT$ be a finite set of finite places of~$E$ that is
stable under complex conjugation. Let $B$ be a real number with $B\ge1$, $j$ a natural number~$j$, $\mathcal{X}$ a
finite set of non-zero $\sT$-integers of~$E$, and let $c=(c_\tau)$ be complex and $x=(x_\tau)$ real numbers, indexed
by the infinite places, with $|c_\tau|^{-1}\le|x_\tau|\le B$ for all~$\tau$. Then, with $b=\bS{S}$,
\[\sum_{y\in\mathcal{X}}\ \prod_\tau\damp{j}{|c_\tau y_\tau|^{-1}}\cdot\frac{\bigl(\gix{y_{\sT}}^{1/2}\bigr)^{j}}{\den{y_{\sT}}^{6}}
\;\le\;\prod_\tau\Bigl(|x_\tau|^{j}\,(bB^{2})^{3}\Bigr)\cdot\bigl(b\,\rdE{E}\bigr)^{3f}.\]
\end{proposition}
\begin{proof}
Weil data include that $E$ is stably unramified, so \Inp{S6} applies. Fix $y\in\mathcal{X}$ and let $L_y\subset E^{2}$ be
the line spanned by $(1,y)$; distinct $y$ give distinct lines. Let $w$ be as in \Inp{L4}(a). Then $w\in\hO$, and
$w\hO$ is the preimage in $\hO$ of $\{t\in O_{\sT}:y_{\sT}t\in O_{\sT}\}$, because $y$ is integral outside~$\sT$;
so $\gix{w}=[\hO:w\hO]=\den{y_{\sT}}$. By \Inp{S6},
\begin{equation}\label{e:line}
\sld{1}{L_y}\le\prod_\tau(1+|y_\tau|^{2})\cdot\den{y_{\sT}} .
\end{equation}
Now bound the term of~$y$. First, $\damp{j}{\rho}=\rho^{j}\min(1,\rho)^{6}$ and $|c_\tau y_\tau|^{-1}\le|x_\tau|/|y_\tau|$,
so the product of the powers $\rho^{j}$ is at most $\prod_\tau|x_\tau|^{j}\cdot\prod_\tau|y_\tau|^{-j}$, and
$\bigl(\gix{y_{\sT}}^{1/2}\bigr)^{j}\prod_\tau|y_\tau|^{-j}\le1$ by the product formula \Inp{L4}(b). Second, for each
$\tau$,
\[\min\bigl(1,|c_\tau y_\tau|^{-1}\bigr)^{6}\le\Bigl(\frac{1000\,B^{2}}{1+|y_\tau|^{2}}\Bigr)^{3}:\]
if $|y_\tau|\le1$ the right side is at least $(500B^{2})^{3}\ge1$; if $|y_\tau|\ge1$ the left side is at most
$(B/|y_\tau|)^{6}\le\bigl(2B^{2}/(1+|y_\tau|^{2})\bigr)^{3}$. Third, $\den{y_{\sT}}\ge1$, so
$\den{y_{\sT}}^{-6}\le\den{y_{\sT}}^{-3}$. Together with \eqref{e:line} the term of $y$ is at most
$\prod_\tau|x_\tau|^{j}(1000B^{2})^{3}\cdot\sld{1}{L_y}^{-3}$. Summing over the distinct lines and using \Inp{S6},
\[\sum_{y\in\mathcal{X}}\cdots\le\prod_\tau|x_\tau|^{j}(1000B^{2})^{3}\cdot(400f\,\rdE{E})^{3f}\le\prod_\tau|x_\tau|^{j}(bB^{2})^{3}\cdot(b\,\rdE{E})^{3f},\]
since $1000\le b$ and $400f\le b$ by \Inp{S1}.
\end{proof}
\leanline{\lean{sumfrequenciesle}}

\begin{proof}[Proof of Theorem C]
Write $\mathrm{pt}(l;t)$ for $\ratpt{\sT}{l}{t}$, $\mathrm{fs}(\tau,a)$ for $\fscale{S}{\tau}{a}$, $g$ for $\gS{S}$, and
$\Lambda^{(i)}=\sum_{p\in\mathcal{A}^{i}}\Lambda_i(\mathrm{hd}_i(p,1))$. We prove the claim for $i$ by induction,
assuming it for all smaller indices and all subsets of~$\mathcal{A}$. Fix $i\le\mathsf{n}$ and
$\mathcal{V}\subseteq\mathcal{A}$.

\emph{The step function.} For $x\in C$ put
\[F(x)=\sum_{p,q\in\mathcal{A}^{i}}J_{i+1}\bigl(\mathrm{hd}_i(p,1)[i\mapsto u(x)],\ \mathrm{hd}_i(q,1)[i\mapsto u(x)]\bigr).\]
By \Inp{S4}(a), $F$ is integrable, $\int F=\sum_{p,q}J_i(\mathrm{hd}_i(p,1),\mathrm{hd}_i(q,1))$, and
$|\int\psi(\xi x)F(x)\,\dd x|\le\beta(\xi)(\Lambda^{(i)})^{2}$ for every unit $\xi$ of~$C$. By \Inp{S4}(b), applied to
the sequences $\mathrm{hd}_i(p,1)[i\mapsto u(x)]$, which share their entry~$i$, $F(x)$ is a non-negative real number,
and $F$ is lower semicontinuous at every $x$ all of whose coordinates are units, since $u$ is continuous there.

\emph{Transport to the adeles.} Let $U_0=\CC^{f}\times\{z\in\AEf:z\text{ integral outside }\sT\}$, an open subset of
$\AdE$, and $\mathrm{m}\colon\AdE\to C$, $\mathrm{m}(z)=\bigl((\mathrm{fs}(\tau,i)z_\tau)_\tau,\ (z_f)_{\sT}\bigr)$. The
map $z_\tau\mapsto\mathrm{fs}(\tau,i)z_\tau$ of $\CC$ divides Lebesgue measure by $\mathrm{fs}(\tau,i)^{2}$, and by
\Inp{L4}(c) the projection to $\ET{\sT}$ multiplies by $\volO{E}$; so $\mathrm{m}$ carries the Haar measure
restricted to $U_0$ to $c_i$ times the Haar measure of~$C$. Put $G=\mathbf{1}_{U_0}\cdot(F\circ\mathrm{m})$. Then
$G\ge0$, $G$ is integrable with $\int G=c_i\int F$, and for $t\in\mathcal{V}$, a $\sT$-unit: $t\in U_0$, all
coordinates of $\mathrm{m}(t)=\mathrm{pt}((\mathrm{fs}(\tau,i))_\tau;t)$ are units, so $G$ is lower semicontinuous
at~$t$, and
\[G(t)=F(\mathrm{m}(t))=\sum_{p,q\in\mathcal{A}^{i}}J_{i+1}\bigl(\mathrm{hd}_i(p,t),\mathrm{hd}_i(q,t)\bigr),\]
because $\mathrm{hd}_i(p,1)[i\mapsto u(\mathrm{m}(t))]=\mathrm{hd}_i(p,t)$. So the left side of the claim is
$\sum_{t\in\mathcal{V}}G(t)$.

\emph{The Fourier coefficients.} Let $y\in E$ and consider the coefficient at the frequency $y/\vpS{S}$,
\[\hat G(y)=\int_{\AdE}\psi_E(-(y/\vpS{S})z)\,G(z)\,\dd z;\]
every element of $E$ is such a frequency. (i) If $y$ is not a $\sT$-integer, $\hat G(y)=0$. Indeed, for $r\in\hO$ the adele
$r'=r-(r_{\sT})^{0}$ is integral outside $\sT$ and vanishes at~$\sT$, so $G(z+(0,r'))=G(z)$, while
$\psi_E(-(y/\vpS{S})(0,r'))=\sdc{S}(yr')$; by translation invariance of the integral, $\hat G(y)=\sdc{S}(yr')\hat G(y)$.
If $\hat G(y)\ne0$, then $\sdc{S}\bigl(r(y-(y_{\sT})^{0})\bigr)=1$ for all $r\in\hO$, so $y-(y_{\sT})^{0}\in\hO$ by
\Inp{S2}, and $y$ is a $\sT$-integer by \Inp{L4}(d). (ii) Let $y\ne0$ be a $\sT$-integer. Since
$\overline{\vpS{S}}=-\vpS{S}$, the number $\vpS{S}_\tau$ is purely imaginary, and a direct computation gives, for
$z\in U_0$ and $x=\mathrm{m}(z)$,
\begin{gather*}\psi_E\bigl(-(y/\vpS{S})z\bigr)=\prod_\tau\exp(\xi_\tau x_\tau-\overline{\xi_\tau x_\tau})\cdot\sdc{S}(yz_f)=\psi(\xi x),\\
\xi=\mathrm{pt}\bigl((c_\tau/\mathrm{fs}(\tau,i))_\tau;\,y\bigr),\qquad c_\tau=\frac{2\pi i}{\vpS{S}_\tau},\end{gather*}
where $\sdc{S}(yz_f)=\charT{S}{\sT}(y_{\sT}x_{\sT})$ by Lemma~\ref{l:charT}, as $yz_f$ is integral outside~$\sT$.
Hence $\hat G(y)=c_i\int_C\psi(\xi x)F(x)\,\dd x$, and $\xi$ is a unit of $C$, so
$|\hat G(y)|\le c_i\,\beta(\xi)(\Lambda^{(i)})^{2}$. By \Inp{S1}, $|c_\tau|^{-1}=|\vpS{S}_\tau|/2\pi=|\kapS{S}{\tau}|/g$,
so the frequency bound, with $a=i$, gives $\sum_{y}\beta(\xi)\le\mathsf{e}$ over every finite set of non-zero
$\sT$-integers. (iii) $\hat G(0)=\int G=c_i\int F\ge0$.

\emph{Poisson summation.} By (i)--(iii), every finite sum of $|\hat G(y)|$ is at most
$c_i\bigl(\int F+\mathsf{e}\,(\Lambda^{(i)})^{2}\bigr)$, and \Inp{L3} gives
\begin{equation}\label{e:poisson}
\sum_{t\in\mathcal{V}}G(t)\le c_i\Bigl(\int F+\mathsf{e}\,(\Lambda^{(i)})^{2}\Bigr)\le c_i\int F+\mathsf{C}\,\mathsf{e}\prod_{a\le i}c_a .
\end{equation}

\emph{The integral of $F$.} If $i=0$, $\int F=J_0(\cdot,\cdot)$ is a single value, of absolute value at most
$\mathrm{En}$, and \eqref{e:poisson} is the claim. Let $i\ge1$. Writing $p\in\mathcal{A}^{i}$ as $(p',s)$ with
$p'\in\mathcal{A}^{i-1}$ and $s\in\mathcal{A}$, we have $\mathrm{hd}_i(p,1)=\mathrm{hd}_{i-1}(p',s)$, so
\[\int F=\sum_{s,s'\in\mathcal{A}}\ \sum_{p',q'\in\mathcal{A}^{i-1}}J_i\bigl(\mathrm{hd}_{i-1}(p',s),\mathrm{hd}_{i-1}(q',s')\bigr).\]
The real part of the terms with $s=s'$ is at most $\prod_{a<i}c_a\cdot(\mathrm{En}+(2i-1)\mathsf{C}\mathsf{e})$ by the
induction hypothesis for $i-1$ and $\mathcal{V}=\mathcal{A}$. For $s\ne s'$ the two sequences have at $i-1$ the entries
$u(\mathrm{pt}((\mathrm{fs}(\tau,i-1))_\tau;s))$ and the same with $s'$, and
\[\iota(\cdot)-\iota(\cdot)=\mathrm{pt}\bigl((g/\mathrm{fs}(\tau,i-1))_\tau;\,y\bigr),\qquad y=g^{-1}(s^{-1}-s'^{-1}),\]
a unit of $C$, where $y$ is a non-zero $\sT$-integer because $g^{-1}$, $s^{-1}$, $s'^{-1}$ are. By the last axiom of
a sieve system and the frequency bound for $a=i-1$, $c_\tau=g$ (admissible: $|g|^{-1}\le|\kapS{S}{\tau}|/g$ as
$|\kapS{S}{\tau}|\ge1$) and $\mathcal{X}=\{y\}$, the term of $(p',s),(q',s')$ is at most
$\mathsf{e}\,\Lambda_i(\mathrm{hd}_{i-1}(p',s))\Lambda_i(\mathrm{hd}_{i-1}(q',s'))$ in absolute value. As
$\Lambda\ge0$, the terms with $s\ne s'$ contribute at most $\mathsf{e}\,(\Lambda^{(i)})^{2}\le\mathsf{C}\mathsf{e}\prod_{a<i}c_a$.
So $\int F\le\prod_{a<i}c_a\cdot(\mathrm{En}+2i\,\mathsf{C}\mathsf{e})$, and \eqref{e:poisson} gives the claim.
\end{proof}

\begin{proof}[Proof of Theorem D]
Let $\rho_\tau=|\kapS{S}{\tau}|^{2\varrho+1}/\gS{S}\NS{S}$, $B_0=b^{23(2\varrho+1)}$, $N=\NS{S}$, $g=\gS{S}$. By \Inp{S5} and
\Inp{S4}(c), the product of the band-limit systems at the infinite places with the system $J'$ is a sieve system on
$C=\CC^{f}\times\ET{\sT}$ of the kind of Theorem~C, with $\mathsf{n}=\varrho$, the forms
$J_i(p,q)=\prod_\tau J^{\tau}_i(p_\tau,q_\tau)\cdot J'_i(p_{\sT},q_{\sT})$, the constant
$\mathrm{En}=e\prod_\tau(1+b^{-10})/\deg_\tau(S)$, and
\[\Lambda_i(p)=\prod_\tau\hw{n^{6}}{i}{p_\tau}\cdot\fhw{n^{6}}{i}{p_{\sT}},\qquad
\beta(x)=\prod_\tau\bandc{\fs}\,b^{2000n^{9}}\damp{k-n^{7}}{|x_\tau|^{-1}}\cdot\fgain{S}{\sT}{x_{\sT}}.\]
We apply Theorem~C with $i=\varrho$ and $\mathcal{V}=\mathcal{A}'$. Its left side is the left side of Theorem~D, the
entries of $\mathrm{hd}_\varrho(p,t)$ being the pairs of the entries of the $x_\tau(p,t)$ and of $y(p,t)$, and
$\prod_{a\le\varrho}c_a$ is the first factor on the right. It remains to verify the two hypotheses with suitable
$\mathsf{e}$, $\mathsf{C}$ and to compare $(2\varrho+1)\mathsf{C}\mathsf{e}$ with the error term. The exponents below
are ours; the formal proof leaves them to unification and to automation.

\emph{The frequency bound.} Let $a\in\NN$, $|c_\tau|^{-1}\le|\kapS{S}{\tau}|/g$, and $\mathcal{X}$ a finite set of
non-zero $\sT$-integers. For $y\in\mathcal{X}$,
$\beta(\ratpt{\sT}{(c_\tau/\fscale{S}{\tau}{a})_\tau}{y})$ is $\prod_\tau\bandc{\fs}\,b^{2000n^{9}}\cdot b^{n^{9}f}$ times
the term of $y$ in Proposition~\ref{p:freq} for the exponent $k-n^{7}$ and the numbers $c_\tau/\fscale{S}{\tau}{a}$.
Since $|\kapS{S}{\tau}|\ge1$,
$|\fscale{S}{\tau}{a}/c_\tau|\le|\kapS{S}{\tau}|^{2(\varrho-a)}N^{-1}\cdot|\kapS{S}{\tau}|/g\le\rho_\tau\le B_0$ by
\Inp{S1}. So Proposition~\ref{p:freq}, with $x_\tau=\rho_\tau$ and $B=B_0$, gives the frequency bound with
\[\mathsf{e}=\prod_\tau\Bigl(\bandc{\fs}\,b^{2000n^{9}}\rho_\tau^{\,k-n^{7}}(bB_0^{2})^{3}\Bigr)\cdot(b\,\rdE{E})^{3f}\,b^{n^{9}f}.\]

\emph{The weights.} Let $i\le\varrho$. For $p\in\mathcal{A}^{i}$ and $1\le j\le i$ let $m_j(p)=p_0\cdots p_{j-1}\in E^{\times}$,
a $\sT$-unit, and $F_{\tau,j}=\prod_{a<j}\fscale{S}{\tau}{a}$. By the definitions of the head weights,
\[\Lambda_i(\mathrm{hd}_i(p,1))=\prod_{j=1}^{i}\Bigl(\prod_\tau\damp{n^{6}}{|F_{\tau,j}\,m_j(p)_\tau|^{-1}}\cdot\omega(m_j(p))\Bigr),\]
with $\omega(y)=\bigl(\gix{y_{\sT}}^{1/2}\bigr)^{n^{6}}$ if $y_{\sT}\in O_{\sT}$ and $\omega(y)=0$ otherwise; in the first
case $\den{y_{\sT}}=1$. The map $p\mapsto(m_1(p),\dots,m_i(p))$ is injective, since $p_{j-1}=m_j(p)/m_{j-1}(p)$. So
the sum over $p$ of this product is at most the product over $j$ of the sums over the finite set
$\{m_j(p):p\in\mathcal{A}^{i}\}$ of non-zero $\sT$-integers, and each of these is bounded by
Proposition~\ref{p:freq} with exponent $n^{6}$, $c_\tau=F_{\tau,j}$, $x_\tau=B=N^{i}$ (indeed
$|F_{\tau,j}|^{-1}=\prod_{a<j}N|\kapS{S}{\tau}|^{-2(\varrho-a)}\le N^{j}\le N^{i}$):
\[\Lambda^{(i)}\le\Bigl(N^{ifn^{6}}\,(bN^{2i})^{3f}\,(b\,\rdE{E})^{3f}\Bigr)^{i}.\]
With $N\le b^{60n}$, $\rdE{E}\le b$, $i\le\varrho<n$ and $n^{6}+6\le7n^{6}$ this gives
$(\Lambda^{(i)})^{2}\le b^{(840n^{9}+18n)f}$. On the other side, $\fscale{S}{\tau}{a}\le|\kapS{S}{\tau}|^{2\varrho}\le b^{46\varrho}$
and $\volO{E}=(2/\rdE{E})^{f}\ge b^{-f}$, so $c_a\ge b^{-(92\varrho+1)f}$ and
$\prod_{a<i}c_a\ge b^{-93n^{2}f}$. Hence the hypothesis of Theorem~C holds with $\mathsf{C}=b^{1000n^{9}f}$.

\emph{The error term.} By the admissibility inequality of \Inp{S1},
$\prod_\tau\bandc{\fs}\,\rho_\tau^{\,k-n^{7}}\le\bdeg{S}^{-1}b^{-7kf}$. Also $2\varrho+1\le b$. So
\begin{gather*}(2\varrho+1)\,\mathsf{C}\,\mathsf{e}\le\bdeg{S}^{-1}\,b^{-7kf}\cdot b^{\,1+1000n^{9}f+2000n^{9}f+3f(1+46(2\varrho+1))+6f+n^{9}f}\\
\le\bdeg{S}^{-1}\,b^{-7kf+1+3286n^{9}f},\end{gather*}
and $1+3286n^{9}f\le kf-10$ because $k\ge1000(2\varrho+2)^{9}=512000\,n^{9}$. This is the error term
$\bdeg{S}^{-1}b^{-(6kf+10)}$.
\end{proof}

\sect{The standard level at a finite place}{s:level}
The tower method of unit~132 bounds integrals over a unitary group floor by floor, given a \emph{level} and a
\emph{gauge} (Input~\Inp{R1}). This section provides them at a finite place, for the lattice of the filling data:
the form is $\Tlatg{\fs}{N}$, the box is the group of the integral isometries, and all sizes are integrality
conditions. The first two results concern any place.

\begin{definition}[the tail vector]\label{d:tail}
Let the ring $A$ be a commutative ring with an involution, $\fs$ a shape, $\pi$ a unit of $A$, $t\in A$ and
$x\in A^{\kappa}$. Then
\[\tailv{\pi}{t}{x}=\blb{\fs}{\pi}^{-1}\cdot\bigl(0,\ (x,\,-(\pi\bar\pi)^{-1}t),\ 0\bigr)\]
is the vector $\sum_jx_jx_j-(\pi\bar\pi)^{-1}t\,z$ of $V_\varrho$, written in the lattice basis.
\end{definition}
\leanline{\lean{tailVector}}

\begin{lemma}\label{l:tail}
Let $A$ be a commutative ring with an involution, $\fs$ a shape, $\pi$ a unit of $A$ and $N\in A$ with
$\pi\bar\pi=N$; let $t\in A$, $x\in A^{\kappa}$, $J=\JT{\Tlatg{\fs}{N}}$ and $e=\tailv{\pi}{t}{x}$. Then
$e^{*}Je=x^{*}x-(\pi\bar\pi)^{-1}t\,\bar t$, and the coordinate of $Je$ at $z$ is~$t$.
\end{lemma}
\begin{proof}
By \Inp{S3}, $J=\blb{\fs}{\pi}^{*}J'\blb{\fs}{\pi}$ with $J'=\JT{\Trg{\fs}{N}}$. So with
$v=(0,(x,-N^{-1}t),0)$ we get $e^{*}Je=v^{*}J'v=x^{*}x-N\cdot N^{-1}\bar t\,N^{-1}t$, as $\bar N=N$. Further
$Je=\blb{\fs}{\pi}^{*}J'v$, and $J'v$ has the coordinate $-N\cdot(-N^{-1}t)=t$ at $z$; the column of $\blb{\fs}{\pi}$ at
$z$ is the basis vector $z$ itself, so the coordinate of $Je$ at $z$ is $t$.
\end{proof}
\leanline{\lean{formtailVector}}

\begin{proposition}[slabs near the quadric]\label{p:slab}
Let $A$ be a place (a commutative ring with an involution, a topology, a Haar measure and a place structure~$P$),
$\Omega$ a finite set, $\mathfrak{m}\in\NN$, $\pi$ a unit of $A$ and $N\in A$ with $\pi\bar\pi=N$ such that
$\JT{\Tlatg{\fs}{N}}$ is a hermitian unit. Let $F\in\HypB{\Omega}{\Tlatg{\fs}{N}}{\mathfrak{m}}$ be a hyperbolic basis,
$\Phi$ a continuous function with compact support on the space of vectors, with values in $[0,\infty)$, and
$c\in[0,\infty]$ with $\int\Phi\,\dd\quadric{F}{\varrho}<c$. Then there is a neighbourhood $\mathcal{N}$ of $0$ in $A$
such that for every measurable set $Z\subseteq\mathcal{N}$
\[\int_A\int_{\{y\in A^{\kappa}\,:\ y^{*}y-((\pi\bar\pi)^{-1}t\bar t-N)\in Z\}}\Phi\bigl(\tailv{\pi}{t}{y}\bigr)\,\dd y\,\dd t\;\le\;
\modA(\pi)\cdot\nu_A\{\beta\in A^{-}:\delta_A^{-1}\beta\in Z\}\cdot c.\]
\end{proposition}
\begin{proof}
Write $q(v)=h(v,v)$ on $V_\varrho$, so that the quadric of $\quadric{F}{\varrho}$ is $\{q=-N\}$, the fibre of
$z$. By \Inp{R1} the Haar measure of $V_\varrho$ disintegrates along $q$: for $\Psi\ge0$ measurable,
$\int_{V_\varrho}\Psi(q(v)+N,v)\,\dd v=\int_{A^{-}}\int\Psi(\delta_A^{-1}\beta,v)\,\dd\delta_\beta(v)\,\dd\nu_A(\beta)$, where
$\delta_\beta$ is the quadric measure on the fibre $\{q+N=\delta_A^{-1}\beta\}$, $\delta_0=\quadric{F}{\varrho}$, and
$\beta\mapsto\int\Phi\,\dd\delta_\beta$ is continuous. Hence there is a neighbourhood $\mathcal{N}$ of $0$ with
$\int\Phi\,\dd\delta_\beta<c$ whenever $\delta_A^{-1}\beta\in\mathcal{N}$. Let $Z\subseteq\mathcal{N}$ be measurable.
Applying the disintegration to $\Psi(s,v)=\mathbf{1}_Z(s)\Phi(v)$,
\[\int_{V_\varrho}\mathbf{1}_Z(q(v)+N)\,\Phi(v)\,\dd v\le\nu_A\{\beta:\delta_A^{-1}\beta\in Z\}\cdot c.\]
The map $(t,y)\mapsto\tailv{\pi}{t}{y}$ from $A\times A^{\kappa}$ to the coordinates of $V_\varrho$ in the lattice
basis is linear and continuous (the formal text records the continuity as a lemma). Its matrix has the columns of
the base change of the $x_j$ and the column $-(\pi\bar\pi)^{-1}z$, so its determinant is $-\pi^{-1}$, and it
multiplies Haar measure by $\modA(\pi)$, in the sense that
$\int\int\Psi'(\tailv{\pi}{t}{y})\,\dd y\,\dd t=\modA(\pi)\int_{V_\varrho}\Psi'(v)\,\dd v$; and by Lemma~\ref{l:tail},
$q(\tailv{\pi}{t}{y})+N=y^{*}y-((\pi\bar\pi)^{-1}t\bar t-N)$. The two displays give the claim.
\end{proof}
\leanline{\lean{eventuallymassnearquadricle}\sep\lean{continuoustailVector}}

From here on $A$ is a semilocal field with an involution and an integral place structure (Input~\Inp{R3}), and
$O_A$ its ring of integers. For an $\fs$-form $T$ with $\JT{T}$ a hermitian unit and a unit~$g$, $\uLev{g}{T}$ is the
group of the $u\in\GT{T}$ that are integral with integral inverse and satisfy $g^{-1}(u-1)$ integral.

\begin{definition}[lattice levels and standard levels]\label{d:stdlevel}
Let $A$ be a semilocal field with an involution and an integral place structure, $\fs$ a shape, $\Omega$ a finite
set, $\mathfrak{m}\in\NN$, and $T$ an $\fs$-form over $A$ with $\JT{T}$ a hermitian unit.
\begin{enumerate}[label=(\alph*),leftmargin=2.2em,itemsep=1pt,topsep=2pt]
\item For a level $\cL\in\Levels{\Omega}{T}{\mathfrak{m}}$, its \emph{congruence subgroup} is
  $\cgp{\cL}=\uLev{\gL{\cL}}{T}$.
\item A level $\cL\in\Levels{\Omega}{T}{\mathfrak{m}}$ is a \emph{lattice level}, \m{\IsLat{\cL}}, if: (LL1) $\JT{T}$ is
  unimodular over $O_A$, the basis matrix $\EL{\cL}$ has entries in $O_A$, $\gL{\cL}\in O_A$ and
  $\overline{\gL{\cL}}=\gL{\cL}$; (LL2)
  $\vol(O_A^{\times})\le\floorMeasure{\cL}{\varrho}(\BoxL{\cL})$, $\BoxL{\cL}=\uLev{1}{T}$,
  $\floorMeasure{\cL}{n}(\BoxL{\cL})=1$, $\Ball{\cL}=O_A^{\Ix{\fs}}$, and
  $\Posn{\cL}=\{y^{-1}uy:y\in\GT{T},\ u\in\UT{T}\}\cdot\cgp{\cL}$; (LL3) there is a vector $x$ with $z$-coordinate $0$
  such that $h_T(x,\cdot)$ is the $z$-coordinate.
\item A level $\cL\in\Levels{\Omega}{T}{\mathfrak{m}+1}$ is a \emph{standard level}, \m{\IsStd{\cL}}, if it is a lattice
  level, $\zetaA{\mathfrak{m}}^{\,n}\le1+2^{-n}$, $\mathfrak{m}<9n$, $13n\le|\Omega|$, and $\GT{T}$ has a Haar measure that
  is right invariant.
\end{enumerate}
\end{definition}
\leanline{\lean{congGroup}\sep\lean{IsLatticeLevel}\sep\lean{IsStdLevel}}

\begin{proof}[Proof of Theorem E]
Let $T=\Tlatg{\fs}{N}$ and $J=\JT{T}$; by \Inp{S3}, $J$ is unimodular over $O_A$. In the lattice basis the middle block
has the orthonormal vectors $x_j$, $j\ne k_0$, the isotropic vector $v_0$ at the index $k_0$, and $z$, with
$h(v_0,z)=-1$, $h(z,z)=-N$. We describe the level and say for each axiom of \Inp{R1} and of
Definition~\ref{d:stdlevel}(b) why it holds; the formal proof verifies them in one construction.

\emph{The hyperbolic basis.} Let $\tau_A$ be the trace-one element and $c\in O_A$ with $\bar cc=-1$ (\Inp{R3}). The
vectors $-v_0$ and $-N\tau_Av_0+z$ form a hyperbolic pair: $h(-v_0,z)=1$ and
$h(-N\tau_Av_0+z,-N\tau_Av_0+z)=N\bar\tau_A+N\tau_A-N=0$. Since $\kk{\fs}-1=|\Omega|+2\mathfrak{m}$, the vectors $x_j$,
$j\ne k_0$, can be split into $\mathfrak{m}$ pairs $(x,x')$ and a family indexed by~$\Omega$; each pair gives the
hyperbolic pair $x+cx'$, $\tau_Ax-\bar\tau_Acx'$. This basis matrix is integral, and $\modA(\det T)=1$ because $T$ is
unimodular; so it is a hyperbolic basis in $\HypB{\Omega}{T}{\mathfrak{m}}$.

\emph{The measure of the floor~$\varrho$.} Put $x=-v_0$: then $h(x,\cdot)$ is the $z$-coordinate and the
$z$-coordinate of $x$ is~$0$, which is (LL3). Let $\varphi$ be the parametrisation of the quadric
$\{v\in V_\varrho:h(v,v)=h(z,z)\}$ by the flat chart of the library, $p=((t,w),\beta)\mapsto\varphi(p)$, in the
coordinates of the hyperbolic basis; the $z$-coordinate of $\varphi(p)$ is~$t$. For a vector $e$ let $\mathrm{E}(e)$
be the Eichler element of \Inp{R3}(e), formed with~$x$. Define $\floorMeasure{\cL}{\varrho}$ as the image of the chart
measure $\dd t\,\dd w\,\dd\nu_A(\beta)$ under $p\mapsto\mathrm{E}(\varphi(p))$, a measurable map. By \Inp{R3}(e) it is
carried by $\Hl{\varrho}{T}$, and since $t$ is a unit for almost every $p$, its image under $\sigma\mapsto\sigma^{-1}z$
is the image of the chart measure under $\varphi$, which is $\quadric{\cL}{\varrho}$.

\emph{The measure of the floor $n$.} The group $\HT{T}$ is closed in $\GT{T}$ and $\HT{T}\cap\uLev{1}{T}$ is a compact
open subgroup of it. Let $\floorMeasure{\cL}{n}$ be the image under inversion of the left Haar measure of $\HT{T}$ that
gives this subgroup mass~$1$. It is carried by $\HT{T}$, invariant under right translations by $\HT{T}$, and
$\floorMeasure{\cL}{n}(\uLev{1}{T})=1$. For $g\in\GT{T}$ let $\Sigma_g=\{h\in\HT{T}:uhg\in\uLev{1}{T}\text{ for some }u\in\UT{T}\}$.
If $h,h'\in\Sigma_g$, with $uhg$ and $u'h'g$ in $\uLev{1}{T}$, then $u'h'h^{-1}u^{-1}\in\uLev{1}{T}$; as $h'h^{-1}$
normalises $\UT{T}$, this is $u''\cdot h'h^{-1}$ with $u''\in\UT{T}$, and then $h'h^{-1}\in\uLev{1}{T}$ (a lemma of the
library on the Bessel subgroup: a product $u''h''\in\uLev{1}{T}$ with $u''\in\UT{T}$, $h''\in\HT{T}$ has
$h''\in\uLev{1}{T}$). So $\Sigma_g$ lies in one right coset of $\HT{T}\cap\uLev{1}{T}$ and has mass at most~$1$.

\emph{The remaining data.} The unit is $g$, with $\bar g=g$; the ball is $O_A^{\Ix{\fs}}$, which is open and compact and
contains~$0$; the box is $\uLev{1}{T}$, which is compact (\Inp{R3}(b)); the positions are
$\{y^{-1}uy\}\cdot\uLev{g}{T}$, an open set, hence measurable. This is a level $\cL$ with $\gL{\cL}=g$, and (LL1) and
the equalities of (LL2) hold by construction. For the inequality of (LL2), consider the chart points $p$ with
$t\in O_A^{\times}$, $w$ integral and $\beta\in A^{-}\cap O_A$; they have chart measure $\vol(O_A^{\times})$, since
$\vol(O_A)=1$ and $\nu_A(A^{-}\cap O_A)=1$. For such $p$ the vector $\varphi(p)$ is integral with $z$-coordinate a
unit of $O_A$, so the matrix $\mathrm{E}(\varphi(p))$ of \Inp{R3}(e) is integral; an integral isometry of a
unimodular form lies in $\uLev{1}{T}$. Hence $\floorMeasure{\cL}{\varrho}(\BoxL{\cL})\ge\vol(O_A^{\times})$.
\end{proof}

\begin{definition}[the gauge of a standard level]\label{d:gauge}
Let $A$ be a semilocal field with an involution and an integral place structure, $\fs$, $\Omega$, $\mathfrak{m}$, $T$
as in Definition~\ref{d:stdlevel}, let $\cL\in\Levels{\Omega}{T}{\mathfrak{m}+1}$ be a standard level, \m{\IsStd{\cL}}, and
let $\alpha$ be a unit of $A$ with $\alpha\in O_A$ and $\bar\alpha=\alpha$ (by default $\alpha=1$). The \emph{gauge of
depth $\alpha$} is
\[\stdg{\cL}{\alpha}\in\Gauges{\cL}:\qquad\Gamma_j=\{x\in A:\alpha^{j}x\in O_A\}\quad(j\in\NN),\qquad
\Lambda=\frac{2}{\modA(\alpha)} .\]
That these data satisfy the axioms of a gauge is part of the definition in the formal text; it is
Lemma~\ref{l:gauge}.
\end{definition}

\begin{lemma}\label{l:gauge}
In the situation of Definition~\ref{d:gauge}, the sets $\Gamma_j=\alpha^{-j}O_A$ and the constant
$\Lambda=2/\modA(\alpha)$ form a gauge of~$\cL$.
\end{lemma}
\begin{proof}
We list the axioms of unit~132 and the reason for each; (G10) is described and not carried out in every detail.
Note first: $\modA(\alpha)\le1$, so $\Lambda\ge2$; $\zetaA{\mathfrak{m}}\le2$ because its $n$-th power is at most
$1+2^{-n}\le2$; and $\modA(\alpha^{j})^{-1}\le\Lambda^{j}$. The sets satisfy the axioms of a \emph{size}:
$\Gamma_i\Gamma_{i'}\subseteq\Gamma_{i+i'}$; each $\Gamma_i$ is an $O_A$-module contained in $\Gamma_{i+1}$, so sums of
its elements lie in $\Gamma_{i+1}$; it is stable under $x\mapsto-x$ and, as $\bar\alpha=\alpha$, under the involution;
it contains $0$ and $1$; the entries of $\JT{T}$ and of its inverse lie in $O_A\subseteq\Gamma_1$ by (LL1).
(G1) $\Gamma_j$ is compact, the preimage of $O_A$ under a homeomorphism. (G2) $\Lambda\ge2$, and the number
$\mathfrak{m}+1$ of pairs of the level is at most $9n$. (G3) The Plancherel constant is $1$ and $\Lambda<\infty$.
(G4) $\gL{\cL}$, the entries of $\EL{\cL}$, and $\tau_A\delta_A^{-1}$ lie in $O_A\subseteq\Gamma_1$, by (LL1) and
\Inp{R3}. (G5), (G6) The elements of the box $\uLev{1}{T}$ and the vectors of the ball $O_A^{\Ix{\fs}}$ are integral.
(G7) If $u$ is a unit in $\Gamma_j$, then $\modA(\alpha^{j}u)\le1$, so $\modA(u)\le\Lambda^{j}$. (G8)
$\vol(\Gamma_j)=\modA(\alpha^{j})^{-1}\le\Lambda^{j+1}$ by \Inp{R3}(a). (G9) For a quadric in $|\Omega'|+1$ hyperbolic
planes with $\mathfrak{m}+1\le|\Omega'|\le\mathfrak{m}+1+\varrho$, the flat quadric measure of the vectors with
coordinates in $\Gamma_j$ is at most $\zetaA{\mathfrak{m}}\,\modA(\alpha^{j})^{-(2|\Omega'|+1)}\le\Lambda^{\,j(2|\Omega'|+1)+1}$
by the cone bound \Inp{R3}(c), and $j(2|\Omega'|+1)+1\le99n(j+1)$ since $|\Omega'|\le10n$. (G10) (the volume
majorant) Let $x$ be in the box, $y$ a position, $v$ a matrix with $\varrho$ rows, and $W$ the set of the
$w\in A^{\Omega}$ with $vw=0$ for which the two vectors $x^{-1}\mathsf{O}w$ and $(xy)^{-1}\mathsf{O}w$ are close in the
sense of the level ($\mathsf{O}$ the orthonormal part of the basis). The axiom asks for a measurable
$\varphi\ge0$ on $A^{\Omega}$ that bounds, at every $u$, the inverse volume of $\{u^{*}w:w\in W\}$ and the volume of
a set of $c\in A$ defined by $W$, and whose power $\volumeExponent{\cL}$ has integral at most $\Lambda^{99n(j+1)}$ over
the vectors with coordinates in~$\Gamma_j$. The formal proof writes $y=y_0^{-1}u_0y_0\cdot k'$ with $u_0\in\UT{T}$,
$k'\in\cgp{\cL}$, and uses the elementary divisor form over $O_A$ (a lemma of the library on valuation pairs) to
find $M\in\GL_\Omega(O_A)$ and a submatrix $Y$ of $M$, with $|\Omega|-(3n+1+\mathfrak{m})+\mathfrak{m}+5$ columns
by the count $13n\le|\Omega|$, such that $Yt\in W$ for every integral~$t$: the columns of $Y$ are killed by $v$, by
the rows of $y_0x^{-1}\mathsf{O}$ at the $e^*_a$ and by the rows of $(u_0^{-1}-1)y_0x^{-1}\mathsf{O}$ at the~$e_a$. Then
$\varphi(u)=\mathrm{ct}(Y^{*}u)$, the content of \Inp{R3}(d), bounds both quantities (by the definition of the
content for the first, and by the integrality of the pairings of integral vectors for the second), and the bound
for the integral is the content moment \Inp{R3}(d), with $\zetaA{\mathfrak{m}}\le2$ and
$\modA(\alpha^{j})^{-1}\le\Lambda^{j}$.
\end{proof}
\leanline{\lean{gauge}}

\begin{proposition}[Heisenberg masses of congruence subgroups]\label{p:heis}
Let $\cL\in\Levels{\Omega}{T}{\mathfrak{m}+1}$ be a standard level, \m{\IsStd{\cL}}, over a semilocal field $A$ with an
involution and an integral place structure, let $g$ be a unit $g$ of $A$ with $g\in O_A$ and $\bar g=g$, let
$a<\varrho$, and let $\mathsf{N}$ be a subgroup of $\GT{T}$ with $\mathsf{N}\subseteq\Nb{a}{T}$. Then
\[\eta_a\bigl\{y:uy\in\uLev{g}{T}\text{ for some }u\in\mathsf{N}\bigr\}=\floorModulus{\cL}{a}{g}
=\modA(g)^{k+1+2(\varrho-a-1)}\,\modA(g)^{1/2}.\]
\end{proposition}
\begin{proof}
$\eta_a$ is the image of $\dd v\,\dd\beta$ on $V_{a+1}\times A^{-}$ under $(v,\beta)\mapsto\mathrm{n}_a(v,\beta)$. We
claim that $u\,\mathrm{n}_a(v,\beta)\in\uLev{g}{T}$ for some $u\in\mathsf{N}$ if and only if $g^{-1}v$ is integral and
$g^{-1}\beta\in O_A$. If the latter holds, then $\mathrm{n}_a(v,\beta)-1$, whose entries are formed from the
coordinates of $v$, from $\beta-\tau_Ah_T(v,v)$ and from the entries of the integral matrix $\JT{T}$, is divisible
by~$g$; an integral isometry of a unimodular form that is congruent to $1$ modulo $g$ lies in $\uLev{g}{T}$, so
$u=1$ works. Conversely, apply \Inp{R2}(b) to the gauge $\stdg{\cL}{1}$, all of whose sets are $O_A$, and to the unit
$g$: the entries of $g^{-1}(u\,\mathrm{n}_a(v,\beta)-1)$ are integral, so $g^{-1}v$ is integral and
$g^{-1}(\beta-\tau_Ah_T(v,v))\in O_A$; since $h_T(v,v)\in g^{2}O_A$, also $g^{-1}\beta\in O_A$. By \Inp{R3}(a), the
set of the $v\in V_{a+1}$ with $g^{-1}v$ integral has volume $\modA(g)^{\dim V_{a+1}}$, with
$\dim V_{a+1}=k+1+2(\varrho-a-1)$, and $\nu_A\{\beta:g^{-1}\beta\in O_A\}=\modA(g)^{1/2}$, which uses $\bar g=g$.
\end{proof}
\leanline{\lean{heisenbergMeasureeqfloorModulus}}

\begin{corollary}\label{c:unfoldU}
Let $\cL\in\Levels{\Omega}{T}{\mathfrak{m}+1}$ be a standard level, \m{\IsStd{\cL}}, over a semilocal field with an
involution and an integral place structure, and let $g$ be a unit $g$ of $A$ with $g\in O_A$ and $\bar g=g$.
\begin{enumerate}[label=(\alph*),leftmargin=2.2em,itemsep=1pt,topsep=2pt]
\item For every $j\le\varrho$: $\etab{j}\bigl(\uLev{g}{T}\bigr)=\prod_{l<j}\floorModulus{\cL}{l}{g}$.
\item For every $c\in\GT{T}$ and every $s$-finite measure $\theta$ on $\GT{T}$:
\[(\etab{\varrho+2}*\theta)\{y:yc\in\uLev{g}{T}\}=\prod_{l<\varrho}\floorModulus{\cL}{l}{g}\cdot
\theta\{y:uyc\in\uLev{g}{T}\text{ for some }u\in\UT{T}\}.\]
\end{enumerate}
\end{corollary}
\begin{proof}
The library proves the unfolding $(\etab{j}*\theta)\{y:yc\in\uLev{g}{T}\}=\etab{j}(\uLev{g}{T})\cdot
\theta\{y:uyc\in\uLev{g}{T}\text{ for some }u\in\Nb{j}{T}\}$ for all $j$, $\theta$, $c$ (unit~123; the measures
$\etab{j}$ are invariant under $\Nb{j}{T}$). (a) By induction on $j$: $\etab{j+1}=\etab{j}*\eta_j$, and the
unfolding with $\theta=\eta_j$, $c=1$ and Proposition~\ref{p:heis} with $\mathsf{N}=\Nb{j}{T}$ give the factor
$\floorModulus{\cL}{j}{g}$. (b) The unfolding for $j=\varrho+2$, where $\Nb{\varrho+2}{T}=\UT{T}$ and
$\etab{\varrho+2}=\etab{\varrho}$, and (a) for $j=\varrho$.
\end{proof}
\leanline{\lean{ascConvunitaryLevel}\sep\lean{unfoldbesselU}}

\begin{proposition}[the mass of the integral group, from above]\label{p:mass}
Let $\cL\in\Levels{\Omega}{T}{\mathfrak{m}+1}$ be a standard level, \m{\IsStd{\cL}}, over a semilocal field with an
involution and an integral place structure. Then $\groupMeasure{\cL}(\BoxL{\cL})\le1+2^{-n}$.
\end{proposition}
\begin{proof}
By \Inp{R2}(a) and Corollary~\ref{c:unfoldU}(b) for $g=1$, $c=1$ and $\theta=\upperMeasure{\cL}{0}$, where all
$\floorModulus{\cL}{l}{1}$ are $1$ and $\uLev{1}{T}=\BoxL{\cL}$,
\begin{gather*}\groupMeasure{\cL}(\BoxL{\cL})=\upperMeasure{\cL}{0}\{y:uy\in\BoxL{\cL}\text{ for some }u\in\UT{T}\}\\
=\upperMeasure{\cL}{0}\{y:\IsNear{\cL}{\varrho+2}{y}\}.\end{gather*}
Apply \Inp{R2}(c) to the gauge $\stdg{\cL}{1}$, whose sets are all $O_A$, of volume~$1$: for $b\le\varrho$ the chart
of $V_b$ has $|\Omega_b|+1$ hyperbolic planes with $|\Omega_b|=\mathfrak{m}+1+\varrho-b\ge\mathfrak{m}+1$, so by the cone
bound \Inp{R3}(c) with $\beta=1$ the integral vectors have flat quadric measure at most $\zetaA{\mathfrak{m}}$.
Hence the right side is at most $\zetaA{\mathfrak{m}}^{\,n}\le1+2^{-n}$.
\end{proof}
\leanline{\lean{groupMeasureboxSetle}}

\begin{proposition}[the mass of the integral group, from below]\label{p:stack}
Let $\cL\in\Levels{\Omega}{T}{\mathfrak{m}+1}$ be a standard level, \m{\IsStd{\cL}}, over a semilocal field with an
involution and an integral place structure, and let $i$ and $l$ be natural numbers $i$, $l$ with $l+i=\varrho$. Then
$\vol(O_A^{\times})^{i+1}\le\stackMeasure{\cL}_{l,i+2}(\BoxL{\cL})$.
\end{proposition}
\begin{proof}
Write $K=\BoxL{\cL}=\uLev{1}{T}$, a group. For $s$-finite measures $\sigma$, $\theta$ on $\GT{T}$ we have
$\sigma(K)\theta(K)\le(\sigma*\theta)(K)$, because the product of two elements of $K$ lies in~$K$. Induction on~$i$.
For $i=0$, $l=\varrho$: $\stackMeasure{\cL}_{\varrho,2}=\floorMeasure{\cL}{n}*\floorMeasure{\cL}{\varrho}$ gives $K$ mass
at least $1\cdot\vol(O_A^{\times})$ by (LL2). For the step, $l+i+1=\varrho$:
$\stackMeasure{\cL}_{l,i+3}=\stackMeasure{\cL}_{l+1,i+2}*(\eta_l*\floorMeasure{\cL}{l})$, so
\[\stackMeasure{\cL}_{l,i+3}(K)\ge\vol(O_A^{\times})^{i+1}\cdot\eta_l(K)\cdot\floorMeasure{\cL}{l}(K),\qquad
\floorMeasure{\cL}{l}(K)\ge\varpi_l(K)\,\eta_l(K).\]
Here $\eta_l(K)=\floorModulus{\cL}{l}{1}=1$ by Proposition~\ref{p:heis} for $g=1$ and the trivial subgroup, and
$\varpi_l(K)\ge\vol(O_A^{\times})$: for $c\in O_A^{\times}$ the Weyl element $\wo{l}{c}$ is an integral isometry, so
its inverse lies in $K$, and the density $\modA(c)^{k+1+2(\varrho-l-1)}$ of $\varpi_l$ is $1$ at such~$c$.
\end{proof}
\leanline{\lean{lestackMeasureboxSet}}

\begin{corollary}\label{c:haar}
For a standard level $\cL$ over a semilocal field with an involution and an integral place structure,
$\groupMeasure{\cL}$ is a Haar measure of $\GT{T}$.
\end{corollary}
\begin{proof}
By Definition~\ref{d:stdlevel}(c), $\GT{T}$ has a Haar measure $\nu$ that is right invariant. The measure
$\groupMeasure{\cL}$ is right invariant and $s$-finite (\Inp{R1}), hence equal to $d\,\nu$ for some
$d\in[0,\infty]$, by the uniqueness of Haar measure. It is not zero, so $d\ne0$; and $d\ne\infty$ because
$\groupMeasure{\cL}(\BoxL{\cL})<\infty$ by Proposition~\ref{p:mass} while $\nu(\BoxL{\cL})>0$, the box being open.
\end{proof}
\leanline{\lean{isHaarMeasuregroupMeasure}}

\sect{The Bessel systems of admissible data}{s:degree}
In this section $P$ is a pair of parameters and $D$ admissible data for~$P$ (Chapter~III of the write-up), with the
notation of Input~\Inp{D1}. The file of the formal text that holds Theorem~F opens with the last step of the seesaw
argument of unit~141. It does not concern Bessel systems; it is stated here, in the words of the note on the seesaw
identity, because it is a member of this unit.

\begin{proposition}[the last stage of the seesaw]\label{p:last}
Let $P$ be parameters, $D$ admissible data for $P$, and $M$ a natural number $M$ with $M=\Dm{D}$. Let $\Stg{D}{M}$ be
the space of the smooth algebraic modular forms of stage $M$ on the definite space $X_M(D)$, $\FormY{D}{M}$ the space
of the forms on the standard definite space $\Ys{D}{M}$ of its type, with the Petersson product, and
$\gYlast{D}$ the grading of $\Ys{D}{M}$ with an empty string (unit~141). There is a $\CC$-linear map
$R\colon\Stg{D}{M}\to\FormY{D}{M}$ with the following two properties.
\begin{enumerate}[label=(\arabic*),leftmargin=2.2em,itemsep=1pt,topsep=2pt]
\item For every tame character $\vartheta$ of~$D$, every finitely supported function $\Psi$ on the tuples of
  translations, every form $y$ of principal level on $\Ys{D}{0}$ and every column $\varphi\in\Col{D}$: if
  $\chainY{D}{\vartheta}{M}(y)=R\bigl(\chcomb{D}{\vartheta}{\Psi}{M}(\varphi)\bigr)$, then
  $\iot{D}{\vartheta}(y)=\Bmu{D}{\vartheta}{\Psi}(\varphi)$.
\item For every linear functional $\mathsf{t}$ on the values of the forms of stage $M$ there is a complex number
  $C_{\mathsf{t}}$ such that for every $\Xi\in\FormY{D}{M}$, every $\psi\in\Stg{D}{M}$ and every function $F$ on
  $G(\Ys{D}{M})$ with $F(\incl{\gYlast{D}}(h))=\coef{\mathsf{t}}{\psi}(\mathrm{diag}(h,1))$ for all~$h$:
  \[\int_{\Gamma(\Ys{D}{M})\backslash G(\Ys{D}{M})}F\,\overline{\Xi}=C_{\mathsf{t}}\cdot\langle\Xi,R\psi\rangle.\]
\end{enumerate}
Here $\chainY{D}{\vartheta}{M}$ is the chain of the first $M$ indefinite lifts, $\iot{D}{\vartheta}$ the same chain read
as an algebraic modular form on $\Ys{D}{M}$, $\chcomb{D}{\vartheta}{\Psi}{M}$ the combination of the chains of definite
lifts with the translated vacua, $\Bmu{D}{\vartheta}{\Psi}$ its restriction to $\Ys{D}{M}$, and
$\coef{\mathsf{t}}{\psi}(g)=\mathsf{t}(g_\infty^{-1}\psi(g_f))$ the coefficient of $\psi$ at~$\mathsf{t}$, all as defined in
the notes of unit~141.
\end{proposition}
\begin{proof}
This is an assembly; its content lies in three inputs, which we did not read beyond their statements. Let
$\mathrm{alg}$ be the identification of the forms on the standard definite space $\Ys{D}{M}$ of its type with
algebraic modular forms (background library; the space is standard and its parameter is a Harish-Chandra parameter by
units 140 and~147), and $\mathrm{res}_D$ the restriction of unit~145 from stage $M$ to $\Ys{D}{M}$: restriction to
the block diagonal elements $\mathrm{diag}(h,1)$, a branching matrix, and a power of the conjugates of the level
element. Put $R=\mathrm{alg}^{-1}\circ\mathrm{res}_D$. (1) By definition
$\iot{D}{\vartheta}=\mathrm{alg}\circ\chainY{D}{\vartheta}{M}$ and
$\Bmu{D}{\vartheta}{\Psi}=\mathrm{res}_D\circ\chcomb{D}{\vartheta}{\Psi}{M}$ on columns, so the hypothesis gives
$\iot{D}{\vartheta}(y)=\mathrm{res}_D(\chcomb{D}{\vartheta}{\Psi}{M}(\varphi))$. (2) The branching formula of the
background library (unit~27) says that the integral of a coefficient of an algebraic form along the block embedding
against an algebraic form is a constant, depending on the functional only, times the inner product with the
branching restriction; the integral over the quotient is identified with the sum over classes that defines the
Petersson product of algebraic forms, and the constants of $\mathrm{res}_D$ and the number of classes are absorbed
in~$C_{\mathsf{t}}$. The hypothesis on $F$ enters through the values of $F$ on the image of $\incl{\gYlast{D}}$, on which
the two integrands agree.
\end{proof}
\leanline{\lean{existslaststagemap}}

\begin{lemma}[differences of the row parameter]\label{l:rowpar}
Let $P$ be parameters, $D$ admissible data for $P$, $\tau$ an infinite place of $\ED{D}$ and $i<l$ natural numbers.
Then
\[1\le\rowPar{D}{\tau}{i}-\rowPar{D}{\tau}{l}\le\wt{D}{l+2-\kD{D}}{\tau}.\]
\end{lemma}
\begin{proof}
Write $k=\kD{D}$, $r=\Dr{D}$, $h_2=\hhD{D}=r^{8}$, $w_j=\wt{D}{j}{\tau}=b_\ast+2jh_2+h_2$ (\Inp{D1}); note $k\le r$ and
$r\ge1$. If $l+1<k$, both entries are in the first range, the difference is $l-i\in[1,l]$, and
$w_{0}\ge h_2\ge r>l$. If $i+1<k\le l+1$, the difference is $k/2-1-i+b_\ast+2(l+2-k)h_2$; it is at most
$w_{l+2-k}$ because $k/2-1-i\le r\le h_2$, and at least $1-k/2+2h_2\ge1$ because $i\le k-2$ and $l+2-k\ge1$. If
$k\le i+1$, the difference is $w_{l+2-k}-w_{i+2-k}=2(l-i)h_2$, which lies between $1$ and $w_{l+2-k}$.
\end{proof}
\leanline{\lean{rowParametersubmem}}

\begin{proof}[Proof of Theorem F]
Let $S=\SD{D}$; by \Inp{D1}(a) it fits $\Ys{D}{0}$, so by (F1) the Bessel system at $\tau$ has the type of the
parameter $l_\tau=\parW{D}{\tau}{\cdot}$ and of the weight of the columns at~$\tau$.

(a) $\#\Bc{S}=\prod_\tau\#I_\tau(S)$. By \Inp{R4}(ii) there is an invertible matrix between the coordinates $I_\tau(S)$
and the Weyl coordinates, so $\#I_\tau(S)$ is the dimension of the representation, which by \Inp{D1}(d) is the
symmetric power of order $a_\tau\le\aD{D}$ of $\CC^{k}$: the binomial coefficient
$\binom{a_\tau+k-1}{k-1}\le(k+\aD{D})^{k}$. With $(k-1)!\le k^{k}$ the left side is at most
$\bigl(((k+\aD{D})k)^{k}\bigr)^{f}$, and $(k+\aD{D})k\le\scb{D}^{5/4}$ by \Inp{D1}(c).

(b) Write $\rho(i)=\rowPar{D}{\tau}{i}$ for a fixed $\tau$. By \Inp{R4}(i), $\deg_\tau(S)$ is Weyl's number of
$l_\tau$. Since $l_{\tau,j}=2\rho(\nu(j))$ and $\rho$ is strictly decreasing (Lemma~\ref{l:rowpar}), the pairs with
$l_{\tau,j}<l_{\tau,i}$ are the pairs with $\nu(i)<\nu(j)$, and
\[\deg_\tau(S)=\prod_{l<r}\ \prod_{i<l}\bigl(\rho(i)-\rho(l)\bigr)\Big/\prod_{l<r}l!\;=\;
\prod_{l=k-1}^{r-1}\Bigl(\prod_{i<l}\bigl(\rho(i)-\rho(l)\bigr)\Big/l!\Bigr),\]
because for $l+1<k$ the differences are $l-i$ and their product over $i<l$ is~$l!$. Taking logarithms and the
product over $\tau$,
\[\log\bdeg{S}+f\sum_{l=k-1}^{r-1}\log l!=\sum_\tau\ \sum_{l=k-1}^{r-1}\ \sum_{i<l}\log\bigl(\rowPar{D}{\tau}{i}-\rowPar{D}{\tau}{l}\bigr).\]
All terms are non-negative by Lemma~\ref{l:rowpar}. \emph{Lower bound.} Keep the terms with $i<k-1$. For these, with
the notation of the proof of Lemma~\ref{l:rowpar}, $\rho(i)-\rho(l)-b_\ast=k/2-1-i+2(l+2-k)h_2\ge0$, and
$\nw{D}{\tau}\le w_0-h_2=b_\ast$ by \Inp{D1}(c); so each of them is at least $\log\nw{D}{\tau}$. There are $k-1$ values
of $i$ and $r-(k-1)=m$ values of~$l$. \emph{Upper bound.} By Lemma~\ref{l:rowpar}, each term is at most
$\log\wt{D}{l+2-k}{\tau}$, which is non-negative, so the sum over $i<l$ is at most $r\log\wt{D}{l+2-k}{\tau}$, and
$\log\wt{D}{l+2-k}{\tau}\le\log\nw{D}{\tau}+1/\scx{D}$ by \Inp{D1}(c), as $l+2-k\le m$. There are $m$ values of~$l$.
\end{proof}

\sect{Notes on the formal text}{s:notes}
\subhead{Arguments that the text reading of the census gets wrong.} The build of this note takes the arguments of
each declaration from a reading of the text, not from Lean. For the following declarations that reading evidently
differs from Lean's signature, and the statements above follow Lean. The Poincar\'e series takes $X$, $\phi$,
$\Delta$, $g$; the census adds the pair~$a$, which the definition binds in its sum. The Poincar\'e map takes only
the hypothesis $X$, implicitly; the census lists the section variables $a$, $\phi$, $g$. For the unfolding
(Proposition~\ref{p:unfold}), for Lemma~\ref{l:tail} and for Theorems C and~D the census stops at the first
\texttt{let} of the statement and misses what follows it: the data $S$, $X$, $\phi$ and the admissibility in the
first; the shape, $t$ and $x$ in the second; the conclusions with their quantifiers in the last two. In the file of
Section~\ref{s:level} the ring $A$ is declared twice; the census lists the instances of the first declaration
(ring, involution, topology, measure) also for the statements about a semilocal field, where Lean takes the
second. For Theorem~E the census lists a form $T$, a level and the hypothesis that it is standard, which are
section variables that the statement does not use: it binds its own level. For the gauge the census lists the
section variables $N$, $t$, $x$, which the definition does not use. None of this was confirmed by running Lean.

\subhead{Definitions that carry a proof.} The Poincar\'e map (Theorem~A(b)) is a definition whose value is a pair,
a linear map with the proof of its formula; the forms $y_a$ are obtained by the axiom of choice from the
existence statement \Inp{L1}. The isomorphism $\wbe{h}$ of Definition~\ref{d:colc} is likewise a pair, with its
transformation rule. The gauge of Definition~\ref{d:gauge} is a record whose fields are the axioms verified in
Lemma~\ref{l:gauge}; its last two explicit arguments are the proofs of $\alpha\in O_A$ and $\bar\alpha=\alpha$, with
default values for $\alpha=1$.

\subhead{The side.} The Bessel functions in the cut Bessel functions, hence in Theorems A and~B, are the full ones:
the formal text declares the full side as the default instance, and no statement of these files names the other.
The band-limit forms of Theorem~D belong to the cut functions; this enters only through Input~\Inp{S5}.

\subhead{Hypotheses and constants.} In Proposition~\ref{p:freq} the filling data enter only through the inequalities
$1000\le b$ and $f\le b/1000$, and the proof uses $\den{y_{\sT}}^{-3}$ where the statement has the sixth power. In
Theorem~C admissibility is used once, for $|\kapS{S}{\tau}|\ge1$, and the hypothesis that the inverses of the elements
of $\mathcal{A}$ are $\sT$-integers once, for the differences $s^{-1}-s'^{-1}$. In the proof of Theorem~D the formal
text does not name the constants $\mathsf{C}$ and $\mathsf{e}$: the first is written as a power of $b$ with an exponent
left to unification, the second is whatever the estimate of the frequency sum produces; the values
$\mathsf{C}=b^{1000n^{9}f}$ and the displayed $\mathsf{e}$ are ours, and so is the final count of exponents, which the
formal text leaves to its automation. Lemma~\ref{l:rowpar} and the two comparisons in the proof of Theorem~F(b) are
proved in the formal text by a decision procedure from the monotonicity of the row parameter (unit~147); the
arguments given here are ours. In Theorem~F(a) the right side is written
$\bigl((\sqrt{\sqrt{b}}^{\,5})^{k}\bigr)^{f}$ in the formal text. The measurability lemma absorbed in
Proposition~\ref{p:unfold} stands in a section that assumes admissibility and stable unramification; its proof uses
neither.

\subhead{Inputs.} Proposition~\ref{p:last} belongs to the seesaw argument (unit~141); the file states it here
because it needs the block embedding of the last grading. The header of the file of the sieve names, as related
work on inductive Fourier transforms of Kloosterman integrals and on spectral large sieves, papers of Jacquet
(2003, 2004), Blomer and Buttcane (2015) and Deshouillers and Iwaniec (1982); we did not compare them with
Theorems C and~D, which are not taken from them. The notes of unit~132 write the head weight, the damped power and
the band constant with other symbols; this note follows the common table of the notes of this unit. Of this unit's declarations, the main-term estimate of Note~F uses Theorem~D,
Proposition~\ref{p:slab} and the masses of Section~\ref{s:level}; Notes T and~G use the standard level and its
gauge; the admissibility of the filling data of $D$ (Note~G) uses Theorem~F.

\subhead{How this note was written.} From the dossier of the unit: the texts of the members of these five files
and their structured proofs. I read the five Lean files completely, line by line, and compared every statement of
this note with the Lean statement, argument by argument, reading in full the statements that the census cuts at a
\texttt{let} (the unfolding, the cut Bessel function of a translate, Lemma~\ref{l:tail}, Proposition~\ref{p:slab},
Theorems C and~D). The readers of the dossier warn that automation hides arguments in three proofs; I read these in
the source: the Fourier coefficients in the proof of Theorem~C (the computation of the character and the vanishing
at frequencies that are not $\sT$-integers are written out above), the single construction in the proof of
Theorem~E (described axiom by axiom), and the reindexing of Weyl's formula in Theorem~F(b) (replaced by the
computation above). The long definition of the gauge and the definitions of $\wbe{h}$ and of the Poincar\'e map,
which have no structured proof or a short one, were read in the source as well. Of the inputs I read the Lean
statements and the definitions they rest on (hermitian spaces with frame, gradings, forms, idempotent kernels,
levels, gauges, hyperbolic bases, integral places, the type of a Bessel system, the row parameter and the weights of
the data) and not the proofs; the inputs from Notes S and~B were compared with their Lean statements. The proofs of
Theorems C and~D, of Propositions \ref{p:freq}, \ref{p:heis}, \ref{p:mass}, \ref{p:stack} and of Theorems A and~B
follow the formal ones step by step, with the exponents of Theorem~D made explicit; the proofs of Theorem~E, of
Lemma~\ref{l:gauge} (axiom (G10)), of Proposition~\ref{p:slab} and of Proposition~\ref{p:last} describe the formal
construction and are less detailed than the formal text. References are those of the catalogue of the write-up and
of the second reader's verdicts; I did not open the cited works. The program check of the unit's dossier covered
\UnitSteps{} leaf steps of \UnitProofs{} proofs.

\head{Appendix A. Concordance with the Lean artifact}
The formal names of the objects: \texttt{S.Fits V l kw Y η} is $\Fits{S}{V}{l}{\mathsf{k}}{\mathcal{Y}}{\eta}$, and the
records named \texttt{X} or \texttt{h} of this type are the hypothesis~$X$; \texttt{V.Γ}, \texttt{V.x}, \texttt{V.xFin},
\texttt{V.J} are $\Gamma(V)$, $\mathrm{fr}_V$, $\mathrm{fr}_{V,f}$, $J_V$; \texttt{G E ε} is $G(V)$; \texttt{Y.unipotent},
\texttt{Y.besselGroup}, \texttt{Y.inclX} are $\Uni{\mathcal{Y}}$, $\Bes{\mathcal{Y}}$, $\incl{\mathcal{Y}}$;
\texttt{typeKernel l ε} is $\epsilon_l$; \texttt{toBesselH} is $\eta\mapsto\dot\eta$; the finite set \texttt{T} of places
is $\sT$, \texttt{LocProd T} is $\ET{\sT}$, \texttt{ratIntegralOff T} the $\sT$-integers, \texttt{ratToPlaces T} the map
$y\mapsto y_{\sT}$, \texttt{adeleOfLocal} the extension by zero, \texttt{padTuple p t} the sequence $(p\,|\,t)$;
\texttt{unitOrOne} is $u$, \texttt{wave} is $c\mapsto\exp(c-\bar c)$; the finite sets \texttt{A}, \texttt{V}, \texttt{U},
\texttt{X} of the sieve are $\mathcal{A}$, $\mathcal{V}$, $\mathcal{A}'$, $\mathcal{X}$, and \texttt{Γ}, \texttt{ε},
\texttt{En} are $\mathsf{C}$, $\mathsf{e}$, $\mathrm{En}$; \texttt{F : Level ι T m} is $\cL\in\Levels{\Omega}{T}{\mathfrak{m}}$,
with \texttt{F.g}, \texttt{F.boxSet}, \texttt{F.𝔅}, \texttt{F.positions}, \texttt{F.θ}, \texttt{F.measureH} the unit, the
box, the ball, the positions and the measures of the floors $\varrho$ and $n$; \texttt{unitaryLevel R g J} is
$\uLev{g}{T}$, \texttt{intUnits} is $O_A^{\times}$, \texttt{idealZeta R m} is $\zetaA{\mathfrak{m}}$, \texttt{md} is $\modA$;
\texttt{D.filling} is $\SD{D}$, \texttt{D.rowParameter} and \texttt{D.paramW} are the row parameter and its double.
Three statements, as they stand in the artifact:
\quoted{peterssonpoincare}
\quoted{sumfrequenciesle}
\quoted{existsisLatticeLevel}

{\fontsize{8.6}{11.6}\selectfont\renewcommand{\leanfont}{\ttfamily\fontsize{7.7}{10}\selectfont\color{leancol}}
\setlength{\LTleft}{0pt}\setlength{\LTright}{\fill}\setlength{\tabcolsep}{0pt}
\begin{longtable}{@{}p{30mm}p{56mm}p{43mm}r@{}}
\textbf{In this note} & \textbf{Lean declaration} & \textbf{File} & \textbf{Lines}\\\midrule\endhead
\ConcordanceRows
\end{longtable}}

\head{Notation}
\begingroup\fontsize{8.6}{10.9}\selectfont\raggedright
\setlength{\columnsep}{9mm}\setlength{\parskip}{0pt}
\begin{multicols}{2}
\nt{$E$, $\tau$, $x_\tau$, $f=\ninf{E}$, $\rdE{E}$}{a CM field, an infinite place, the image of $x$ at $\tau$, the number of infinite places, the root discriminant}
\nt{$w$, $E_w$, $\F_w$, $\AdE$, $\AEf$, $\hO$, $\volO{E}$}{a finite place, completion, residue field, adeles, finite adeles, their integers and their volume}
\nt{$\fs=(\varrho,\kappa,k_0)$, $\rh{\fs}$, $\kk{\fs}$, $\nn{\fs}$, $\rr{\fs}$, $\Ix{\fs}$}{a shape and its numbers $\varrho$, $k$, $n=\varrho+1$, $r$; the index set of the basis $e_a$, $x_j$, $z$, $e^*_a$}
\nt{$\fr{a}$, $V_l$, $h_T$, $\JT{T}$, $\GT{T}$, $\UT{T}$, $\HT{T}$, $\ST{T}$, $\Hl{l}{T}$, $\Ql{l}{T}$, $\Nb{l}{T}$, $\Nf{l}{T}$}{the frame, the span of the floors $\ge l$, the form, its split Gram matrix and unitary group, the Bessel subgroups}
\nt{$\Trg{\fs}{N}$, $\Tlatg{\fs}{N}$, $\blb{\fs}{p}$, $\Gone{\fs}$, $\sd{u}$}{the forms in the rational and the lattice basis, the lattice basis, the complex group, the simple root sum}
\nt{$S\in\Fill{E}{\fs}$}{filling data, with $\NS{S}$, $\qS{S}{j}$, $\gS{S}$, $\piS{S}$, $\vpS{S}$, $\bS{S}$, $\ChS{S}{j}$, $\KcolS{S}$, $\BsysS{S}{\tau}$}
\nt{$\Trat{S}{R}$, $\Tlat{S}{R}$, $\GS{S}{R}$, $\GFin{S}$, $\latB{S}$, $\tolb{S}$}{forms and groups of $S$ over a ring $R$; the lattice basis and the change to it}
\nt{$\KS{S}$, $\GRes{S}$, $\redS{S}$, $\KgS{S}$, $\stdK{S}$}{the integral subgroup, the finite group, the reduction, its kernel; the principal level at the chosen places}
\nt{$\radl{S}$, $\ainf{S}$, $\afin{S}$, $\rinf{S}$, $\rfin{S}$, $\incH{S}$, $\bsGL{S}{\tau}$}{the maps between rational, adelic, real and finite points; the inclusion of the definite group; the scales}
\nt{$\Cl{S}$, $\hc{S}{c}$, $\pos{S}{c}{c'}{\gamma}$, $\terms{S}{c}{s}{c'}$, $\clU{S}{c}$, $\cutoff{S}{c}$}{classes, representatives, positions, terms, the arithmetic unipotent group, the cutoff}
\nt{$\Bc{S}$, $I_\tau(S)$, $W^{\sigma}_{S,\tau}$, $\besfn{\sigma}{S}{y}{i}$, $\besker{\sigma}{S}{\gamma}{y}$, $\deg_\tau(S)$, $\bdeg{S}$}{Bessel coordinates, Bessel functions for the side $\sigma$, the kernel, the degrees}
\nt{$\amp{\sigma}{S}{c}{c'}{\gamma}$, $\ampabs{\sigma}{S}{c}{c'}{\gamma}$, $\ampsum{\sigma}{S}{\Xi}{c}{s}{c'}$}{amplitudes and amplitude sums}
\nt{$\sdc{S}$, $\psi_f$, $\psi_E$, $\kapS{S}{\tau}$, $\fscale{S}{\tau}{a}$}{the self-dual character, the standard characters, the numbers $\kappa_\tau$ and the frame scales}
\nt{$\damp{j}{\rho}$, $\hw{K}{i}{p}$, $\fhw{K}{i}{p}$, $\gix{\xi}$, $\den{\xi}$, $\bandc{\fs}$}{damped power, head weights, generalized index, denominator, band constant}
\nt{$\Sieve{u}{\iota}{\mathsf{n}}{J}{\psi}{\mathrm{En}}{\Lambda}{\beta}$}{a sieve system of length $\mathsf{n}$ on a ring $C$ with units $U$; $p[i\mapsto x]$ a sequence with one entry replaced}
\nt{$V$, $I$, $I_0$, $J_V$, $\mathrm{fr}_V$, $G(V)$, $G_\infty(V)$, $G_f(V)$, $\Gamma(V)$, $\Delta$}{a hermitian space with frame, its index sets, Gram matrix, frame, groups, rational points, a subgroup}
\nt{$\Gf{E}{I}$, $\Kf{E}{I}$, $R(k)$, $\epsilon$, $\epsilon_l$, $\langle F,F'\rangle$}{finite adelic unitary group, its integral points, right translation, a type, the type of a parameter, the Petersson product}
\nt{$\mathcal{Y}$, $\lambda$, $\Uni{\mathcal{Y}}$, $\Bes{\mathcal{Y}}$, $\incl{\mathcal{Y}}$}{a grading with a string of length $\lambda$, its unipotent subgroup, Bessel group, inclusion of the definite group}
\nt{$\Fits{S}{V}{l}{\mathsf{k}}{\mathcal{Y}}{\eta}$, $X$, $\mathsf{p}$, $\dot\eta$, $k_0$, $u$}{$S$ fits $V$; the hypothesis; the frame of the signs; $\eta$ on the definite part; the matrix and the idele of (F4)}
\nt{$\ratEq{X}$, $\adEq{X}$, $\reH{X}$, $\finEq{X}$}{the four dictionaries}
\nt{$\cutB{X}{a}{g}$, $\Poin{X}{\phi}{\Delta}{g}$, $\Pmap{X}$, $e_a$, $y_a$, $W_q$, $\gamma_q$}{cut Bessel function, Poincar\'e series, Poincar\'e map, unit vectors, the forms and the partial sums of the proof}
\nt{$\wbe{h}$, $\colc{h}$, $\mathbb{S}_{\mathsf{k}}$, $\mathsf{A}^{(\tau)}$, $g_c$}{Weyl to Bessel coordinates, column coordinates, the representation of weight $\mathsf{k}$, representatives}
\nt{$\sT$, $\ET{\sT}$, $O_{\sT}$, $y_{\sT}$, $x^{0}$, $\charT{S}{\sT}$, $\fgain{S}{\sT}{\xi}$, $\ratpt{\sT}{l}{t}$}{a finite set of places, the product of completions, integers, image, extension by zero, the character, the gain, a point}
\nt{$\mathcal{A}$, $\mathcal{A}'$, $\mathcal{V}$, $\mathcal{X}$, $\mathcal{N}$, $Z$, $B$, $B_0$, $\rho_\tau$}{finite sets of $\sT$-units and of frequencies; a neighbourhood and a measurable set; bounds}
\nt{$C$, $U$, $\mathrm{hd}_i(p,t)$, $(p\,|\,t)$, $c_a$, $\mathsf{C}$, $\mathsf{e}$, $\Lambda^{(i)}$, $F$, $G$, $\hat G$, $\mathrm{m}$, $U_0$, $L_y$}{the objects of Theorem~C and of its proof}
\nt{$J^{\tau}$, $J'$, $e$, $x_\tau(p,t)$, $y(p,t)$, $m_j(p)$, $F_{\tau,j}$, $\omega$}{the objects of Theorem~D and of its proof}
\nt{$A$, $O_A$, $A^{-}$, $\modA$, $\psi_A$, $\tau_A$, $\delta_A$, $\nu_A$, $\zetaA{\mu}$, $\mathrm{ct}$}{a place, its integers, skew elements, module, character, trace-one element, skew unit, measure on $A^{-}$, ideal zeta sum, content}
\nt{$\Omega$, $\Omega'$, $\Omega_b$, $\mathfrak{m}$, $\cL\in\Levels{\Omega}{T}{\mathfrak{m}}$, $\HypB{\Omega}{T}{\mathfrak{m}}$}{finite sets, the number of pairs, a level, hyperbolic bases}
\nt{$\EL{\cL}$, $\gL{\cL}$, $\Ball{\cL}$, $\BoxL{\cL}$, $\Posn{\cL}$, $\floorMeasure{\cL}{l}$, $\eta_l$, $\etab{j}$, $\varpi_l$, $\quadric{\cL}{b}$}{the data of a level, floor measures, Heisenberg measures, Weyl laws, quadric measures}
\nt{$\stackMeasure{\cL}_{a,i}$, $\groupMeasure{\cL}$, $\upperMeasure{\cL}{l}$, $\floorModulus{\cL}{l}{g}$, $\volumeExponent{\cL}$}{stack, group and upper measures, floor module, volume exponent}
\nt{$\uLev{g}{T}$, $\cgp{\cL}$, \m{\IsLat{\cL}}, \m{\IsStd{\cL}}, $\stdg{\cL}{\alpha}$, $\Gauges{\cL}$, $\Gamma_j$, $\Lambda$}{congruence subgroups, lattice and standard levels, the gauge of depth $\alpha$}
\nt{$\tailv{\pi}{t}{x}$, $\mathrm{n}_a(v,\beta)$, $\wo{l}{c}$, $\mathrm{E}(e)$, $\varsigma'$, $\varphi$, $\Phi$, $\Psi$, $q$, $\mathsf{N}$, $\Sigma_g$, $K$, $M$, $M'$, $Y$, $W$, $d$}{tail vector, Heisenberg and Weyl elements, Eichler elements, charts and functions of Section~\ref{s:level}}
\nt{$P$, $D$, $\SD{D}$, $\ED{D}$, $\scx{D}$, $\scb{D}$, $\kD{D}$, $\fD{D}$, $\Dm{D}$, $\Dr{D}$, $\aD{D}$, $\hhD{D}$, $\hD{D}$}{parameters, admissible data and their quantities}
\nt{$\wt{D}{j}{\tau}$, $\nw{D}{\tau}$, $\rowPar{D}{\tau}{i}$, $\parW{D}{\tau}{j}$, $\Ridx{D}$, $\nu$, $b_\ast$, $a_\tau$, $\rho(i)$}{weights, nominal weight, row parameter and its double, row indices, their numbering}
\nt{$\Ys{D}{M}$, $X_M(D)$, $\Stg{D}{M}$, $\FormY{D}{M}$, $\Col{D}$, $\gYlast{D}$, $\vartheta$, $\Psi$, $\chainY{D}{\vartheta}{M}$, $\iot{D}{\vartheta}$, $\chcomb{D}{\vartheta}{\Psi}{M}$, $\Bmu{D}{\vartheta}{\Psi}$, $\coef{\mathsf{t}}{\psi}$, $R$, $C_{\mathsf{t}}$, $\Xi$}{the objects of Proposition~\ref{p:last}}
\nt{$b$, $g$, $k$, $l$, $n$, $r$, $s$, $t$, $c$, $c'$, $h$, $i$, $i'$, $j$, $m$, $p$, $q$, $v$, $y$, $z$, $\alpha$, $\beta$, $\gamma$, $\delta$, $\theta$, $\mu$, $\xi$, $\sigma$, $\phi$, $\varphi$, $\eta$, $\pi$, $\psi$, $\rho$, $\iota$, $N$, $T$, $H$, $L$, $Q$}{local letters, as each statement says}
\end{multicols}
\endgroup

\head{References}
{\fontsize{9.4}{12.6}\selectfont
\nt{[Bo]}{A.~Borel, \emph{Introduction to automorphic forms}, Proc.\ Sympos.\ Pure Math.~9 (1966). Named by the formal
text as the source for Poincar\'e series; place not confirmed.}
\nt{[Bou]}{N.~Bourbaki, \emph{Int\'egration}, Chap.~VII (1963). Named by the catalogue of the write-up for Weil's
integration formula; place not confirmed.}
\nt{[GL]}{D.~Gentner, G.~Last, \emph{Palm pairs and the general mass-transport principle}, §2, equations (2.5)--(2.9)
(arXiv:0902.0068). The second reader's source for the unfolding \Inp{L2}.}
\nt{[Ta]}{J.~Tate, \emph{Fourier analysis in number fields and Hecke's zeta-functions} (1950), Theorem~4.1.4 and
Lemmas 4.2.1--4.2.4. The catalogue's and the second reader's source for \Inp{L3}, which is a weaker form.}
}

\end{document}
